% Cinétique chimique — Chimie Tle E — Cours signé OnBuch+
% Programme officiel MINESEC, Terminale C, D, E. Compiler avec : tectonic L11-cinetique-chimique.tex
\newcommand{\DOCMATIERE}{Chimie}
\newcommand{\DOCNIVEAU}{Tle E}
\newcommand{\DOCMODULE}{Module 3 · Cinétique chimique}
\newcommand{\DOCLECON}{11}
\newcommand{\DOCTITRE}{Cinétique chimique}
% OnBuch+ — préambule « cours signés » ENRICHI (compilé avec tectonic / XeLaTeX)
% Système visuel « Pop » d'OnBuch+ : fond crème (jamais blanc), bordures encre
% épaisses, ombres « dures » décalées, titres Archivo Black en capitales.
% Version enrichie pour les cours de chimie : chimie (mhchem, chemfig),
% graphes (pgfplots), boîtes pédagogiques supplémentaires (définition,
% propriété, méthode, attention, le savais-tu, exercice résolu, exercice,
% TP, bilan), page de garde refaite et signature OnBuch+ ancrée en bas.
%
% Macros attendues dans main.tex AVANT \input{preamble} :
%   \DOCMATIERE  \DOCNIVEAU  \DOCMODULE  \DOCLECON (numéro)  \DOCTITRE
\documentclass[11pt]{article}

\usepackage{fontspec}
\usepackage[french]{babel}
\usepackage[a4paper,margin=2cm,top=3.4cm,bottom=2.8cm,headheight=1.4cm,footskip=1.3cm]{geometry}
\usepackage{amsmath,amssymb,amsfonts}
\usepackage[table]{xcolor} % [table] : fournit aussi \rowcolors (alternance de lignes)
\usepackage{pagecolor}
\usepackage{tikz}
\usetikzlibrary{arrows.meta,positioning,calc,shapes.geometric,shapes.misc,shapes.arrows,decorations.pathmorphing,decorations.markings,patterns,shadows,mindmap,trees,fit,backgrounds,matrix,intersections}
\usepackage{pgfplots}
\pgfplotsset{compat=1.18}
\usepackage[version=4]{mhchem}
\usepackage{chemfig}
\usepackage{enumitem}
\usepackage{fancyhdr}
% [most] charge toutes les bibliothèques tcolorbox (listings, minted, raster,
% documentation...) et gonfle inutilement la mémoire TeX sur un document long
% (~300 boîtes) : on ne charge que ce qui est réellement utilisé.
\usepackage[skins,breakable]{tcolorbox}
\usepackage{graphicx}
\usepackage{colortbl}
\usepackage{booktabs}
\usepackage{tabularx}
\usepackage{array}
\usepackage{multicol}
\usepackage{siunitx}
% parse-numbers=false : accepte aussi \SI{2,5 \times 10^{-3}}{...}
\sisetup{locale=FR,output-decimal-marker={,},group-minimum-digits=4,parse-numbers=false}
\usepackage{qrcode}
% \degree : commande fréquemment utilisée par les modèles pour un angle ou une
% température en dehors de \SI{}{} (non fournie par les paquets chargés).
\providecommand{\degree}{^{\circ}}
\usepackage{lastpage}
\usepackage{hyperref}
\hypersetup{hidelinks}

% ============================================================
% Palette « Pop » OnBuch+ — mêmes hex que la classe P (plus_ui.dart)
% ============================================================
\definecolor{poporange}{HTML}{F59321}
\definecolor{popdark}{HTML}{E07A0C}
\definecolor{popcream}{HTML}{FFF6E8}
\definecolor{popink}{HTML}{1C1714}
\definecolor{poppurple}{HTML}{7A5AE0}
\definecolor{popgreen}{HTML}{2E9E6A}
\definecolor{poppink}{HTML}{E0457B}
\definecolor{popblue}{HTML}{2F7DE1}
\definecolor{popgold}{HTML}{C98A12}
\definecolor{popmuted}{HTML}{8A7F76}
\definecolor{popline}{HTML}{EADFD2}
\definecolor{popgray}{HTML}{8A7F76}
\colorlet{popgrey}{popgray}
% Alias tolérants (le modèle nomme parfois les couleurs différemment)
\colorlet{popwhite}{white}
\colorlet{popblack}{popink}
\colorlet{popred}{poppink}
\colorlet{popyellow}{popgold}
% Teintes claires pour fonds de tableaux / zones
\colorlet{popblueL}{popblue!10!white}
\colorlet{poporangeL}{poporange!14!white}
\colorlet{popgreenL}{popgreen!12!white}
\colorlet{poppurpleL}{poppurple!12!white}
\colorlet{poppinkL}{poppink!12!white}
\colorlet{popgoldL}{popgold!14!white}

\pagecolor{popcream}

% ============================================================
% Typographie
% ============================================================
\defaultfontfeatures{Ligatures=TeX}
\setmainfont{PlusJakartaSans-Medium.ttf}[Path=../../../fonts/,BoldFont=PlusJakartaSans-Bold.ttf]
% Maths en sans-serif (Fira Math) avec chiffres et lettres droites de Plus
% Jakarta, pour que les formules se fondent dans le texte.
\usepackage[math-style=ISO,bold-style=ISO]{unicode-math}
\setmathfont{FiraMath-Regular.otf}[Path=../../../fonts/]
\setmathfont{PlusJakartaSans-Medium.ttf}[Path=../../../fonts/,range={up/num,up/{Latin,latin},\mathup}]
\usepackage{icomma} % virgule décimale sans espace en mode math
% Caractères absents de la police de texte (sinon carré vide) : rendus par
% leur équivalent mathématique, en texte comme en formule.
\usepackage{newunicodechar}
\newunicodechar{α}{\ensuremath{\alpha}}
\newunicodechar{Α}{\ensuremath{\alpha}}
\newunicodechar{β}{\ensuremath{\beta}}
\newunicodechar{δ}{\ensuremath{\delta}}
\newunicodechar{✓}{\popcheck{popgreen}}
\newfontfamily\popdisplay{ArchivoBlack-Regular.ttf}[Path=../../../fonts/]
\newfontfamily\poplogo{SpaceGrotesk-Bold.ttf}[Path=../../../fonts/]
\newfontfamily\popsemi{PlusJakartaSans-SemiBold.ttf}[Path=../../../fonts/]
\newfontfamily\popxbold{PlusJakartaSans-ExtraBold.ttf}[Path=../../../fonts/]
\newfontfamily\popxboldls{PlusJakartaSans-ExtraBold.ttf}[Path=../../../fonts/,LetterSpace=3.0]

\setlist[enumerate]{leftmargin=1.7em,itemsep=5pt,topsep=4pt}
\setlist[itemize]{leftmargin=1.7em,itemsep=4pt,topsep=4pt,label={\color{poporange}\textbullet}}
\setlength{\parindent}{0pt}
\setlength{\parskip}{5pt}
\renewcommand{\arraystretch}{1.25}

% chemfig : liaisons un peu plus épaisses, encre Pop
\setchemfig{atom sep=2em,bond style={line width=0.9pt,popink},cram width=3pt}

% pgfplots : style Pop par défaut
\pgfplotsset{
  every axis/.append style={axis line style={popink,line width=1pt},tick style={popink},
    grid style={popline,line width=0.5pt},label style={font=\small},tick label style={font=\footnotesize}},
  popcurve/.style={poporange,line width=1.8pt,no markers,smooth},
}

% ============================================================
% Pill / squiggle
% ============================================================
\newcommand{\poppill}[3]{%
  \tikz[baseline=-0.55ex]\node[draw=popink,line width=1.2pt,rounded corners=2.6pt,%
    fill=#2,inner xsep=6.5pt,inner ysep=3pt]{%
    \popxboldls\fontsize{7}{7}\selectfont\color{#3}\MakeUppercase{#1}};%
}
\newcommand{\popsquiggle}[1]{%
  \tikz[baseline=-0.4ex]{
    \draw[#1,line width=2.6pt,line cap=round]
      plot [smooth] coordinates {(0,0.09) (0.34,0.01) (0.68,0.21) (1.02,0.01) (1.36,0.10)};
  }%
}
% Mot-clé surligné (marqueur orange sous le texte)
\newcommand{\cle}[1]{{\popxbold\color{popdark}#1}}

% ============================================================
% En-tête / pied
% ============================================================
\pagestyle{fancy}
\fancyhf{}
\renewcommand{\headrulewidth}{0pt}
\renewcommand{\footrulewidth}{0pt}
\lhead{\raisebox{-0.15ex}{\poplogo\fontsize{13}{13}\selectfont\color{popink}OnBuch\color{poporange}+}}
\rhead{\poppill{\DOCMATIERE\ \textbullet\ \DOCNIVEAU\ \textbullet\ Leçon \DOCLECON}{white}{popink}}
\lfoot{\popxbold\fontsize{7.6}{7.6}\selectfont\color{popmuted}\MakeUppercase{Cours signé OnBuch+}\ \textbullet\ {\popsemi\DOCTITRE}}
\rfoot{\tikz[baseline=-0.6ex]\node[rounded corners=8.5pt,fill=popink,minimum height=17pt,minimum width=17pt,inner xsep=5pt,inner sep=0]{\popxbold\fontsize{8}{8}\selectfont\color{popcream}\thepage\,/\,\pageref*{LastPage}};}

% ============================================================
% Titres
% ============================================================
\newcommand{\chaptitre}[2]{%
  \begin{tcolorbox}[enhanced,colback=poporange,colframe=popink,boxrule=2.6pt,arc=11pt,
      drop shadow=popink,left=15pt,right=15pt,top=12pt,bottom=11pt,before skip=3pt,after skip=18pt]
    \poppill{#2}{popink}{popcream}\par\vspace{9pt}
    {\raggedright\hyphenpenalty=10000\exhyphenpenalty=10000\popdisplay\fontsize{22}{24}\selectfont\color{popcream}\MakeUppercase{#1}\par}
  \end{tcolorbox}
}
% Section numérotée : pastille orange avec le numéro + titre Archivo
\newcounter{popsec}
\newcommand{\coursec}[1]{%
  \stepcounter{popsec}\setcounter{popsub}{0}%
  \par\vspace{16pt}\noindent
  \begin{minipage}{\linewidth}
  \tikz[baseline=-0.7ex]\node[draw=popink,line width=1.6pt,fill=poporange,rounded corners=4pt,minimum size=22pt,inner sep=0]{\popdisplay\fontsize{12}{12}\selectfont\color{popcream}\thepopsec};%
  \hspace{8pt}{\popdisplay\fontsize{15}{17}\selectfont\color{popink}\MakeUppercase{#1}}\par
  \vspace{4pt}\noindent{\color{poporange}\rule{36pt}{3.6pt}}
  \end{minipage}\par\nopagebreak\vspace{8pt}
}
\newcounter{popsub}
\newcommand{\courssub}[1]{%
  \stepcounter{popsub}%
  \par\vspace{9pt}\noindent{\popxbold\fontsize{11.5}{13}\selectfont\color{popink}{\color{poporange}\thepopsec.\thepopsub}\hspace{6pt}#1}\par\nopagebreak\vspace{4pt}
}

% ============================================================
% Boîtes pédagogiques — même recette : fond blanc, bordure épaisse colorée,
% ombre dure encre, badge plein. Titre optionnel [#1].
% ============================================================
\tcbset{popbase/.style={breakable,enhanced,colback=white,boxrule=2.3pt,arc=9pt,
  shadow={3pt}{-3pt}{0pt}{popink},left=14pt,right=14pt,top=11pt,bottom=11pt,before skip=11pt,after skip=15pt,
  fontupper=\popsemi\fontsize{10.6}{14.5}\selectfont\color{popink},
  fontlower=\popsemi\fontsize{10.6}{14.5}\selectfont\color{popink}}}
% Coche dessinée (le glyphe ✓ n'existe pas dans les polices du document)
\newcommand{\popcheck}[1]{\tikz[baseline=-0.5ex]\draw[#1,line width=1.6pt,line cap=round,line join=round](-0.13,0.0)--(-0.04,-0.09)--(0.13,0.1);}
\newcommand{\popboxhead}[3]{\poppill{#1}{#2}{white}\ifx\relax#3\relax\else\hspace{6pt}{\popxbold\fontsize{10.6}{12}\selectfont\color{popink}#3}\fi\par\vspace{7pt}}

\newtcolorbox{aretenir}[1][]{popbase,colframe=popgreen,before upper={\popboxhead{À retenir}{popgreen}{#1}}}
\newtcolorbox{exemplebox}[1][]{popbase,colframe=poporange,before upper={\popboxhead{Exemple}{poporange}{#1}}}
\newtcolorbox{definition}[1][]{popbase,colframe=popblue,before upper={\popboxhead{Définition}{popblue}{#1}}}
\newtcolorbox{propriete}[1][]{popbase,colframe=poppurple,before upper={\popboxhead{Propriété}{poppurple}{#1}}}
\newtcolorbox{methode}[1][]{popbase,colframe=popgold,before upper={\popboxhead{Méthode}{popgold}{#1}}}
\newtcolorbox{attention}[1][]{popbase,colframe=poppink,before upper={\popboxhead{Attention}{poppink}{#1}}}
\newtcolorbox{experience}[1][]{popbase,colframe=popblue,colback=popblueL,before upper={\popboxhead{Expérience}{popblue}{#1}}}
% « Le savais-tu ? » : carte violette pleine (contexte, culture, Cameroun)
\newtcolorbox{savaistu}[1][]{popbase,colback=poppurple,colframe=popink,
  fontupper=\popsemi\fontsize{10.4}{14.2}\selectfont\color{white},
  before upper={\poppill{Le savais-tu\,?}{white}{poppurple}\ifx\relax#1\relax\else\hspace{6pt}{\popxbold\color{white}#1}\fi\par\vspace{7pt}}}
% Exercice résolu : énoncé en haut, \tcblower puis la solution sur fond crème
\newcounter{popexo}\newcounter{popexr}
\newtcolorbox{exoresolu}[1][]{popbase,colframe=poporange,colbacklower=popcream,
  segmentation style={popink,line width=1.2pt,dashed},
  before upper={\stepcounter{popexr}\popboxhead{Exercice résolu \thepopexr}{poporange}{#1}},
  before lower={\poppill{Solution}{popink}{popcream}\par\vspace{6pt}}}
% Exercice (non résolu en place) — la correction est dans la partie Corrigés
\newtcolorbox{exercice}[1][]{popbase,colframe=popink,
  before upper={\stepcounter{popexo}\popboxhead{Exercice \thepopexo}{popink}{#1}}}
% Corrigé
\newtcolorbox{corrige}[1][]{popbase,colframe=popgreen,colback=popgreenL,
  before upper={\popboxhead{Corrigé}{popgreen}{#1}}}
% Objectifs / prérequis / bilan
\newtcolorbox{objectifs}{popbase,colframe=popink,colback=poporangeL,
  before upper={\popboxhead{Objectifs de la leçon}{popink}{}}}
\newtcolorbox{prerequis}{popbase,colframe=popink,colback=popblueL,
  before upper={\popboxhead{Prérequis}{popblue}{}}}
\newtcolorbox{bilan}{popbase,colframe=popink,colback=white,boxrule=2.8pt,
  before upper={\popboxhead{Fiche bilan}{poporange}{L'essentiel à connaître pour l'examen}}}
% Figure encadrée avec légende
\newtcolorbox{popfigure}[1][]{enhanced,breakable=false,colback=white,colframe=popline,boxrule=1.4pt,arc=8pt,
  left=10pt,right=10pt,top=10pt,bottom=8pt,before skip=10pt,after skip=12pt,halign=center,
  lower separated=false,halign lower=center,
  fontlower=\popsemi\fontsize{9}{11.5}\selectfont\color{popmuted},
  #1}
\newcommand{\legende}[1]{\par\vspace{6pt}{\popsemi\fontsize{9}{11.5}\selectfont\color{popmuted}#1}}

% ============================================================
% Signature OnBuch+
% ============================================================
\newcommand{\poplogoinline}[1]{{\poplogo\fontsize{#1}{#1}\selectfont\color{popink}OnBuch\color{poporange}+}}
\newcommand{\signatureonbuch}{%
  \begin{tcolorbox}[enhanced,colback=popink,colframe=popink,boxrule=2.2pt,arc=10pt,
      drop shadow=poporange,left=14pt,right=14pt,top=10pt,bottom=10pt,before skip=0pt,after skip=0pt]
    \tikz[baseline=-0.7ex]\node[circle,fill=popgreen,draw=popcream,line width=1.3pt,minimum size=22pt,inner sep=0]{\popcheck{white}};%
    \hspace{9pt}\begin{minipage}[c]{0.62\linewidth}
      {\popxbold\fontsize{12}{13}\selectfont\color{popcream}Signé\ \textbullet\ {\poplogo OnBuch\color{poporange}+}}\par\vspace{2pt}
      {\raggedright\popsemi\fontsize{8}{10.5}\selectfont\color{popline}\MakeUppercase{Équipe pédagogique OnBuch+}\\\MakeUppercase{Conforme au programme officiel MINESEC}\par}
    \end{minipage}\hfill
    {\poplogo\fontsize{22}{22}\selectfont\color{popcream}OnBuch\color{poporange}+}
  \end{tcolorbox}%
}

% ============================================================
% Page de garde
% #1 titre · #2 matière · #3 niveau · #4 module · #5 n° leçon · #6 durée ·
% #7 description / objectifs (texte ou itemize)
% ============================================================
\newcommand{\pagegarde}[7]{%
  \thispagestyle{empty}
  \newgeometry{a4paper,margin=1.7cm,top=1.6cm,bottom=1.4cm}%
  \begin{tikzpicture}[remember picture,overlay]
    \fill[poporange!18] ([xshift=-3.2cm,yshift=-3.6cm]current page.north east) circle (4.6cm);
    \fill[poppurple!14] ([xshift=2.4cm,yshift=7.6cm]current page.south west) circle (3.8cm);
  \end{tikzpicture}%
  {\poplogo\fontsize{30}{30}\selectfont\color{popink}OnBuch\color{poporange}+}\hfill
  \raisebox{0.6ex}{\poppill{Cours signé \textbullet\ Premium}{popink}{popcream}}\par
  \vspace{1pt}\popsquiggle{poppurple}\par
  \vspace{8pt}
  \begin{tcolorbox}[enhanced,colback=poporange,colframe=popink,boxrule=2.8pt,arc=13pt,
      drop shadow=popink,left=18pt,right=18pt,top=12pt,bottom=13pt,after skip=10pt]
    \poppill{#2 \textbullet\ #3}{popink}{popcream}\hspace{5pt}\poppill{#4}{white}{popink}\par\vspace{8pt}
    {\popdisplay\fontsize{14}{14}\selectfont\color{popink}LEÇON #5}\par\vspace{4pt}
    {\raggedright\hyphenpenalty=10000\exhyphenpenalty=10000\popdisplay\fontsize{25}{27}\selectfont\color{popcream}\MakeUppercase{#1}\par}
    \vspace{8pt}
    {\popxbold\fontsize{9.5}{9.5}\selectfont\color{popink}\MakeUppercase{Durée conseillée : #6}}
  \end{tcolorbox}
  \begin{tcolorbox}[enhanced,colback=white,colframe=popink,boxrule=2.2pt,arc=10pt,
      drop shadow=popink,left=16pt,right=16pt,top=10pt,bottom=10pt,after skip=8pt]
    \poppill{Dans cette leçon}{poppurple}{white}\par\vspace{5pt}
    \popsemi\fontsize{9.3}{12.6}\selectfont\color{popink}#7
  \end{tcolorbox}
  \noindent\begin{minipage}[t]{0.487\linewidth}
    \begin{tcolorbox}[enhanced,colback=popgreen,colframe=popink,boxrule=2.2pt,arc=10pt,
        drop shadow=popink,left=11pt,right=11pt,top=10pt,bottom=10pt,height=3.1cm,valign=center]
      \tcbox[colback=white,colframe=popink,boxrule=1.3pt,arc=5pt,boxsep=0pt,nobeforeafter,
        left=3pt,right=3pt,top=3pt,bottom=3pt,valign=center]{\qrcode[height=1.9cm]{https://play.google.com/store/apps/details?id=com.luvvix.onbuch&hl=fr}}%
      \hspace{8pt}\begin{minipage}[c]{0.5\linewidth}
        {\popdisplay\fontsize{11}{12.5}\selectfont\color{white}\MakeUppercase{Télécharge OnBuch}}\par\vspace{4pt}
        {\raggedright\popsemi\fontsize{8.4}{11}\selectfont\color{white}Gratuit \textbullet\ Google Play\\Tuteur IA, annales, concours, résultats\par}
      \end{minipage}
    \end{tcolorbox}
  \end{minipage}\hfill
  \begin{minipage}[t]{0.487\linewidth}
    \begin{tcolorbox}[enhanced,colback=poppurple,colframe=popink,boxrule=2.2pt,arc=10pt,
        drop shadow=popink,left=11pt,right=11pt,top=10pt,bottom=10pt,height=3.1cm,valign=center]
      \tikz[baseline=-0.7ex]\node[circle,fill=white,draw=popink,line width=1.3pt,minimum size=1.6cm,inner sep=0]{\popdisplay\fontsize{24}{24}\selectfont\color{poppurple}+};%
      \hspace{8pt}\begin{minipage}[c]{0.6\linewidth}
        {\popdisplay\fontsize{11}{12.5}\selectfont\color{white}\MakeUppercase{Passe à OnBuch+}}\par\vspace{4pt}
        {\raggedright\popsemi\fontsize{8.4}{11}\selectfont\color{white}Tous les cours signés, Tuteur IA illimité, fiches PDF — onglet OnBuch+ de l'app\par}
      \end{minipage}
    \end{tcolorbox}
  \end{minipage}
  \vspace{8pt}
  \popcontenu
  \vfill
  \signatureonbuch
  \restoregeometry
  \newpage
}

% Rangée « Ce cours contient » (page de garde)
\newcommand{\poptile}[3]{% #1 couleur #2 gros texte #3 légende
  \begin{tcolorbox}[enhanced,colback=#1,colframe=popink,boxrule=2pt,arc=9pt,shadow={2.5pt}{-2.5pt}{0pt}{popink},
      left=6pt,right=6pt,top=9pt,bottom=9pt,halign=center,valign=center,height=1.75cm,width=\linewidth,nobeforeafter]
    {\popdisplay\fontsize{12.5}{13}\selectfont\color{white}\MakeUppercase{#2}}\par\vspace{3pt}
    {\centering\popsemi\fontsize{7.8}{9.5}\selectfont\color{white}#3\par}
  \end{tcolorbox}}
\newcommand{\popcontenu}{%
  \par\noindent{\popxboldls\fontsize{7.5}{7.5}\selectfont\color{popmuted}CE COURS SIGNÉ CONTIENT}\par\vspace{7pt}
  \noindent\begin{minipage}[t]{0.235\linewidth}\poptile{popblue}{Cours}{complet, illustré, pas à pas}\end{minipage}\hfill
  \begin{minipage}[t]{0.235\linewidth}\poptile{poporange}{Méthodes}{et exercices résolus}\end{minipage}\hfill
  \begin{minipage}[t]{0.235\linewidth}\poptile{poppink}{Exercices}{type examen + corrigés}\end{minipage}\hfill
  \begin{minipage}[t]{0.235\linewidth}\poptile{popgreen}{Fiche bilan}{l'essentiel à retenir}\end{minipage}\par}

% Sommaire
\newtcolorbox{sommaire}{enhanced,colback=white,colframe=popink,boxrule=2.2pt,arc=10pt,
  drop shadow=popink,left=18pt,right=18pt,top=13pt,bottom=13pt,before skip=6pt,after skip=20pt,
  before upper={\poppill{Sommaire}{poppurple}{white}\par\vspace{9pt}\popsemi\fontsize{10.6}{14.5}\selectfont\color{popink}}}

% Signature de fin de cours — poussée tout en bas de la dernière page
\newcommand{\findecours}{%
  \par\vfill\nopagebreak
  \begin{center}\popsquiggle{poporange}\end{center}\vspace{4pt}
  \signatureonbuch
}

%% FIGFIX-BEGIN v4 — correctifs globaux des figures (docs/onbuch-generation/tools/figfix)
\usepackage{adjustbox}\usepackage{etoolbox}
% toute figure TikZ/pgfplots plus large que la ligne est réduite à la largeur disponible
\BeforeBeginEnvironment{tikzpicture}{\begin{adjustbox}{max width=\linewidth}}
\AfterEndEnvironment{tikzpicture}{\end{adjustbox}}
% légendes pgfplots lisibles par-dessus les courbes
\pgfplotsset{every axis/.append style={legend style={fill=white,fill opacity=0.92,text opacity=1,draw=black!15,inner sep=3pt}}}
% étiquettes d'arêtes (syntaxe edge["..."]) avec fond blanc
\tikzset{every edge quotes/.append style={fill=white,inner sep=1.2pt}}
%% FIGFIX-END
\begin{document}

\pagegarde{Cinétique chimique}{Chimie}{Tle E}{Module 3 · Cinétique chimique}{11}{6 h (cours + TP)}{Cette leçon étudie l'évolution temporelle des systèmes chimiques : classification des réactions selon leur rapidité, définition et détermination graphique des vitesses de formation et de disparition, méthodes de suivi cinétique et facteurs permettant d'accélérer ou de ralentir une réaction. Elle s'appuie sur des exemples classiques du programme (I-/S2O8\^{}2-, estérification) et sur des situations concrètes du quotidien camerounais.
\vspace{4pt}
\begin{multicols}{2}\small
\begin{itemize}
\item À la fin de cette leçon, je sais classer des réactions chimiques en réactions rapides, lentes ou très lentes en citant des exemples.
\item À la fin de cette leçon, je sais tracer la courbe de concentration du diiode formé en fonction du temps à partir de résultats expérimentaux.
\item À la fin de cette leçon, je sais déterminer graphiquement une vitesse moyenne et une vitesse instantanée de formation d'un produit ou de disparition d'un réactif.
\item À la fin de cette leçon, je sais définir la vitesse volumique de réaction et la relier à l'avancement.
\item À la fin de cette leçon, je sais déterminer le temps de demi-réaction sur une courbe cinétique.
\item À la fin de cette leçon, je sais décrire les méthodes de suivi temporel d'une réaction (trempe et dosage, colorimétrie, conductimétrie, pression).
\end{itemize}\end{multicols}}

\begin{sommaire}
\begin{multicols}{2}
\begin{enumerate}
\item Réactions rapides, lentes et très lentes
\item Vitesses de formation et de disparition
\item Vitesse volumique, avancement et temps de demi-réaction
\item Suivi temporel d'une réaction
\item Facteurs cinétiques et catalyse
\item Étude cinétique de l'estérification
\item Travaux pratiques
\item Méthodes et exercices résolus
\item Exercices
\item Corrigés des exercices
\item Fiche bilan
\end{enumerate}
\end{multicols}
\end{sommaire}

\begin{objectifs}
\begin{itemize}
\item À la fin de cette leçon, je sais classer des réactions chimiques en réactions rapides, lentes ou très lentes en citant des exemples.
\item À la fin de cette leçon, je sais tracer la courbe de concentration du diiode formé en fonction du temps à partir de résultats expérimentaux.
\item À la fin de cette leçon, je sais déterminer graphiquement une vitesse moyenne et une vitesse instantanée de formation d'un produit ou de disparition d'un réactif.
\item À la fin de cette leçon, je sais définir la vitesse volumique de réaction et la relier à l'avancement.
\item À la fin de cette leçon, je sais déterminer le temps de demi-réaction sur une courbe cinétique.
\item À la fin de cette leçon, je sais décrire les méthodes de suivi temporel d'une réaction (trempe et dosage, colorimétrie, conductimétrie, pression).
\item À la fin de cette leçon, je sais montrer expérimentalement le rôle catalytique des ions Fe3+ et du platine.
\item À la fin de cette leçon, je sais réaliser l'étude cinétique d'une estérification et interpréter l'asymptote de la courbe.
\item À la fin de cette leçon, je sais vérifier et expliquer l'influence de la température, de la concentration, d'un catalyseur et de la surface de contact sur la vitesse d'une réaction.
\end{itemize}
\end{objectifs}

\begin{prerequis}
\begin{itemize}
\item Tableau d'avancement et notion d'avancement x (Première)
\item Concentration molaire, quantité de matière et dilutions (Première)
\item Réactions acide-base et réactions de précipitation (Première)
\item Lecture de graphiques, pente d'une droite et tangente à une courbe (mathématiques)
\item Réaction d'estérification : équation et caractères (début d'année de Terminale)
\end{itemize}
\end{prerequis}

\newpage

\coursec{Réactions rapides, lentes et très lentes}

Au marché Mokolo de Yaoundé, le poisson exposé en plein soleil se gâte en quelques heures, alors qu'au frais il se conserve plusieurs jours. Pourquoi la température change-t-elle la rapidité des transformations ? Pourquoi le cuisinier ajoute-t-il du sel ou couvre-t-il la marmite pour accélérer la cuisson ? En laboratoire, certaines réactions comme la précipitation du chlorure d'argent sont instantanées, tandis que la rouille sur la tôle d'une case met des mois à se former. Comment mesurer, suivre et surtout \cle{contrôler} la rapidité d'une réaction chimique ? C'est tout l'objet de la \cle{cinétique chimique}.

\courssub{Qu'est-ce que la cinétique chimique ?}

Jusqu'ici, tu as appris à écrire et à équilibrer des équations de réactions : on te disait \emph{ce qui} se transforme et \emph{en quoi}. Mais aucune équation ne te dit \emph{combien de temps} cela prend. Or, dans la vie réelle, cette question est essentielle : une bonne réaction trop lente est inutilisable en industrie ; une mauvaise réaction trop rapide (dégradation d'un aliment, corrosion d'un pont) est un désastre.

\begin{definition}[Cinétique chimique]
La \cle{cinétique chimique} est la partie de la chimie qui étudie l'\cle{évolution des systèmes chimiques au cours du temps} : elle s'intéresse à la \textbf{vitesse} à laquelle les réactifs disparaissent et les produits apparaissent, ainsi qu'aux \textbf{facteurs} (température, concentration, catalyseur, surface de contact) qui modifient cette vitesse.
\end{definition}

L'outil central de la cinétique est donc la \textbf{courbe d'évolution} : on porte en ordonnée une grandeur qui caractérise l'avancement de la réaction (concentration d'un produit, d'un réactif, pression, absorbance...) et en abscisse le temps. Toute la suite de cette leçon consistera à construire, lire et exploiter de telles courbes.

\begin{attention}[Rapidité et évolution spontanée : deux notions différentes]
Erreur fréquente : croire qu'une réaction thermodynamiquement favorisée (spontanée) est forcément rapide. \textbf{C'est faux !} La thermodynamique dit si une réaction \emph{peut} avoir lieu ; la cinétique dit \emph{à quelle vitesse} elle a lieu. La formation de la rouille est spontanée, mais extrêmement lente. Le mélange \ce{H2} + \ce{O2} devrait réagir violemment... et pourtant il peut rester inchangé pendant des années à température ambiante. Ne confonds jamais « possible » et « rapide ».
\end{attention}

\courssub{Classer les réactions selon leur durée}

Le classement « rapide / lente / très lente » est \textbf{qualitatif} : il repose sur la comparaison entre la durée de la réaction et l'\cle{échelle de temps humaine} (la seconde, la minute, l'heure, le jour). Une réaction est dite rapide si elle semble terminée dès que les réactifs sont mis en contact ; lente si on peut suivre son évolution à l'œil nu ou avec des instruments pendant l'expérience ; très lente si elle dure des jours, des mois, des années.

\begin{popfigure}
\begin{center}
\begin{tikzpicture}[scale=1.0]
  % Axe du temps
  \draw[-{Stealth[length=3mm]}, popdark, very thick] (0,0) -- (15.5,0);
  \node[popdark, font=\small\bfseries] at (15.2,-0.5) {durée};
  % Graduations
  \foreach \x/\t in {0.5/{instantané}, 3.2/{milliseconde}, 5.9/{seconde}, 8.6/{minute}, 11.3/{heure}, 14/{jour -- année}} {
    \draw[popdark, thick] (\x,0.12) -- (\x,-0.12);
    \node[popmuted, font=\scriptsize, align=center] at (\x,-0.55) {\t};
  }
  % Réactions rapides
  \node[draw=popblue, fill=popblueL, rounded corners=2pt, font=\scriptsize, align=center, inner sep=3pt] (a) at (1.6,1.5) {Précipitation \ce{AgCl}\\ \ce{Ag+ + Cl- -> AgCl}};
  \draw[popblue, thick] (a.south) -- (1.6,0.12);
  \node[draw=popblue, fill=popblueL, rounded corners=2pt, font=\scriptsize, align=center, inner sep=3pt] (b) at (4.3,2.4) {Neutralisation acide--base\\ \ce{H3O+ + HO- -> 2 H2O}};
  \draw[popblue, thick] (b.south) -- (4.3,0.12);
  % Réactions lentes
  \node[draw=poporange, fill=poporangeL, rounded corners=2pt, font=\scriptsize, align=center, inner sep=3pt] (c) at (7.4,1.5) {\ce{2 I- + S2O8^2- -> I2 + 2 SO4^2-}\\ (quelques minutes)};
  \draw[poporange, thick] (c.south) -- (7.4,0.12);
  \node[draw=poporange, fill=poporangeL, rounded corners=2pt, font=\scriptsize, align=center, inner sep=3pt] (d) at (10.1,2.4) {Estérification, saponification\\ (heures)};
  \draw[poporange, thick] (d.south) -- (10.1,0.12);
  % Très lentes
  \node[draw=poppink, fill=poppinkL, rounded corners=2pt, font=\scriptsize, align=center, inner sep=3pt] (e) at (13.3,1.5) {Rouille, corrosion\\ synthèse de l'eau sans catalyseur};
  \draw[poppink, thick] (e.south) -- (13.3,0.12);
  % Étiquettes de zones
  \node[popblue, font=\small\bfseries] at (2.9,3.4) {RAPIDES};
  \node[poporange, font=\small\bfseries] at (8.7,3.4) {LENTES};
  \node[poppink, font=\small\bfseries] at (13.3,3.4) {TRÈS LENTES};
\end{tikzpicture}
\end{center}
\legende{Frise chronologique : classement qualitatif des réactions chimiques selon leur durée, de l'instantané (précipitation de \ce{AgCl}) au très lent (corrosion).}
\end{popfigure}

\courssub{Les réactions rapides (quasi instantanées)}

\begin{definition}[Réaction rapide]
Une réaction est dite \cle{rapide} (ou quasi instantanée) lorsqu'elle est terminée dès que les réactifs sont mis en contact : sa durée est inférieure à environ une seconde, souvent de l'ordre de la milliseconde. On ne peut pas suivre son évolution à l'œil nu : on n'observe que l'état final.
\end{definition}

\textbf{Exemple 1 : la précipitation du chlorure d'argent.} Dans un tube à essai, versons quelques gouttes de solution de nitrate d'argent (\ce{Ag+} + \ce{NO3-}) dans une solution de chlorure de sodium (\ce{Na+} + \ce{Cl-}). Immédiatement, un \textbf{précipité blanc} de chlorure d'argent apparaît, qui noircit à la lumière par photodécomposition :

\[ \ce{Ag+ (aq) + Cl- (aq) -> AgCl (s)} \]

Cette réaction est totale et instantanée : c'est le test classique d'identification des ions chlorure.

\begin{popfigure}
\begin{center}
\begin{tikzpicture}[scale=0.9]
  % Tube à essai
  \draw[popdark, thick] (-1,0) -- (-1,-4) arc[start angle=180, end angle=360, radius=1] -- (1,0);
  % Liquide
  \fill[popblueL] (-0.92,-1.2) -- (-0.92,-3.9) arc[start angle=180, end angle=360, radius=0.92] -- (0.92,-1.2) -- cycle;
  \draw[popblue, thin] (-0.92,-1.2) -- (0.92,-1.2);
  % Précipité au fond
  \fill[white, opacity=0.9] (-0.75,-3.75) arc[start angle=180, end angle=360, radius=0.75] -- cycle;
  \fill[popmuted!40] (-0.75,-3.75) arc[start angle=180, end angle=360, radius=0.75] -- cycle;
  % Flocons de précipité en suspension
  \foreach \x/\y in {-0.4/-2.2, 0.3/-2.6, -0.1/-3.0, 0.5/-3.3, -0.5/-3.4, 0.1/-1.8} {
    \node[popmuted, font=\tiny] at (\x,\y) {$\bullet$};
  }
  % Pipette qui verse
  \draw[popdark, thick] (1.6,1.2) -- (0.4,0.2);
  \draw[popdark, thick] (1.85,1.45) -- (0.65,0.45);
  \draw[popdark, thick] (1.6,1.2) -- (1.85,1.45);
  \draw[popblue, -{Stealth}, thick] (0.5,0.1) -- (0.3,-0.9);
  \node[popblue, font=\scriptsize, align=left] at (3.6,1.0) {solution de \ce{AgNO3}\\ (ions \ce{Ag+})};
  % Annotation
  \node[popdark, font=\scriptsize, align=left] at (3.6,-1.6) {solution de \ce{NaCl}\\ (ions \ce{Cl-})};
  \draw[popmuted, thin] (2.2,-1.5) -- (0.95,-1.5);
  \node[popdark, font=\scriptsize, align=left] at (3.6,-3.6) {précipité blanc\\ de \ce{AgCl}};
  \draw[popmuted, thin] (2.2,-3.7) -- (0.6,-3.9);
\end{tikzpicture}
\end{center}
\legende{Test d'identification des ions chlorure : dès l'ajout des ions \ce{Ag+}, un précipité blanc de \ce{AgCl} se forme instantanément.}
\end{popfigure}

\textbf{Exemple 2 : les réactions acide--base.} Le mélange d'un acide fort et d'une base forte donne lieu à la neutralisation :

\[ \ce{H3O+ (aq) + HO- (aq) -> 2 H2O (l)} \]

Cette réaction est totale, exothermique et \textbf{extrêmement rapide} : elle se produit en moins d'une milliseconde. C'est pourquoi, lors d'un dosage acide--base, on considère que chaque goutte versée réagit immédiatement.

\textbf{Exemple 3 : les précipitations en général.} La formation de tout précipité (sulfate de baryum \ce{BaSO4}, hydroxyde de fer(III) \ce{Fe(OH)3}...) est une réaction rapide, car elle met en jeu des \textbf{ions en solution} qui se rencontrent sans obstacle : aucune liaison covalente n'a besoin d'être rompue.

\begin{aretenir}[Pourquoi ces réactions sont-elles rapides ?]
Les réactions entre \textbf{ions en solution aqueuse} (précipitations, neutralisations) sont quasi instantanées car les ions sont déjà « prêts à réagir » : il suffit qu'ils se rencontrent. En revanche, les réactions qui exigent la \textbf{rupture de liaisons covalentes} (oxydoréductions entre molécules, estérification...) sont généralement lentes.
\end{aretenir}

\courssub{Les réactions lentes (secondes à heures)}

\begin{definition}[Réaction lente]
Une réaction est dite \cle{lente} lorsque son évolution peut être \textbf{suivie au cours du temps} : elle dure de quelques secondes à quelques heures. On peut alors prélever des échantillons ou mesurer une grandeur physique à différents instants et tracer une courbe d'évolution.
\end{definition}

\textbf{Exemple 1 : l'oxydation des ions iodure par les ions peroxodisulfate.} C'est \emph{la} réaction d'étude de cette leçon. En solution aqueuse, les ions peroxodisulfate \ce{S2O8^2-} oxydent lentement les ions iodure \ce{I-} :

\[ \ce{2 I- (aq) + S2O8^2- (aq) -> I2 (aq) + 2 SO4^2- (aq)} \]

Le \cle{diiode} \ce{I2} formé colore progressivement la solution en \textbf{brun-jaune}, de plus en plus foncé. À température ambiante, la réaction dure \textbf{plusieurs minutes} : on peut donc suivre l'apparition du diiode et tracer la courbe $[\ce{I2}] = f(t)$, ce que nous ferons dans les sections suivantes.

\textbf{Exemple 2 : l'estérification.} La réaction entre l'acide éthanoïque et l'éthanol :

\[ \ce{CH3COOH + C2H5OH <=> CH3COOC2H5 + H2O} \]

est lente (plusieurs heures, voire plusieurs jours à température ambiante), \textbf{limitée} (équilibre) et faiblement exothermique (souvent traitée comme athermique en première approche cinétique car la constante d'équilibre varie peu avec la température). Nous lui consacrerons toute la dernière section de cette leçon.

\begin{popfigure}
\begin{center}
\chemfig{CH_3-C(=[2]O)-[6]OH + H_3C-CH_2-OH <=> CH_3-C(=[2]O)-[6]O-CH_2-CH_3 + H_2O}
\end{center}
\legende{Équation de l'estérification de l'acide éthanoïque par l'éthanol (représentation semi-développée).}
\end{popfigure}

\textbf{Exemple 3 : la saponification.} La réaction d'un corps gras (huile de palme, triglycérides) avec la soude :

\[ \ce{ester + HO- -> alcool + carboxylate (savon)} \]

est lente mais, contrairement à l'estérification, elle est \textbf{totale}. C'est la réaction utilisée dans les savonneries artisanales : on chauffe l'huile de palme avec la lessive de soude pendant de longues minutes pour obtenir le savon.

\begin{exemplebox}[Comparaison chiffrée de deux durées caractéristiques]
À température ambiante (\SI{25}{\celsius}) :
\begin{itemize}
  \item la neutralisation \ce{H3O+ + HO- -> 2 H2O} est achevée en moins de $10^{-3}$ s (\textbf{moins d'une milliseconde}) ;
  \item la réaction \ce{2 I- + S2O8^2- -> I2 + 2 SO4^2-}, avec des concentrations usuelles de l'ordre de \SI{0,10}{mol.L^{-1}}, atteint sa moitié d'avancement en \textbf{quelques minutes} (environ \SI{10}{min}) et n'est totale qu'au bout d'une heure environ.
\end{itemize}
Le rapport des durées est donc supérieur à $10^{5}$ : la seconde réaction est \textbf{plus de cent mille fois plus lente} que la première. C'est toute la différence entre une réaction « rapide » et une réaction « lente ».
\end{exemplebox}

\courssub{Les réactions très lentes (jours, mois, années)}

\begin{definition}[Réaction très lente]
Une réaction est dite \cle{très lente} lorsqu'elle s'étend sur des jours, des mois ou des années. Son évolution est imperceptible à l'échelle d'une séance de travaux pratiques.
\end{definition}

\textbf{Exemple 1 : la formation de la rouille.} Le fer, exposé à l'air humide, s'oxyde lentement en oxyde de fer(III) hydraté (la rouille) :

\[ \ce{4 Fe (s) + 3 O2 (g) -> 2 Fe2O3 (s)} \quad \text{(puis hydratation : \ce{Fe2O3 . $n$ H2O})} \]

Sur la tôle d'une case ou la carrosserie d'une voiture, cette \cle{corrosion} met des mois ou des années à devenir visible. Pourtant, elle est thermodynamiquement spontanée !

\textbf{Exemple 2 : la synthèse de l'eau sans catalyseur.} La réaction :

\[ \ce{2 H2 (g) + O2 (g) -> 2 H2O (g)} \]

est, à température ambiante et sans catalyseur, \textbf{infiniment lente} : un mélange de dihydrogène et de dioxygène peut rester des années sans changement observable.

\begin{savaistu}[Un mélange explosif... qui n'explose pas tout seul]
Le mélange \ce{H2} + \ce{O2} est si lent à réagir à température ambiante qu'il peut rester inchangé pendant des années. Pourtant, il suffit d'une \textbf{flamme} (apport de température) ou d'un peu de \textbf{platine} (catalyseur) pour que la réaction devienne brutalement explosive ! Ce contraste spectaculaire illustre parfaitement la différence entre thermodynamique (la réaction est très favorisée) et cinétique (elle est bloquée). Les chimistes disent que le système est dans un état « métastable » : stable en apparence, mais prêt à évoluer dès qu'on lui donne un « coup de pouce ».
\end{savaistu}

\courssub{La cinétique dans la vie quotidienne au Cameroun}

La cinétique chimique n'est pas qu'une affaire de laboratoire : elle gouverne de nombreuses transformations de notre quotidien.

\begin{itemize}
  \item \textbf{Le rancissement de l'huile de palme} : l'oxydation des corps gras par le dioxygène de l'air est une réaction très lente, accélérée par la chaleur et la lumière. C'est pourquoi on conserve l'huile dans un récipient fermé, à l'abri de la lumière.
  \item \textbf{La fermentation du vin de palme} : la transformation des sucres en éthanol par les levures, \ce{C6H12O6 -> 2 C2H5OH + 2 CO2}, prend quelques heures à quelques jours ; la chaleur du jour l'accélère, ce qui explique que le vin de palme « tourne » vite en plein soleil.
  \item \textbf{La dégradation du poisson} : les réactions de décomposition (enzymatiques et bactériennes) sont lentes au froid, rapides à la chaleur : d'où la situation du marché Mokolo évoquée en introduction.
  \item \textbf{La cuisson des aliments} : couvrir la marmite augmente la température et accélère les transformations chimiques de la cuisson.
\end{itemize}

\courssub{Pourquoi contrôler la vitesse des réactions ?}

L'intérêt pratique de la cinétique est immense : il s'agit de savoir \textbf{réaliser volontairement une réaction rapide} ou, au contraire, \textbf{ralentir une réaction gênante}.

\begin{center}
\begin{tabularx}{\linewidth}{|l|X|X|}
\hline
\rowcolor{popblueL} \textbf{Objectif} & \textbf{Exemples} & \textbf{Moyen utilisé} \\
\hline
Accélérer une réaction utile & Cuisson des aliments ; synthèses industrielles (ammoniac, savon) & Chauffer, utiliser un catalyseur, concentrer les réactifs \\
\hline
Ralentir une réaction nuisible & Conservation des aliments (poisson, huile) ; protection contre la corrosion & Réfrigérer, congeler, peindre ou galvaniser les métaux \\
\hline
\end{tabularx}
\end{center}

\begin{aretenir}[L'essentiel de cette section]
\begin{itemize}
  \item La \cle{cinétique chimique} étudie l'évolution des systèmes chimiques au cours du temps.
  \item Les réactions \textbf{rapides} (moins d'une seconde) : précipitations, réactions acide--base — ce sont des réactions entre ions en solution.
  \item Les réactions \textbf{lentes} (secondes à heures) : \ce{2 I- + S2O8^2- -> I2 + 2 SO4^2-}, estérification, saponification — on peut suivre leur évolution et tracer des courbes.
  \item Les réactions \textbf{très lentes} (jours à années) : corrosion, synthèse de l'eau sans catalyseur, rancissement des huiles.
  \item Ce classement est \textbf{qualitatif} : il dépend de l'échelle de temps humaine.
  \item Une réaction spontanée n'est pas forcément rapide : la cinétique et la thermodynamique répondent à deux questions différentes.
\end{itemize}
\end{aretenir}

\begin{exoresolu}[Classer des réactions selon leur durée]
\textbf{Énoncé.} On réalise les quatre expériences suivantes : (a) on verse une solution de nitrate d'argent dans de l'eau salée ; (b) on mélange des solutions d'iodure de potassium et de peroxodisulfate de sodium ; (c) on laisse un clou en fer à l'air humide ; (d) on mélange de l'acide chlorhydrique et de la soude. Pour chaque expérience, écrire l'équation de la réaction et la classer (rapide, lente ou très lente) en justifiant.
\tcblower
\textbf{Solution.}
\begin{itemize}
  \item (a) \ce{Ag+ + Cl- -> AgCl (s)} : réaction entre ions en solution, précipitation \textbf{instantanée} — réaction \textbf{rapide}.
  \item (b) \ce{2 I- + S2O8^2- -> I2 + 2 SO4^2-} : la coloration brune du diiode apparaît progressivement en quelques minutes — réaction \textbf{lente}.
  \item (c) \ce{4 Fe + 3 O2 -> 2 Fe2O3} : la rouille n'apparaît qu'au bout de plusieurs jours ou semaines — réaction \textbf{très lente}.
  \item (d) \ce{H3O+ + HO- -> 2 H2O} : neutralisation entre ions, achevée en moins d'une milliseconde — réaction \textbf{rapide}.
\end{itemize}
\textbf{Justification du critère :} les réactions (a) et (d) mettent en jeu des ions qui réagissent dès leur rencontre ; les réactions (b) et (c) exigent la rupture de liaisons covalentes (liaison O--O dans \ce{S2O8^2-}, liaison O=O dans \ce{O2}), ce qui demande du temps.
\end{exoresolu}

\begin{exercice}[À toi de jouer]
Pour chacune des transformations suivantes, indique si elle est rapide, lente ou très lente, et propose un moyen pratique d'en modifier la vitesse : (1) la fermentation du vin de palme ; (2) la formation d'un précipité de sulfate de baryum ; (3) la corrosion d'un portail en fer à Douala ; (4) la saponification de l'huile de palme par la soude.
\end{exercice}

\begin{corrige}[Exercice 1]
(1) Lente (heures à jours) : on l'accélère à la chaleur, on la ralentit au frais. (2) Rapide (instantanée) : réaction entre ions en solution (\ce{Ba^2+ + SO4^2- -> BaSO4}). (3) Très lente (mois, années) : on la ralentit en peignant ou galvanisant le fer ; l'air humide et salé de Douala l'accélère. (4) Lente (dizaines de minutes) : on l'accélère en chauffant et en agitant, comme dans les savonneries artisanales.
\end{corrige}

Maintenant que tu sais distinguer les réactions selon leur durée, il reste à répondre à la question centrale : \emph{comment mesurer cette rapidité ?} C'est l'objet de la section suivante, où nous définirons rigoureusement les vitesses de formation et de disparition.

\coursec{Vitesses de formation et de disparition}

Dans la section précédente, nous avons classé les réactions chimiques selon leur durée : certaines, comme la précipitation de \ce{AgCl}, sont quasi instantanées ; d'autres, comme l'oxydation des ions iodure par les ions peroxodisulfate, se déroulent en plusieurs minutes ; d'autres encore, comme la formation de la rouille sur un portail à Douala, s'étalent sur des mois. Dire qu'une réaction est « lente » ou « rapide » reste toutefois qualitatif. Le physicien ou le chimiste veut un \cle{nombre} : il veut \emph{mesurer} la rapidité d'une transformation et prévoir comment elle évolue. C'est l'objet de cette section : définir rigoureusement la \cle{vitesse} de formation d'un produit et la vitesse de disparition d'un réactif, et apprendre à les déterminer à partir d'une courbe expérimentale.

\courssub{La réaction de référence : formation du diiode}

Pour étudier une vitesse, il faut une réaction suffisamment lente pour être observée à l'œil, et dont l'avancement se \emph{voit}. La réaction entre les ions iodure \ce{I-} et les ions peroxodisulfate \ce{S2O8^2-} est idéale :

\[ \ce{2 I- + S2O8^2- -> I2 + 2 SO4^2-} \]

Commentons cette équation. Les ions iodure (apportés par exemple par l'iodure de potassium \ce{KI}) sont oxydés en diiode \ce{I2} ; les ions peroxodisulfate (apportés par le peroxodisulfate de potassium \ce{K2S2O8}) sont réduits en ions sulfate \ce{SO4^2-}. Il s'agit d'une réaction d'\cle{oxydoréduction} (couples \ce{S2O8^2-}/\ce{SO4^2-} et \ce{I2}/\ce{I-}), \textbf{lente} à température ambiante, et \textbf{totale}.

Pourquoi est-elle si commode à suivre ? Parce que le diiode formé colore progressivement le mélange en \cle{brun} (jaune-orangé au début, puis brun foncé), alors que tous les réactifs et les ions sulfate sont incolores. L'intensité de la coloration est donc une image directe de la quantité de diiode formé : en mesurant cette coloration au spectrophotomètre à différentes dates, on obtient la concentration $[\ce{I2}]$ en fonction du temps $t$.

\begin{experience}[Observation au laboratoire]
Dans un bécher, on mélange \SI{50}{mL} de solution d'iodure de potassium de concentration \SI{0,20}{mol.L^{-1}} et \SI{50}{mL} de solution de peroxodisulfate de potassium de concentration \SI{0,10}{mol.L^{-1}}. On déclenche le chronomètre au moment du mélange. \textbf{Observation} : la coloration brune apparaît lentement et s'intensifie pendant plusieurs dizaines de minutes. \textbf{Interprétation} : le diiode se forme progressivement ; la réaction est lente, ce qui laisse le temps de prélever des échantillons et de mesurer $[\ce{I2}]$ à différentes dates.
\end{experience}

\courssub{Vitesse moyenne de formation d'un produit}

Commençons par l'intuition. Si un automobiliste parcourt \SI{100}{km} en \SI{2}{h}, sa vitesse moyenne est \SI{50}{km.h^{-1}} : c'est la distance parcourue divisée par la durée. En cinétique, la « distance parcourue » est la variation de concentration du produit formé.

\begin{definition}[Vitesse moyenne de formation]
La \cle{vitesse moyenne de formation} d'un produit $P$ entre les dates $t_1$ et $t_2$ est :
\[ v_{\text{moy}}(P) = \frac{[P]_2 - [P]_1}{t_2 - t_1} = \frac{\Delta [P]}{\Delta t} \]
Elle s'exprime en \si{mol.L^{-1}.s^{-1}} (ou en \si{mol.L^{-1}.min^{-1}}, \si{mmol.L^{-1}.min^{-1}} selon l'échelle de temps de la réaction).
\end{definition}

Graphiquement, $v_{\text{moy}}$ est le \textbf{coefficient directeur de la corde} (segment de droite) joignant les points de la courbe $[P] = f(t)$ d'abscisses $t_1$ et $t_2$. Une vitesse moyenne dépend donc de l'intervalle choisi : elle ne renseigne pas sur ce qui se passe à une date précise.

\begin{exemplebox}[Calcul de vitesse moyenne]
Au cours du suivi de la réaction \ce{2 I- + S2O8^2- -> I2 + 2 SO4^2-}, on mesure $[\ce{I2}] = \SI{0}{mmol.L^{-1}}$ à $t_1 = 0$ et $[\ce{I2}] = \SI{4,0}{mmol.L^{-1}}$ à $t_2 = \SI{8}{min}$. La vitesse moyenne de formation du diiode entre ces deux dates vaut :
\[ v_{\text{moy}}(\ce{I2}) = \frac{4{,}0 - 0}{8 - 0} = \SI{0,50}{mmol.L^{-1}.min^{-1}} \]
Entre 0 et 8 minutes, la concentration en diiode augmente \emph{en moyenne} de \SI{0,50}{mmol.L^{-1}} chaque minute.
\end{exemplebox}

\courssub{Vitesse instantanée de formation : la tangente à la courbe}

Revenons à l'automobiliste : sa vitesse moyenne sur \SI{2}{h} ne dit rien de la valeur affichée par son compteur à un instant précis. De même, la vitesse moyenne entre $t_1$ et $t_2$ ne dit pas à quelle allure le diiode se forme \emph{exactement} à la date $t$. Pour cela, on définit la vitesse instantanée.

L'idée est de resserrer l'intervalle $[t_1, t_2]$ autour de la date $t$ : la corde se rapproche alors de la \cle{tangente} à la courbe au point considéré. À la limite, la pente de la corde devient la pente de la tangente.

\begin{definition}[Vitesse instantanée de formation]
La \cle{vitesse instantanée de formation} d'un produit $P$ à la date $t$ est le coefficient directeur de la tangente à la courbe $[P] = f(t)$ au point d'abscisse $t$ :
\[ v(P, t) = \left( \frac{\mathrm{d}[P]}{\mathrm{d}t} \right)_t \]
Elle s'exprime en \si{mol.L^{-1}.s^{-1}} (ou \si{mmol.L^{-1}.min^{-1}}).
\end{definition}

\begin{methode}[Déterminer graphiquement une vitesse instantanée]
Pour déterminer $v_{\text{inst}}$ à la date $t_0$ à partir de la courbe expérimentale :
\begin{enumerate}
\item Repérer le point $M$ de la courbe d'abscisse $t_0$.
\item Tracer la \textbf{tangente} à la courbe en $M$ : poser la règle de façon qu'elle « effleure » la courbe en $M$ sans la traverser, en équilibrant les angles de part et d'autre.
\item Choisir sur cette tangente \textbf{deux points éloignés} $A$ et $B$ (plus ils sont éloignés, plus le calcul est précis) et relever leurs coordonnées $(t_A, [P]_A)$ et $(t_B, [P]_B)$.
\item Calculer la pente : $v_{\text{inst}} = \dfrac{[P]_B - [P]_A}{t_B - t_A}$.
\end{enumerate}
\end{methode}

\begin{attention}[Corde ou tangente ?]
Erreur fréquente au baccalauréat : confondre la pente de la \textbf{corde} (qui donne la vitesse \emph{moyenne} entre deux dates) avec la pente de la \textbf{tangente} (qui donne la vitesse \emph{instantanée} à une date). Lis bien la question : « entre $t_1$ et $t_2$ » appelle une corde ; « à la date $t$ » appelle une tangente.
\end{attention}

\begin{popfigure}
\begin{center}
\begin{tikzpicture}
\begin{axis}[
  width=0.92\linewidth, height=7cm,
  xlabel={$t$ (min)}, ylabel={$[\ce{I2}]$ (mmol.L$^{-1}$)},
  xmin=0, xmax=30, ymin=0, ymax=11,
  xtick={0,5,10,15,20,25,30}, ytick={0,2,4,6,8,10},
  grid=both, grid style={popline!40},
  axis lines=left, tick label style={font=\small}, label style={font=\small}
]
% Courbe [I2] = 10(1 - exp(-t/12))
\addplot[popcurve,domain=0:30,samples=200]{10*(1-exp(-x/12))};
% Tangente en t=10 : pente = 10/12*exp(-10/12) ≈ 0,3624 ; f(10) ≈ 5,654
\addplot[poporange, thick, domain=0:22, samples=2]{5.654 + 0.3624*(x-10)};
% Point M
\addplot[only marks, mark=*, popdark, mark size=2pt] coordinates {(10,5.654)};
\node[popdark, font=\small, above right] at (axis cs:10,5.654) {$M$};
% Triangle de pente
\draw[popgreen, thick, dashed] (axis cs:4,3.48) -- (axis cs:18,8.55);
\draw[popgreen, thick, dashed] (axis cs:4,3.48) -- (axis cs:18,3.48);
\draw[popgreen, thick, dashed] (axis cs:18,3.48) -- (axis cs:18,8.55);
\node[popgreen, font=\small, below] at (axis cs:11,3.48) {$\Delta t$};
\node[popgreen, font=\small, right] at (axis cs:18,5.9) {$\Delta[\ce{I2}]$};
\node[poporange, font=\small, rotate=14, anchor=south] at (axis cs:16,7.4) {tangente en $t = 10$ min};
\end{axis}
\end{tikzpicture}
\end{center}
\legende{Concentration du diiode formé en fonction du temps. La vitesse instantanée à $t = \SI{10}{min}$ est le coefficient directeur de la tangente (orange), calculé grâce au triangle de pente (vert) : $v = \Delta[\ce{I2}]/\Delta t$.}
\end{popfigure}

\begin{exoresolu}[Vitesse moyenne puis vitesse instantanée du diiode]
\textbf{Énoncé.} Le suivi colorimétrique de la réaction \ce{2 I- + S2O8^2- -> I2 + 2 SO4^2-} donne les résultats suivants :

\begin{center}
\begin{tabularx}{0.85\linewidth}{|l|*{7}{>{\centering\arraybackslash}X|}}
\hline
\rowcolor{popblueL}
$t$ (min) & 0 & 5 & 10 & 15 & 20 & 25 & 30 \\
\hline
$[\ce{I2}]$ (mmol.L$^{-1}$) & 0 & 3,4 & 5,7 & 7,1 & 8,1 & 8,8 & 9,2 \\
\hline
\end{tabularx}
\end{center}

\begin{enumerate}
\item Calculer la vitesse moyenne de formation de \ce{I2} entre $t = 0$ et $t = \SI{5}{min}$.
\item Sur la courbe tracée, la tangente au point d'abscisse $t = \SI{10}{min}$ passe par les points $A(4\ \text{min}\ ;\ 3{,}5\ \text{mmol.L}^{-1})$ et $B(18\ \text{min}\ ;\ 8{,}6\ \text{mmol.L}^{-1})$. Calculer la vitesse instantanée à $t = \SI{10}{min}$.
\item Comparer et conclure.
\end{enumerate}
\tcblower
\textbf{Solution.}
\begin{enumerate}
\item Vitesse moyenne entre 0 et 5 min :
\[ v_{\text{moy}} = \frac{[\ce{I2}]_5 - [\ce{I2}]_0}{5 - 0} = \frac{3{,}4 - 0}{5} = \SI{0,68}{mmol.L^{-1}.min^{-1}} \]
\item Vitesse instantanée à $t = \SI{10}{min}$, pente de la tangente :
\[ v_{\text{inst}} = \frac{[\ce{I2}]_B - [\ce{I2}]_A}{t_B - t_A} = \frac{8{,}6 - 3{,}5}{18 - 4} = \frac{5{,}1}{14} \approx \SI{0,36}{mmol.L^{-1}.min^{-1}} \]
\item On constate que $v_{\text{inst}}(10\ \text{min}) < v_{\text{moy}}(0 \to 5\ \text{min})$ : la réaction \textbf{ralentit} au cours du temps. C'est un résultat général que nous expliquerons plus bas : les réactifs s'épuisent, leurs concentrations diminuent, et la vitesse diminue avec elles.
\end{enumerate}
\end{exoresolu}

\courssub{Vitesse de disparition d'un réactif}

Jusqu'ici, nous avons suivi l'apparition du diiode. Mais on peut tout aussi bien suivre la \emph{disparition} d'un réactif, par exemple l'ion peroxodisulfate \ce{S2O8^2-}. Sa concentration $[\ce{S2O8^2-}]$ \textbf{décroît} au cours du temps : la courbe $[\ce{S2O8^2-}] = g(t)$ est décroissante et ses tangentes ont des coefficients directeurs \emph{négatifs}.

Or une vitesse, par convention, doit être une grandeur \textbf{positive}. On introduit donc un signe moins.

\begin{definition}[Vitesse de disparition d'un réactif]
La \cle{vitesse moyenne de disparition} d'un réactif $R$ entre $t_1$ et $t_2$ est :
\[ v_{\text{moy}}(R) = -\frac{[R]_2 - [R]_1}{t_2 - t_1} = -\frac{\Delta [R]}{\Delta t} \]
La \cle{vitesse instantanée de disparition} à la date $t$ est l'opposé du coefficient directeur de la tangente à la courbe $[R] = g(t)$ :
\[ v(R, t) = -\left( \frac{\mathrm{d}[R]}{\mathrm{d}t} \right)_t \]
Comme $[R]$ diminue, $\Delta[R] < 0$ et le signe moins garantit une vitesse positive.
\end{definition}

\begin{attention}[Le signe moins est obligatoire]
Erreur classique : oublier le signe moins dans la vitesse de disparition et trouver une vitesse \emph{négative}. Une vitesse de réaction est \textbf{toujours positive}. Si tu obtiens un résultat négatif, vérifie immédiatement ton signe. Astuce : pour un réactif, on calcule souvent $\dfrac{[R]_1 - [R]_2}{t_2 - t_1}$ (le « grand moins le petit »), ce qui revient au même.
\end{attention}

\begin{popfigure}
\begin{center}
\begin{tikzpicture}
\begin{axis}[
  width=0.92\linewidth, height=7cm,
  xlabel={$t$ (min)}, ylabel={$[\ce{S2O8^2-}]$ (mmol.L$^{-1}$)},
  xmin=0, xmax=30, ymin=0, ymax=11,
  xtick={0,5,10,15,20,25,30}, ytick={0,2,4,6,8,10},
  grid=both, grid style={popline!40},
  axis lines=left, tick label style={font=\small}, label style={font=\small}
]
% Courbe décroissante : 10*exp(-t/12)
\addplot[popcurve,domain=0:30,samples=200]{10*exp(-x/12)};
% Tangente en t=10 : f(10)=4,346 ; pente = -10/12*exp(-10/12) = -0,3624
\addplot[poporange, thick, domain=0:24, samples=2]{4.346 - 0.3624*(x-10)};
\addplot[only marks, mark=*, popdark, mark size=2pt] coordinates {(10,4.346)};
\node[popdark, font=\small, above left] at (axis cs:10,4.346) {$M$};
% Triangle
\draw[popgreen, thick, dashed] (axis cs:4,6.52) -- (axis cs:16,2.17);
\draw[popgreen, thick, dashed] (axis cs:4,6.52) -- (axis cs:16,6.52);
\draw[popgreen, thick, dashed] (axis cs:16,6.52) -- (axis cs:16,2.17);
\node[popgreen, font=\small, above] at (axis cs:10,6.52) {$\Delta t$};
\node[popgreen, font=\small, right] at (axis cs:16,4.3) {$\Delta[\ce{S2O8^2-}] < 0$};
\node[poporange, font=\small, rotate=-14, anchor=south] at (axis cs:17,2.6) {tangente en $t = 10$ min};
\end{axis}
\end{tikzpicture}
\end{center}
\legende{Disparition des ions peroxodisulfate. La tangente a une pente négative ; la vitesse de disparition est son opposé : $v = -\Delta[\ce{S2O8^2-}]/\Delta t > 0$.}
\end{popfigure}

\begin{exemplebox}[Vitesse moyenne de disparition]
Entre $t_1 = \SI{5}{min}$ et $t_2 = \SI{15}{min}$, la concentration en ions peroxodisulfate passe de \SI{6,6}{mmol.L^{-1}} à \SI{2,9}{mmol.L^{-1}}. La vitesse moyenne de disparition vaut :
\[ v_{\text{moy}}(\ce{S2O8^2-}) = -\frac{2{,}9 - 6{,}6}{15 - 5} = -\frac{-3{,}7}{10} = \SI{0,37}{mmol.L^{-1}.min^{-1}} \]
Le résultat est bien positif.
\end{exemplebox}

\courssub{Lien stœchiométrique entre vitesses de formation et de disparition}

Regardons à nouveau l'équation :
\[ \ce{2 I- + S2O8^2- -> I2 + 2 SO4^2-} \]
Chaque fois qu'\textbf{un} ion \ce{S2O8^2-} disparaît, \textbf{un} diiode \ce{I2} apparaît : les vitesses de disparition de \ce{S2O8^2-} et de formation de \ce{I2} sont donc égales. En revanche, il faut \textbf{deux} ions \ce{I-} pour former un \ce{I2} : les ions iodure disparaissent deux fois plus vite que le diiode n'apparaît.

\begin{propriete}[Relations entre vitesses]
Pour une réaction d'équation $\ce{a A + b B -> c C + d D}$, les vitesses sont liées par les coefficients stœchiométriques :
\[ \frac{v(A)}{a} = \frac{v(B)}{b} = \frac{v(C)}{c} = \frac{v(D)}{d} \]
Pour la réaction de référence : $v(\ce{I2}) = v(\ce{S2O8^2-}) = \dfrac{v(\ce{I-})}{2} = \dfrac{v(\ce{SO4^2-})}{2}$.
\end{propriete}

\begin{exemplebox}[Application]
Si à une date donnée le diiode se forme à $v(\ce{I2}) = \SI{0,36}{mmol.L^{-1}.min^{-1}}$, alors les ions iodure disparaissent à :
\[ v(\ce{I-}) = 2 \times v(\ce{I2}) = \SI{0,72}{mmol.L^{-1}.min^{-1}} \]
\end{exemplebox}

\begin{attention}[Diviser ou multiplier ?]
Piège fréquent : se tromper de sens dans le facteur stœchiométrique. Pose-toi la question concrète : « pour fabriquer 1 mol de \ce{I2}, combien de mol de \ce{I-} consomme-t-on ? ». Réponse : 2. Donc \ce{I-} disparaît \emph{plus vite} que \ce{I2} n'apparaît : $v(\ce{I-}) = 2\,v(\ce{I2})$. La vérification par le bon sens évite bien des erreurs.
\end{attention}

\courssub{Comment évolue la vitesse au cours du temps ?}

Observons nos courbes : au début, elles sont très inclinées (tangentes raides) ; puis elles s'aplatissent progressivement jusqu'à devenir horizontales lorsque la réaction s'arrête. La pente de la tangente — donc la \cle{vitesse instantanée} — \textbf{diminue continuellement} au cours du temps.

\begin{propriete}[Évolution de la vitesse]
Au cours d'une transformation chimique, la vitesse instantanée est \textbf{maximale au début} ($t = 0$) puis \textbf{diminue} au fur et à mesure que les réactifs sont consommés, car la vitesse dépend des concentrations des réactifs : moins il reste de réactifs, plus la réaction est lente. La vitesse s'annule lorsque la réaction est terminée (réactif limitant épuisé) ou lorsque l'équilibre est atteint.
\end{propriete}

C'est une observation que tu connais déjà sans le savoir : le feu de bois flambe vigoureusement au début, quand le combustible est abondant, puis faiblit à mesure que les braises se consomment. Nous étudierons en détail ces \emph{facteurs cinétiques} — concentration, température, catalyseur — dans la section 5.

\begin{aretenir}[L'essentiel de la section]
\begin{itemize}
\item \textbf{Vitesse moyenne de formation} d'un produit : $v_{\text{moy}} = \dfrac{[P]_2 - [P]_1}{t_2 - t_1}$ — pente de la \textbf{corde}.
\item \textbf{Vitesse instantanée de formation} à la date $t$ : coefficient directeur de la \textbf{tangente} à la courbe $[P] = f(t)$ ; on la calcule avec deux points éloignés de la tangente.
\item \textbf{Vitesse de disparition} d'un réactif : $v = -\dfrac{\mathrm{d}[R]}{\mathrm{d}t}$ ; le signe moins garantit une vitesse \textbf{positive}.
\item Les vitesses des différentes espèces sont liées par les \textbf{coefficients stœchiométriques} : $\dfrac{v(A)}{a} = \dfrac{v(C)}{c}$.
\item La vitesse instantanée \textbf{diminue au cours du temps} car les concentrations des réactifs diminuent.
\end{itemize}
\end{aretenir}

\begin{exercice}[À toi de jouer]
Lors du suivi de la réaction \ce{2 I- + S2O8^2- -> I2 + 2 SO4^2-}, on relève $[\ce{I2}] = \SI{7,1}{mmol.L^{-1}}$ à $t = \SI{15}{min}$ et $[\ce{I2}] = \SI{8,8}{mmol.L^{-1}}$ à $t = \SI{25}{min}$.
\begin{enumerate}
\item Calculer la vitesse moyenne de formation du diiode entre ces deux dates.
\item En déduire la vitesse moyenne de disparition des ions iodure sur le même intervalle.
\item La vitesse instantanée à $t = \SI{20}{min}$ est-elle supérieure, égale ou inférieure à la valeur trouvée à la question 1 ? Justifier sans calcul.
\end{enumerate}
\end{exercice}

\begin{corrige}[Exercice : À toi de jouer]
\begin{enumerate}
\item $v_{\text{moy}}(\ce{I2}) = \dfrac{8{,}8 - 7{,}1}{25 - 15} = \dfrac{1{,}7}{10} = \SI{0,17}{mmol.L^{-1}.min^{-1}}$.
\item D'après la stœchiométrie, $v(\ce{I-}) = 2\,v(\ce{I2}) = \SI{0,34}{mmol.L^{-1}.min^{-1}}$.
\item Inférieure : la vitesse instantanée diminue au cours du temps car les réactifs s'épuisent ; la pente de la tangente à $t = \SI{20}{min}$ est plus faible que la pente moyenne sur l'intervalle $[15\ ;\ 25]$ min, qui intègre des dates plus précoces où la réaction est plus rapide.
\end{enumerate}
\end{corrige}

\coursec{Vitesse volumique, avancement et temps de demi-réaction}

Dans la section précédente, nous avons appris à mesurer la vitesse de formation d'un produit (comme le diiode \ce{I2}) ou la vitesse de disparition d'un réactif (comme les ions \ce{S2O8^2-}) à partir des courbes de concentration. Mais une question se pose naturellement : si deux expérimentateurs, l'un à Douala avec \SI{500}{mL} de mélange, l'autre à Yaoundé avec \SI{50}{mL} du même mélange, étudient la même réaction, obtiendront-ils les mêmes vitesses ? Pas nécessairement, car les quantités de matière mises en jeu diffèrent ! Les chimistes ont donc défini une grandeur plus fondamentale, qui caractérise la réaction elle-même et non l'expérience particulière : la \cle{vitesse volumique de réaction}. C'est elle que nous allons construire pas à pas, avant de découvrir une grandeur remarquablement pratique : le \cle{temps de demi-réaction}.

\courssub{De l'avancement à la vitesse volumique}

\textbf{L'avancement : la variable universelle de la réaction}\par

Considérons la réaction entre les ions iodure et les ions peroxodisulfate, réaction lente que nous suivons depuis le début de cette leçon :

\[ \ce{2I- (aq) + S2O8^2- (aq) -> I2 (aq) + 2SO4^2- (aq)} \]

Cette réaction est lente à température ambiante, ce qui permet de la suivre confortablement ; elle est totale et ne nécessite pas de catalyseur (même si les ions \ce{Fe^3+} l'accélèrent, comme nous le verrons dans la section 5).

Pour décrire l'état d'avancement de cette réaction, on utilise l'\cle{avancement} $x$, exprimé en mol : c'est la quantité de matière de \ce{S2O8^2-} disparue, ou de \ce{I2} apparue, depuis le début. Le tableau d'avancement relie $x$ à toutes les quantités de matière présentes à chaque instant.

\begin{popfigure}
\begin{center}
\begin{tikzpicture}[font=\small]
\draw[rounded corners=2pt, fill=popcream, draw=popline] (-7.4,-3.3) rectangle (7.4,2.9);
\node[font=\bfseries, popdark] at (0,2.5) {Tableau d'avancement : \ce{2I- + S2O8^2- -> I2 + 2SO4^2-}};
\node[align=left] at (-4.6,1.7) {État initial ($t=0$)};
\node[align=left] at (-4.6,0.9) {État intermédiaire ($t$)};
\node[align=left] at (-4.6,0.1) {État final ($t_f$)};
\node[popblue] at (-1.2,1.7) {$n_0(\ce{I-})$};
\node[popblue] at (0.9,1.7) {$n_0(\ce{S2O8^2-})$};
\node[popgreen] at (3.1,1.7) {$0$};
\node[popgreen] at (5.5,1.7) {$0$};
\node[popblue] at (-1.2,0.9) {$n_0(\ce{I-}) - 2x$};
\node[popblue] at (0.9,0.9) {$n_0(\ce{S2O8^2-}) - x$};
\node[popgreen] at (3.1,0.9) {$x$};
\node[popgreen] at (5.5,0.9) {$2x$};
\node[poporange] at (-1.2,0.1) {$n_0(\ce{I-}) - 2x_f$};
\node[poporange] at (0.9,0.1) {$n_0(\ce{S2O8^2-}) - x_f$};
\node[poporange] at (3.1,0.1) {$x_f$};
\node[poporange] at (5.5,0.1) {$2x_f$};
\draw[popline] (-7.0,-0.4) -- (7.0,-0.4);
\node[align=center, font=\footnotesize\itshape, popmuted] at (0,-1.0) {Pour une réaction \textbf{totale} : $x_f = x_{\max}$, atteint quand le réactif limitant est épuisé.\\ Pour une réaction \textbf{limitée} (estérification) : $x_f < x_{\max}$, l'équilibre est atteint avant épuisement.};
\node[align=center, font=\footnotesize, popdark] at (0,-2.3) {Lecture clé : $n(\ce{I2}) = x$ \quad et \quad $[\ce{I2}] = \dfrac{x}{V}$ \quad (avec $V$ = volume de la solution)};
\end{tikzpicture}
\end{center}
\legende{Le tableau d'avancement relie l'avancement $x$ aux quantités de matière de toutes les espèces. La dernière ligne distingue réaction totale et réaction limitée.}
\end{popfigure}

\textbf{Pourquoi diviser par le volume ?}\par

Imaginons que tu suives la formation du diiode dans un bécher de \SI{100}{mL} : la quantité $x$ augmente au cours du temps. Si ton voisin réalise exactement la même expérience mais dans \SI{1}{L} de solution (avec dix fois plus de réactifs), sa quantité $x$ augmentera environ dix fois plus vite... pourtant c'est la même réaction, à la même température, avec les mêmes concentrations ! Pour comparer les deux expériences, il faut raisonner \emph{par litre de solution}, c'est-à-dire sur les concentrations et non sur les quantités de matière.

\begin{definition}[Vitesse volumique de réaction]
La \cle{vitesse volumique de réaction} $v$ est la dérivée de l'avancement $x$ par rapport au temps, divisée par le volume $V$ de la solution :
\[ v = \frac{1}{V}\,\frac{\mathrm{d}x}{\mathrm{d}t} \]
Elle s'exprime en \si{mol.L^{-1}.s^{-1}} (ou en \si{mol.L^{-1}.min^{-1}} si le temps est en minutes). C'est une grandeur toujours \textbf{positive}, qui décroît généralement au cours du temps car les réactifs s'épuisent.
\end{definition}

\begin{aretenir}[Lien avec les vitesses de formation et de disparition]
Pour une équation-bilan $\ce{aA + bB -> cC + dD}$, la vitesse volumique est liée aux vitesses de disparition des réactifs et de formation des produits par les coefficients stœchiométriques :
\[ v = -\frac{1}{a}\frac{\mathrm{d}[\ce{A}]}{\mathrm{d}t} = -\frac{1}{b}\frac{\mathrm{d}[\ce{B}]}{\mathrm{d}t} = +\frac{1}{c}\frac{\mathrm{d}[\ce{C}]}{\mathrm{d}t} = +\frac{1}{d}\frac{\mathrm{d}[\ce{D}]}{\mathrm{d}t} \]
Le signe « $-$ » devant les réactifs garantit que $v$ reste positive (leur concentration diminue). Pour notre réaction de référence :
\[ v = \frac{\mathrm{d}[\ce{I2}]}{\mathrm{d}t} = -\frac{\mathrm{d}[\ce{S2O8^2-}]}{\mathrm{d}t} = -\frac{1}{2}\,\frac{\mathrm{d}[\ce{I-}]}{\mathrm{d}t} \]
\end{aretenir}

\begin{attention}[Le piège des coefficients stœchiométriques]
L'erreur la plus fréquente au Baccalauréat est d'oublier de diviser par le coefficient stœchiométrique. Ici, la vitesse de disparition des ions iodure est \emph{deux fois} la vitesse volumique : $-\dfrac{\mathrm{d}[\ce{I-}]}{\mathrm{d}t} = 2v$, car deux ions \ce{I-} disparaissent pour chaque mole d'avancement. Vérifie toujours l'équation-bilan avant d'écrire la relation !
\end{attention}

\begin{exemplebox}[Calcul d'une vitesse volumique]
On mélange $V_1 = \SI{20,0}{mL}$ d'une solution d'iodure de potassium à \SI{0,20}{mol.L^{-1}} avec $V_2 = \SI{20,0}{mL}$ d'une solution de peroxodisulfate de potassium à \SI{0,10}{mol.L^{-1}}. Le volume total est $V = \SI{40,0}{mL}$. À l'instant $t = \SI{4}{min}$, on mesure $[\ce{I2}] = \SI{2,0e-3}{mol.L^{-1}}$ et la tangente à la courbe a pour coefficient directeur \SI{8,0e-4}{mol.L^{-1}.min^{-1}}.

La vitesse volumique à cet instant vaut :
\[ v = \frac{\mathrm{d}[\ce{I2}]}{\mathrm{d}t} = \SI{8,0e-4}{mol.L^{-1}.min^{-1}} = \SI{1,3e-5}{mol.L^{-1}.s^{-1}} \]

La vitesse de disparition des ions iodure au même instant est le double : \SI{1,6e-3}{mol.L^{-1}.min^{-1}}. Remarque que ce résultat ne dépend pas du volume $V$ : un autre élève travaillant avec \SI{400}{mL} du même mélange trouverait exactement la même valeur. C'est tout l'intérêt de la vitesse volumique : elle caractérise \textbf{la réaction elle-même}, dans les conditions de l'expérience.
\end{exemplebox}

\courssub{Le temps de demi-réaction}

\textbf{Une intuition d'abord}\par

Pour comparer la rapidité de deux réactions, les chimistes aiment utiliser un repère unique, comme le « chrono » d'un coureur. Plutôt que d'attendre la fin complète de la réaction (qui peut être très longue, voire infinie pour une réaction limitée), on regarde simplement : \emph{en combien de temps la réaction a-t-elle fait la moitié du chemin ?} Ce temps, c'est le temps de demi-réaction.

\begin{definition}[Temps de demi-réaction]
Le \cle{temps de demi-réaction}, noté $t_{1/2}$, est la durée nécessaire pour que l'avancement atteigne la moitié de sa valeur finale :
\[ x(t_{1/2}) = \frac{x_f}{2} \]
Pour une réaction totale, $x_f = x_{\max}$ et l'on écrit $x(t_{1/2}) = \dfrac{x_{\max}}{2}$. De manière équivalente, pour le diiode formé : $[\ce{I2}](t_{1/2}) = \dfrac{[\ce{I2}]_f}{2}$.
\end{definition}

\begin{methode}[Déterminer graphiquement le temps de demi-réaction]
\begin{enumerate}
\item Déterminer la valeur finale : $x_f$ (ou $[\ce{I2}]_f$), c'est l'ordonnée de l'asymptote horizontale de la courbe.
\item Calculer la moitié : $x_f/2$ (ou $[\ce{I2}]_f/2$).
\item Reporter cette valeur sur l'axe des ordonnées et tracer une horizontale en pointillés jusqu'à la courbe.
\item Du point d'intersection, abaisser une verticale en pointillés jusqu'à l'axe des abscisses.
\item Lire l'abscisse correspondante : c'est $t_{1/2}$.
\end{enumerate}
\end{methode}

\begin{popfigure}
\begin{center}
\begin{tikzpicture}
\begin{axis}[
  width=0.85\linewidth, height=7cm,
  xlabel={$t$ (min)}, ylabel={$[\ce{I2}]$ (mmol.L$^{-1}$)},
  xmin=0, xmax=30, ymin=0, ymax=9.5,
  xtick={0,5,10,15,20,25,30}, ytick={0,2,4,6,8},
  grid=major, grid style={popline!40},
  axis lines=left, tick label style={font=\small}, label style={font=\small}
]
\addplot[popcurve, domain=0:30, samples=200] {8*(1-exp(-x/8.66))};
\addplot[dashed, poporange, thick] coordinates {(0,4) (6,4)};
\addplot[dashed, poporange, thick] coordinates {(6,0) (6,4)};
\addplot[dotted, popmuted, thick] coordinates {(0,8) (30,8)};
\node[popmuted, font=\footnotesize, anchor=west] at (axis cs:20,8.4) {$[\ce{I2}]_{\max} = 8{,}0$};
\node[poporange, font=\footnotesize, anchor=west] at (axis cs:0.5,4.4) {$\dfrac{[\ce{I2}]_{\max}}{2} = 4{,}0$};
\node[poporange, font=\footnotesize, anchor=north] at (axis cs:6,-0.4) {$t_{1/2} = 6$ min};
\addplot[only marks, mark=*, mark size=2pt, poporange] coordinates {(6,4)};
\end{axis}
\end{tikzpicture}
\end{center}
\legende{Lecture graphique du temps de demi-réaction sur la courbe $[\ce{I2}] = f(t)$ : on repère la moitié de la valeur finale sur l'axe des ordonnées, puis on lit l'abscisse correspondante.}
\end{popfigure}

\begin{propriete}[$t_{1/2}$, jauge de rapidité]
Le temps de demi-réaction est une \textbf{grandeur caractéristique de la rapidité} d'une réaction dans des conditions données :
\begin{itemize}
\item plus $t_{1/2}$ est \textbf{petit}, plus la réaction est \textbf{rapide} ;
\item plus $t_{1/2}$ est \textbf{grand}, plus la réaction est \textbf{lente}.
\end{itemize}
Ainsi, la précipitation \ce{Ag+ + Cl- -> AgCl} a un $t_{1/2}$ de l'ordre de la fraction de seconde (réaction rapide), la réaction \ce{I-}/\ce{S2O8^2-} un $t_{1/2}$ de quelques minutes (réaction lente), et la formation de la rouille un $t_{1/2}$ de plusieurs mois ou années (réaction très lente).
\end{propriete}

\begin{attention}[Réaction totale ou limitée : quel dénominateur ?]
Erreur classique au Baccalauréat : pour une réaction \textbf{limitée} comme l'estérification, l'avancement final $x_f$ est inférieur à l'avancement maximal théorique $x_{\max}$ (l'équilibre est atteint avant épuisement du réactif limitant). La définition impose alors :
\[ x(t_{1/2}) = \frac{x_f}{2} \quad \text{et non pas} \quad \frac{x_{\max}}{2} \]
Retiens : on prend toujours la moitié de la valeur \emph{réellement atteinte} à l'état final, lue sur l'asymptote de la courbe expérimentale.
\end{attention}

\begin{exoresolu}[Temps de demi-réaction de la réaction \ce{I-}/\ce{S2O8^2-}]
\textbf{Énoncé.} On suit par spectrophotométrie la formation du diiode lors de la réaction \ce{2I- + S2O8^2- -> I2 + 2SO4^2-} à \SI{25}{\celsius}. La courbe $[\ce{I2}] = f(t)$ donne les résultats suivants :

\begin{center}
\begin{tabularx}{0.9\linewidth}{|l|*{7}{>{\centering\arraybackslash}X|}}
\hline
\rowcolor{popblueL} $t$ (min) & 0 & 2 & 6 & 10 & 16 & 24 & 40 \\
\hline
$[\ce{I2}]$ (mmol.L$^{-1}$) & 0 & 1,6 & 4,0 & 5,5 & 6,8 & 7,5 & 8,0 \\
\hline
\end{tabularx}
\end{center}

\begin{enumerate}
\item Déterminer la concentration finale $[\ce{I2}]_{\max}$ en justifiant que la réaction est totale.
\item Déterminer graphiquement (ou par lecture du tableau) le temps de demi-réaction.
\item La même expérience réalisée à \SI{45}{\celsius} donne $t'_{1/2} = \SI{2}{min}$. Conclure.
\end{enumerate}

\tcblower

\textbf{Solution détaillée.}

\textbf{1.} La concentration en diiode se stabilise à \SI{8,0}{mmol.L^{-1}} : la courbe atteint une asymptote horizontale. Comme la réaction est totale, cette valeur limite correspond à $[\ce{I2}]_{\max} = \dfrac{x_{\max}}{V} = \SI{8,0}{mmol.L^{-1}}$.

\textbf{2.} On calcule la moitié de la valeur finale :
\[ \frac{[\ce{I2}]_{\max}}{2} = \frac{8{,}0}{2} = \SI{4,0}{mmol.L^{-1}} \]
On reporte \SI{4,0}{mmol.L^{-1}} sur l'axe des ordonnées ; la lecture du tableau (ou de la courbe) montre que cette valeur est atteinte à $t = \SI{6}{min}$. Donc :
\[ \boxed{t_{1/2} = \SI{6}{min}} \]

\textbf{3.} À \SI{45}{\celsius}, $t'_{1/2} = \SI{2}{min} < t_{1/2} = \SI{6}{min}$ : la moitié du chemin est parcourue trois fois plus vite. La température est donc bien un \cle{facteur cinétique} : son augmentation accélère la réaction. Nous étudierons ce point en détail dans la section 5.
\end{exoresolu}

\begin{savaistu}[Le temps de demi-réaction dans la vie quotidienne]
La notion de « temps de moitié » dépasse la chimie de laboratoire ! En pharmacie, la \emph{demi-vie} d'un médicament indique en combien de temps sa concentration dans le sang est divisée par deux : c'est elle qui détermine la fréquence des prises (toutes les 6 h, toutes les 12 h...). En archéologie, la demi-vie du carbone 14 (\SI{5730}{ans}) permet de dater les objets anciens. Et au quotidien, quand le jus de bissap préparé la veille fermente, ou quand le vin de palme se transforme peu à peu en vinaigre au contact de l'air, il existe aussi un temps caractéristique au bout duquel la moitié de la transformation est accomplie. La cinétique chimique est partout autour de toi !
\end{savaistu}

\begin{exercice}[À toi de jouer]
Lors de la décomposition lente de l'eau oxygénée selon $\ce{2H2O2 (aq) -> 2H2O (l) + O2 (g)}$, on mesure $[\ce{H2O2}]_0 = \SI{0,80}{mol.L^{-1}}$ et $[\ce{H2O2}] = \SI{0,40}{mol.L^{-1}}$ à $t = \SI{15}{min}$.
\begin{enumerate}
\item Exprimer la vitesse volumique de réaction en fonction de $[\ce{H2O2}]$ et de $[\ce{O2}]$.
\item En admettant que la réaction soit totale, déterminer le temps de demi-réaction.
\end{enumerate}
\end{exercice}

\begin{corrige}[Exercice 1]
\textbf{1.} En divisant par les coefficients stœchiométriques :
\[ v = -\frac{1}{2}\,\frac{\mathrm{d}[\ce{H2O2}]}{\mathrm{d}t} = +\frac{\mathrm{d}[\ce{O2}]}{\mathrm{d}t} \]
\textbf{2.} La réaction étant totale, la moitié de l'avancement final correspond à la disparition de la moitié du réactif : $[\ce{H2O2}](t_{1/2}) = \dfrac{0{,}80}{2} = \SI{0,40}{mol.L^{-1}}$. Cette concentration est atteinte à $t = \SI{15}{min}$, donc $t_{1/2} = \SI{15}{min}$.
\end{corrige}

\begin{aretenir}[L'essentiel de la section]
\begin{itemize}
\item La \cle{vitesse volumique} $v = \dfrac{1}{V}\,\dfrac{\mathrm{d}x}{\mathrm{d}t}$ caractérise la réaction elle-même, indépendamment du volume utilisé ; elle s'exprime en \si{mol.L^{-1}.s^{-1}} ou \si{mol.L^{-1}.min^{-1}}.
\item Elle est reliée aux vitesses de formation et de disparition par les coefficients stœchiométriques : $v = \dfrac{1}{\nu_i}\,\dfrac{\mathrm{d}[i]}{\mathrm{d}t}$ (avec un signe « $-$ » pour les réactifs).
\item Le \cle{temps de demi-réaction} $t_{1/2}$ est la durée au bout de laquelle $x = \dfrac{x_f}{2}$ ; il se lit graphiquement en reportant la moitié de la valeur finale (l'asymptote) sur la courbe.
\item Plus $t_{1/2}$ est petit, plus la réaction est rapide : c'est l'outil idéal pour comparer des réactions entre elles.
\item Pour une réaction limitée, on utilise $x_f$ (valeur expérimentale finale) et non $x_{\max}$ théorique.
\end{itemize}
\end{aretenir}

\coursec{Suivi temporel d'une réaction}

Dans la section précédente, nous avons défini la vitesse volumique de réaction à partir de la courbe $x = f(t)$ et introduit le temps de demi-réaction $t_{1/2}$. Mais une question pratique demeure : \emph{comment obtient-on cette courbe en laboratoire ?} Comment connaître, minute après minute, la quantité de diiode formé ou la quantité d'ions peroxodisulfate restant ? C'est tout l'objet du \cle{suivi temporel} d'une réaction chimique : mesurer, à différentes dates, la concentration (ou la quantité de matière) d'une espèce chimique du système, afin de tracer la courbe cinétique et d'en tirer les vitesses et le temps de demi-réaction.

\courssub{Objectif du suivi temporel}

\begin{definition}[Suivre l'évolution d'un système chimique]
\cle{Suivre temporellement} une réaction chimique, c'est déterminer à différentes dates $t$ la concentration $[X]$ d'une espèce chimique $X$ (réactif qui disparaît ou produit qui apparaît) présente dans le milieu réactionnel. On en déduit l'\cle{avancement} $x(t)$, puis on trace la courbe cinétique $[X] = f(t)$ ou $x = f(t)$.
\end{definition}

Pour cela, deux grandes familles de méthodes existent :

\begin{itemize}
\item les \cle{méthodes chimiques} : on prélève un échantillon du milieu à une date $t$, on \emph{bloque} la réaction dans cet échantillon, puis on \emph{dose} l'espèce recherchée ;
\item les \cle{méthodes physiques} : on mesure en continu une grandeur physique (absorbance, conductivité, pression) reliée à la concentration, \emph{sans prélever ni détruire} l'échantillon.
\end{itemize}

\courssub{Méthode chimique : prélèvements, trempe et dosage}

\textbf{Le principe.}\par
Imagine que tu suives la réaction lente entre les ions iodure et les ions peroxodisulfate :
\[ \ce{S2O8^2- (aq) + 2 I- (aq) -> 2 SO4^2- (aq) + I2 (aq)} \]
(réaction lente, totale, à température ambiante). Le diiode \ce{I2} formé colore progressivement le milieu en brun. Pour connaître $[\ce{I2}]$ à une date $t$, on prélève un volume $V_p$ du mélange réactionnel et on dose le diiode qu'il contient.

\textbf{Le problème.}\par
Pendant que tu effectues le dosage — ce qui prend plusieurs minutes — la réaction \emph{continue} dans le prélèvement ! Le résultat mesuré ne correspondrait donc pas à la date $t$ du prélèvement, mais à une date ultérieure. Il faut \emph{arrêter le temps chimique} dans l'échantillon : c'est le rôle de la \cle{trempe}.

\begin{definition}[La trempe]
La \cle{trempe} est un \textbf{refroidissement brutal} du prélèvement (par versement dans de l'eau glacée) accompagné d'une \textbf{dilution} importante. Elle provoque le \emph{quasi-arrêt} de la réaction dans le prélèvement, car :
\begin{itemize}
\item l'abaissement de la température ralentit fortement la réaction (la température est un facteur cinétique) ;
\item la dilution diminue la concentration des réactifs, ce qui ralentit encore la réaction (la concentration est un facteur cinétique).
\end{itemize}
La composition du prélèvement est ainsi \og figée \fg{} : le dosage peut être réalisé tranquillement, son résultat correspond bien à la date $t$ du prélèvement.
\end{definition}

\begin{methode}[Suivre une réaction par dosage : la démarche]
\begin{enumerate}
\item Préparer le mélange réactionnel et déclencher le chronomètre à l'instant du mélange ($t = 0$).
\item Préparer \textbf{à l'avance} une série de béchers contenant chacun un volume d'eau distillée glacée (avec quelques glaçons).
\item À des dates $t_1, t_2, \ldots$ régulièrement espacées, prélever rapidement un volume $V_p$ connu du milieu réactionnel à la pipette.
\item Verser \textbf{immédiatement} le prélèvement dans un bécher d'eau glacée : c'est la trempe.
\item Doser l'espèce recherchée dans chaque prélèvement trempé.
\item Reporter les résultats dans un tableau $t$ / concentration, puis tracer la courbe.
\end{enumerate}
\end{methode}

\begin{attention}[L'erreur classique : doser sans tremper]
Si tu doses le prélèvement \emph{sans trempe}, la réaction se poursuit pendant toute la durée du dosage : la quantité dosée est \textbf{trop grande} pour un produit (ou trop petite pour un réactif). Tous les points de la courbe sont alors faux, et les vitesses calculées sont surestimées. La trempe n'est pas un détail : c'est la condition de validité de toute la méthode chimique.
\end{attention}

\begin{popfigure}
\begin{tikzpicture}[scale=0.95, font=\small]
% Bécher réactionnel
\draw[thick, popdark] (-0.9,0) -- (-0.9,2.2) (0.9,0) -- (0.9,2.2);
\draw[thick, popdark] (-0.9,0) -- (0.9,0);
\fill[poporange!45] (-0.85,0.05) rectangle (0.85,1.3);
\node[popdark] at (0,-0.45) {Mélange réactionnel};
\node[popdark] at (0,-0.85) {\ce{S2O8^2-} + \ce{I-} (brunit)};
% Pipette
\draw[thick, popblue] (0.5,1.5) -- (2.6,3.4);
\draw[thick, popblue] (0.75,1.35) -- (2.85,3.25);
\fill[poporange!45] (0.62,1.42) -- (2.72,3.32) -- (2.85,3.25) -- (0.75,1.35) -- cycle;
\node[popblue, rotate=38] at (2.4,3.75) {prélèvement $V_p$ à la date $t$};
% Flèche vers bécher glacé
\draw[-{Stealth}, very thick, poppurple] (3.2,2.6) .. controls (4.0,2.2) .. (4.6,1.6);
\node[poppurple] at (4.1,2.9) {versement rapide};
% Bécher d'eau glacée
\draw[thick, popdark] (4.0,0) -- (4.0,2.0) (6.4,0) -- (6.4,2.0);
\draw[thick, popdark] (4.0,0) -- (6.4,0);
\fill[popblue!25] (4.05,0.05) rectangle (6.35,1.5);
% glaçons
\draw[popblue, fill=popblue!15] (4.4,1.15) rectangle (4.9,1.6);
\draw[popblue, fill=popblue!15, rotate=20] (5.6,1.1) rectangle (6.05,1.55);
\draw[popblue, fill=popblue!15, rotate=-15] (5.0,1.2) rectangle (5.45,1.62);
\node[popdark] at (5.2,-0.45) {Bécher d'eau glacée};
\node[popdark] at (5.2,-0.85) {\textbf{TREMPE} : réaction bloquée};
% Flèche vers dosage
\draw[-{Stealth}, very thick, popgreen!70!black] (6.7,1.0) -- (7.6,1.0);
\node[popgreen!60!black] at (7.15,1.35) {dosage};
% Erlenmeyer de dosage
\draw[thick, popdark] (7.8,0) -- (8.35,1.3) -- (8.35,1.9);
\draw[thick, popdark] (9.6,0) -- (9.05,1.3) -- (9.05,1.9);
\draw[thick, popdark] (7.8,0) -- (9.6,0);
\fill[poppink!35] (7.95,0.05) -- (9.45,0.05) -- (9.2,0.7) -- (8.2,0.7) -- cycle;
\node[popdark] at (8.7,-0.45) {Dosage de \ce{I2}};
\node[popdark] at (8.7,-0.85) {par \ce{S2O3^2-}};
\end{tikzpicture}
\legende{Protocole de suivi par dosage : le prélèvement est versé dans l'eau glacée (trempe) puis dosé.}
\end{popfigure}

\textbf{Exemple de dosage : le diiode par le thiosulfate.}\par
Le diiode formé est dosé par une solution de \cle{thiosulfate de sodium} $(2\,\ce{Na+} + \ce{S2O3^2-})$ de concentration $C$ connue. L'équation de la réaction de dosage, rapide et totale, est :
\[ \ce{I2 (aq) + 2 S2O3^2- (aq) -> 2 I- (aq) + S4O6^2- (aq)} \]
Couples mis en jeu : \ce{I2}/\ce{I-} et \ce{S4O6^2-}/\ce{S2O3^2-}. L'équivalence est repérée par la \textbf{décoloration} de la solution (disparition de la teinte brune de \ce{I2}), plus nettement encore en ajoutant quelques gouttes d'empois d'amidon (le complexe bleu-noir amidon-diiode disparaît à l'équivalence).

À l'équivalence : $n(\ce{S2O3^2-}) = 2\,n(\ce{I2})$, soit $C \cdot V_E = 2\,[\ce{I2}] \cdot V_p$, d'où :
\[ [\ce{I2}] = \frac{C \cdot V_E}{2\,V_p} \]

\begin{exoresolu}[Dosage du diiode formé au cours du temps]
On mélange \SI{50}{mL} d'une solution d'iodure de potassium à \SI{0,20}{mol.L^{-1}} et \SI{50}{mL} d'une solution de peroxodisulfate de potassium à \SI{0,10}{mol.L^{-1}}. À la date $t = \SI{21}{min}$, on prélève $V_p = \SI{10,0}{mL}$ du mélange, on le trempe, puis on dose le diiode par une solution de thiosulfate de sodium de concentration $C = \SI{2,0e-2}{mol.L^{-1}}$. L'équivalence est atteinte pour $V_E = \SI{8,0}{mL}$.
\begin{enumerate}
\item Calculer la concentration du diiode dans le prélèvement à $t = \SI{21}{min}$.
\item En déduire l'avancement volumique $\dfrac{x}{V}$ du mélange réactionnel à cette date.
\item La concentration maximale théorique de \ce{I2} est \SI{5,0e-2}{mol.L^{-1}}. La réaction est-elle terminée à $t = \SI{21}{min}$ ?
\end{enumerate}
\tcblower
\textbf{1.} À l'équivalence du dosage : $C \cdot V_E = 2\,[\ce{I2}] \cdot V_p$, donc
\[ [\ce{I2}] = \frac{C \cdot V_E}{2\,V_p} = \frac{\num{2,0e-2} \times \num{8,0}}{2 \times 10,0} = \SI{8,0e-3}{mol.L^{-1}}. \]
\textbf{2.} Le dosage détermine la quantité de matière de diiode $n(\ce{I2})$ présente dans le prélèvement trempé. Comme la trempe ne consomme pas le diiode, cette quantité est égale à celle prélevée dans le milieu réactionnel à l'instant $t$. La concentration dans le milieu est donc $[\ce{I2}] = n(\ce{I2}) / V_p$. D'après le tableau d'avancement, $[\ce{I2}] = \dfrac{x}{V}$, donc
\[ \frac{x}{V} = \SI{8,0e-3}{mol.L^{-1}}. \]
\textbf{3.} $\dfrac{x}{V} = \SI{8,0e-3}{mol.L^{-1}} < \SI{5,0e-2}{mol.L^{-1}} = \dfrac{x_{\max}}{V}$ : la réaction n'est \textbf{pas terminée} à $t = \SI{21}{min}$. On est même loin de $t_{1/2}$ (qui correspondrait à \SI{2,5e-2}{mol.L^{-1}}).
\end{exoresolu}

\courssub{Méthodes physiques : mesurer sans prélever}

Les méthodes physiques reposent sur une idée élégante : au lieu de prélever et doser, on mesure \emph{en continu} une grandeur physique du milieu réactionnel qui varie de façon connue avec la concentration d'une espèce.

\textbf{Colorimétrie et spectrophotométrie.}\par
Le diiode \ce{I2} en solution présente une couleur \textbf{brune} (jaune-orangé à brun foncé selon la concentration) : plus il se forme de diiode, plus la solution absorbe la lumière. Un \cle{spectrophotomètre} mesure l'\cle{absorbance} $A$ de la solution à une longueur d'onde choisie. D'après la \cle{loi de Beer-Lambert}, pour une solution diluée :
\[ A = \varepsilon \cdot \ell \cdot [\ce{I2}] \]
où $\varepsilon$ est le coefficient d'absorption molaire et $\ell$ la largeur de la cuve. L'absorbance est donc \textbf{proportionnelle} à la concentration du diiode : la courbe $A = f(t)$ donne directement l'allure de $[\ce{I2}] = f(t)$, sans aucun prélèvement.

\begin{popfigure}
\begin{tikzpicture}[scale=0.9, font=\small]
% Source
\draw[thick, popdark, fill=popgoldL] (0,0.6) rectangle (1.2,1.8);
\node[popdark] at (0.6,1.2) {source};
\node[popdark] at (0.6,0.85) {lumineuse};
% Faisceau avant cuve
\draw[-{Stealth}, very thick, popgold!80!black] (1.2,1.2) -- (2.6,1.2);
% Cuve
\draw[thick, popdark] (2.6,0.4) -- (2.6,2.0) (4.2,0.4) -- (4.2,2.0);
\draw[thick, popdark] (2.6,0.4) -- (4.2,0.4);
\fill[poporange!50] (2.65,0.45) rectangle (4.15,1.7);
\node[popdark] at (3.4,0.05) {cuve : solution brune de \ce{I2}};
% Faisceau après cuve (atténué)
\draw[-{Stealth}, thick, popgold!80!black, dashed] (4.2,1.2) -- (5.6,1.2);
\node[popgold!70!black] at (4.9,1.55) {lumière};
\node[popgold!70!black] at (4.9,1.3) {atténuée};
% Détecteur
\draw[thick, popdark, fill=popblueL] (5.6,0.6) rectangle (6.8,1.8);
\node[popdark] at (6.2,1.2) {détecteur};
% Affichage
\draw[thick, popdark, fill=popcream] (7.6,0.5) rectangle (9.6,1.9);
\node[popdark] at (8.6,1.55) {afficheur};
\node[popink] at (8.6,1.05) {$A = 0{,}42$};
\draw[-{Stealth}, popmuted] (6.8,1.2) -- (7.6,1.2);
\end{tikzpicture}
\legende{Principe du spectrophotomètre : la solution brune de diiode absorbe une partie de la lumière ; l'absorbance mesurée est proportionnelle à $[\ce{I2}]$.}
\end{popfigure}

\textbf{Conductimétrie.}\par
Si la réaction modifie la concentration totale des ions en solution, la \cle{conductivité} $\sigma$ du milieu (mesurée avec un conductimètre) varie au cours du temps. Comme $\sigma$ est une fonction connue des concentrations ioniques, le suivi de $\sigma(t)$ permet de remonter à l'avancement $x(t)$. Cette méthode convient par exemple à l'hydrolyse basique d'un ester (saponification), où les ions \ce{HO-} très conducteurs sont remplacés par des ions carboxylate.

\textbf{Mesure de pression.}\par
Si la réaction produit (ou consomme) un \textbf{gaz} en milieu fermé, la pression $P$ du récipient varie au cours du temps. En assimilant le gaz à un gaz parfait, $n_{\text{gaz}} = \dfrac{PV}{RT}$ : le suivi de $P(t)$ donne la quantité de gaz formé, donc l'avancement. Exemple : la décomposition de l'eau oxygénée \ce{2 H2O2 -> 2 H2O + O2}, suivie par la pression du dioxygène dégagé.

\begin{propriete}[Avantage décisif des méthodes physiques]
Les méthodes physiques permettent une mesure \textbf{continue} et \textbf{rapide}, \textbf{sans prélèvement} et \textbf{sans destruction} de l'échantillon : le milieu réactionnel reste intact, on obtient un grand nombre de points, et la courbe cinétique est plus précise. Elles ne nécessitent ni trempe ni dosage.
\end{propriete}

\begin{savaistu}[Suivre une réaction... à l'œil nu]
La réaction entre les ions iodure et peroxodisulfate est si commode que la couleur brune du diiode permet de la suivre \emph{sans aucun instrument} : on compare simplement la teinte du mélange à des étalons de couleur, ou on la mesure au colorimètre du laboratoire. Mieux : en ajoutant au mélange une petite quantité connue de thiosulfate et un peu d'empois d'amidon, la solution devient \textbf{brutalement bleu-noir} à une date précise (quand tout le thiosulfate est consommé) : c'est le principe des \og réactions horloge \fg{}, spectaculaires en travaux pratiques !
\end{savaistu}

\courssub{Choisir la méthode adaptée au système étudié}

Le choix de la méthode de suivi dépend des propriétés des espèces mises en jeu :

\begin{center}
\begin{tabularx}{\linewidth}{|l|X|X|}
\hline
\rowcolor{popblueL}
\textbf{Situation} & \textbf{Méthode conseillée} & \textbf{Exemple} \\
\hline
Espèce colorée formée ou consommée & Spectrophotométrie (absorbance) & Formation de \ce{I2} (brun) \\
\hline
Concentration en ions variable & Conductimétrie & Saponification d'un ester \\
\hline
Dégagement (ou consommation) d'un gaz & Mesure de pression ou de volume de gaz & Décomposition de \ce{H2O2} \\
\hline
Aucune propriété physique exploitable & Prélèvements + trempe + dosage & Estérification (dosage de l'acide restant par \ce{NaOH}) \\
\hline
\end{tabularx}
\end{center}

\courssub{Traitement des données : du tableau à la courbe}

Quelle que soit la méthode, l'exploitation suit toujours la même démarche.

\begin{methode}[Exploiter un suivi temporel]
\begin{enumerate}
\item Rassembler les mesures dans un tableau à deux lignes : dates $t$ et concentration $[X]$ (ou absorbance, conductivité, pression...).
\item Si nécessaire, convertir la grandeur mesurée en concentration (loi de Beer-Lambert, relation de dosage, gaz parfait...).
\item Tracer la courbe $[X] = f(t)$ : dates en abscisse, concentration en ordonnée, points reliés par une courbe lisse.
\item Exploiter la courbe : vitesse moyenne entre deux dates (pente du segment), vitesse instantanée à une date (pente de la tangente), temps de demi-réaction $t_{1/2}$ (abscisse du point d'ordonnée $\dfrac{[X]_{\text{finale}}}{2}$ pour un produit).
\end{enumerate}
\end{methode}

\begin{exemplebox}[Tableau de suivi de la formation du diiode]
Pour le mélange de l'exercice résolu précédent, les dosages successifs donnent :
\begin{center}
\begin{tabularx}{\linewidth}{|l|X|X|X|X|X|X|}
\hline
\rowcolor{popblueL}
$t$ (min) & 0 & 7 & 14 & 21 & 30 & 60 \\
\hline
$[\ce{I2}]$ (\SI{e-3}{mol.L^{-1}}) & 0 & 2{,}7 & 5{,}3 & 8{,}0 & 11{,}0 & 19{,}0 \\
\hline
\end{tabularx}
\end{center}
Le tracé de ces points donne une courbe croissante, de pente décroissante : la vitesse de formation du diiode \emph{diminue} au cours du temps, car les concentrations des réactifs diminuent. La vitesse moyenne entre $t = 0$ et $t = \SI{21}{min}$ vaut :
\[ v_{\text{moy}} = \frac{[\ce{I2}]_{21} - [\ce{I2}]_{0}}{21 - 0} = \frac{\num{8,0e-3}}{21} \approx \SI{3,8e-4}{mol.L^{-1}.min^{-1}}. \]
\end{exemplebox}

\begin{attention}[Unités et conversions]
Avant tout calcul de vitesse, vérifie la cohérence des unités : concentration en \si{mol.L^{-1}} et temps en minutes donnent une vitesse en \si{mol.L^{-1}.min^{-1}} ; en secondes, en \si{mol.L^{-1}.s^{-1}}. Autre piège : la trempe dilue le prélèvement, mais la concentration calculée par le dosage est celle du prélèvement \emph{avant} dilution si l'on raisonne en quantités de matière — raisonne toujours sur $n = C \cdot V$ pour éviter toute confusion.
\end{attention}

\begin{exercice}[Choix d'une méthode de suivi]
Pour chacune des réactions suivantes, indique la méthode de suivi la plus adaptée et justifie :
\begin{enumerate}
\item l'estérification de l'éthanol par l'acide éthanoïque ;
\item la décomposition de l'eau oxygénée en présence de ions \ce{Fe^3+} ;
\item la réaction entre les ions iodure et les ions peroxodisulfate.
\end{enumerate}
\end{exercice}

\begin{corrige}[Exercice 1]
\begin{enumerate}
\item Estérification : aucune espèce colorée, pas de gaz, variation ionique négligeable $\Rightarrow$ \textbf{prélèvements + trempe + dosage} de l'acide éthanoïque restant par la soude.
\item Décomposition de \ce{H2O2} : dégagement de dioxygène $\Rightarrow$ \textbf{mesure de pression} (ou de volume de gaz) en récipient fermé.
\item \ce{I-} + \ce{S2O8^2-} : formation de diiode brun $\Rightarrow$ \textbf{spectrophotométrie} (suivi continu de l'absorbance) ou, classiquement, prélèvements + trempe + dosage par le thiosulfate.
\end{enumerate}
\end{corrige}

\begin{aretenir}[L'essentiel de la section]
\begin{itemize}
\item Suivre une réaction, c'est mesurer la concentration d'une espèce à différentes dates pour tracer la courbe cinétique.
\item Méthode chimique : prélèvements à dates fixes, \textbf{trempe} (refroidissement brutal + dilution) pour bloquer la réaction, puis dosage (ex. : \ce{I2} dosé par \ce{S2O3^2-}).
\item Méthodes physiques : spectrophotométrie (espèce colorée), conductimétrie (ions), pression (gaz) : mesures \textbf{continues et non destructives}.
\item Exploitation : tableau $t$ / concentration $\rightarrow$ courbe $\rightarrow$ vitesses (moyenne, instantanée) et temps de demi-réaction.
\end{itemize}
\end{aretenir}

\coursec{Facteurs cinétiques et catalyse}

Dans la section précédente, nous avons appris à \emph{suivre} l'évolution temporelle d'un système chimique : tremper les prélèvements, doser le diiode formé, tracer la courbe $[\ce{I2}] = f(t)$ et en déduire les vitesses. Une question surgit alors naturellement : \textbf{peut-on agir sur cette vitesse ?} Peut-on rendre une réaction lente plus rapide — ou ralentir une réaction indésirable ? C'est toute la problématique de cette section, et elle a des enjeux immenses : conserver les aliments au frais, accélérer la production d'engrais à la SONARA, empêcher la corrosion des ponts sur le Wouri, ou encore équiper les motos de Douala de pots catalytiques. Les grandeurs sur lesquelles on peut agir s'appellent les \cle{facteurs cinétiques}.

\courssub{Notion de facteur cinétique}

\begin{definition}[Facteur cinétique]
Un \cle{facteur cinétique} est toute grandeur physique ou chimique dont la modification entraîne une modification de la \textbf{vitesse} d'une réaction chimique.
\end{definition}

L'expérience et l'observation quotidienne en fournissent quatre, que nous allons étudier en détail :

\begin{itemize}
\item la \cle{concentration} des réactifs ;
\item la \cle{température} du milieu réactionnel ;
\item la présence d'un \cle{catalyseur} ;
\item la \cle{surface de contact} entre les réactifs (lorsque l'un d'eux est solide).
\end{itemize}

\begin{attention}[Vitesse et état final : deux notions distinctes]
Modifier un facteur cinétique change la \textbf{vitesse} à laquelle le système évolue, mais \textbf{pas nécessairement} l'état final. Certains facteurs (concentration, température) peuvent aussi déplacer un équilibre ; en revanche, un catalyseur ne modifie \emph{jamais} l'état final : il ne fait que raccourcir le temps pour l'atteindre. Ne confonds jamais « réaction plus rapide » et « réaction plus complète ».
\end{attention}

\courssub{Influence de la concentration des réactifs}

\textbf{Intuition.} Imagine un marché bondé le samedi à Mokolo : plus il y a de monde, plus les rencontres (et les bousculades !) sont fréquentes. De même, une réaction chimique résulte de \emph{chocs efficaces} entre les entités réactives. Plus les réactifs sont concentrés, plus les chocs par unité de temps sont nombreux, et plus la réaction avance vite.

\begin{propriete}[Influence de la concentration]
La vitesse d'une réaction \textbf{augmente} lorsque la \textbf{concentration} d'un ou de plusieurs réactifs \textbf{augmente}.
\end{propriete}

\textbf{Interprétation graphique.} Sur la courbe $[\ce{I2}] = f(t)$ de la réaction entre les ions iodure et les ions peroxodisulfate :
\[ \ce{S2O8^{2-} + 2 I- -> I2 + 2 SO4^{2-}} \quad \text{(réaction lente, totale, à température ambiante)} \]
la vitesse instantanée est la pente de la tangente. Si l'on refait l'expérience avec une concentration initiale en ions \ce{S2O8^{2-}} deux fois plus grande, la tangente à l'origine est \textbf{plus pentue} : la vitesse initiale est plus grande. Remarque aussi que sur une même courbe, la pente \emph{diminue} au cours du temps : les réactifs se consomment, leur concentration décroît, donc la vitesse décroît. C'est une preuve graphique directe du rôle de la concentration.

\begin{exemplebox}[Dilution et vitesse initiale]
On réalise deux mélanges identiques en volume, mais le second contient deux fois moins d'ions iodure. On mesure la vitesse initiale de formation du diiode :
\begin{center}
\begin{tabularx}{\linewidth}{|l|X|X|}
\hline
\rowcolor{popblueL} Expérience & $[\ce{I-}]_0$ (\si{\mol\per\liter}) & $v_0$ (\si{\milli\mol\per\liter\per\minute}) \\
\hline
1 & 0,20 & 1,6 \\
\hline
2 & 0,10 & 0,8 \\
\hline
\end{tabularx}
\end{center}
En divisant la concentration par 2, la vitesse initiale a été divisée par 2 : la concentration est bien un facteur cinétique.
\end{exemplebox}

\courssub{Influence de la température}

\textbf{Intuition.} La température traduit l'agitation des entités. En chauffant, on donne aux ions et molécules plus d'énergie : les chocs deviennent à la fois plus fréquents et plus « violents », donc plus souvent efficaces.

\begin{propriete}[Influence de la température]
La vitesse d'une réaction \textbf{augmente} lorsque la \textbf{température} \textbf{augmente}. Règle qualitative usuelle : une élévation de température de l'ordre de \SI{10}{\celsius} multiplie souvent la vitesse par un facteur 2 à 3 (ordre de grandeur, pas une loi exacte).
\end{propriete}

\begin{exemplebox}[Chiffres parlants]
Pour la réaction \ce{S2O8^{2-}}/\ce{I-}, le temps de demi-réaction vaut environ \SI{8}{min} à \SI{20}{\celsius} ; il tombe à environ \SI{2}{min} vers \SI{40}{\celsius}. La réaction est environ 4 fois plus rapide pour \SI{20}{\celsius} de plus.
\end{exemplebox}

\textbf{Applications quotidiennes.} On exploite ce facteur dans les deux sens :
\begin{itemize}
\item \emph{Ralentir} : la conservation des aliments au réfrigérateur (\SI{4}{\celsius}) ou au congélateur (\SI{-18}{\celsius}) ralentit les réactions de décomposition et de fermentation. C'est pourquoi le jus de bissap se conserve plusieurs jours au frais mais tourne en une journée à la chaleur de Maroua.
\item \emph{Accélérer} : la cuisson des aliments, la stérilisation, les synthèses industrielles à chaud.
\item La \cle{trempe} vue dans la section précédente (dilution + refroidissement brutal du prélèvement) repose précisément sur ces deux facteurs : on « fige » la réaction pour doser tranquillement.
\end{itemize}

\courssub{Le catalyseur : définition et propriétés}

\textbf{Intuition.} Certaines réactions sont possibles mais désespérément lentes. Un catalyseur agit comme un « entremetteur » : il participe au mécanisme, ouvre un chemin plus facile, puis est restitué intact à la fin.

\begin{definition}[Catalyseur]
Un \cle{catalyseur} est une espèce chimique qui \textbf{augmente la vitesse} d'une réaction chimique \textbf{sans être consommée} : il est \textbf{régénéré} en fin de réaction et n'apparaît pas dans l'équation bilan. Il ne modifie \textbf{ni l'état final, ni la quantité de produit} obtenue : il permet seulement d'atteindre cet état \textbf{plus vite}.
\end{definition}

\begin{propriete}[Propriétés d'un catalyseur]
\begin{enumerate}
\item Il est utilisé en \textbf{faible quantité} (quelques gouttes, quelques milligrammes suffisent).
\item Il est \textbf{régénéré} : on le retrouve intact en fin de réaction.
\item Il \textbf{ne modifie pas l'état final} : même asymptote sur la courbe, même quantité de produit.
\item Il est \textbf{sélectif} : un catalyseur donné n'accélère qu'une réaction (ou une famille de réactions) bien précise.
\end{enumerate}
\end{propriete}

\begin{attention}[Erreur fréquente à éviter absolument]
Un catalyseur \textbf{ne rend pas possible une réaction impossible}. Il accélère une réaction \emph{déjà possible} mais lente. Dire « sans catalyseur la réaction n'a pas lieu » est faux : sans catalyseur, elle a lieu... très lentement. Autre piège : ne pas écrire le catalyseur comme réactif dans l'équation bilan ; on l'indique au-dessus de la flèche.
\end{attention}

\courssub{La catalyse homogène : les ions \ce{Fe^{3+}} sur la réaction \ce{I-}/\ce{S2O8^{2-}}}

\begin{definition}[Catalyse homogène]
La catalyse est \cle{homogène} lorsque le catalyseur et les réactifs se trouvent dans la \textbf{même phase} (généralement en solution aqueuse).
\end{definition}

\begin{experience}[Mise en évidence du rôle des ions \ce{Fe^{3+}}]
\textbf{Dispositif et protocole.} On prépare deux erlenmeyers \emph{identiques} contenant chacun \SI{20}{mL} de solution d'iodure de potassium (\ce{K+} + \ce{I-}) à \SI{0,20}{\mol\per\liter} et \SI{20}{mL} de solution de peroxodisulfate de potassium (\ce{2 K+} + \ce{S2O8^{2-}}) à \SI{0,10}{\mol\per\liter}, à la même température. Dans le seul erlenmeyer B, on ajoute \textbf{quelques gouttes} de solution de chlorure de fer(III) (\ce{Fe^{3+}} + \ce{3 Cl-}). On déclenche le chronomètre au mélange.

\textbf{Observations.}
\begin{itemize}
\item Erlenmeyer A (témoin) : la teinte brune du diiode apparaît progressivement, en plusieurs minutes.
\item Erlenmeyer B (avec \ce{Fe^{3+}}) : la teinte brune apparaît presque immédiatement et s'intensifie en quelques secondes.
\item Au bout d'un temps long, les \textbf{deux} erlenmeyers présentent la \textbf{même teinte finale}.
\end{itemize}

\textbf{Interprétation.} La réaction est :
\[ \ce{S2O8^{2-} + 2 I- -> I2 + 2 SO4^{2-}} \quad \text{(lente ; accélérée par \ce{Fe^{3+}})} \]
Les ions \ce{Fe^{3+}} participent à un mécanisme en deux étapes rapides :
\[ \ce{2 Fe^{3+} + 2 I- -> 2 Fe^{2+} + I2} \qquad \text{puis} \qquad \ce{2 Fe^{2+} + S2O8^{2-} -> 2 Fe^{3+} + 2 SO4^{2-}} \]
Les ions \ce{Fe^{3+}} consommés dans la première étape sont \textbf{régénérés} dans la seconde : la somme des deux étapes redonne exactement le bilan global, sans \ce{Fe^{3+}}. La teinte finale identique prouve que l'état final est inchangé.
\end{experience}

\begin{methode}[Montrer expérimentalement le rôle catalytique de \ce{Fe^{3+}}]
\begin{enumerate}
\item Réaliser \textbf{deux milieux réactionnels rigoureusement identiques} (mêmes volumes, mêmes concentrations, même température) : c'est le principe du \textbf{témoin}.
\item Ajouter quelques gouttes de solution de \ce{Fe^{3+}} dans l'un seulement.
\item Comparer les \textbf{durées d'apparition de la teinte brune} du diiode (ou mesurer les temps de demi-réaction).
\item Vérifier que la \textbf{teinte finale est la même} dans les deux béchers : le catalyseur n'a pas changé l'état final.
\item Conclure : \ce{Fe^{3+}} est un catalyseur (accélère sans être consommé).
\end{enumerate}
\end{methode}

\begin{exemplebox}[Exemple chiffré]
Dans les conditions de l'expérience précédente à \SI{25}{\celsius}, le temps de demi-réaction vaut environ $t_{1/2} \approx \SI{8}{min}$ sans catalyseur, et seulement $t'_{1/2} \approx \SI{20}{s}$ avec quelques gouttes de \ce{Fe^{3+}}. La vitesse a été multipliée par un facteur de l'ordre de 24, alors que la concentration en \ce{Fe^{3+}} est cent fois plus faible que celle des réactifs.
\end{exemplebox}

\begin{popfigure}
\begin{center}
\begin{tikzpicture}
\begin{axis}[
width=0.92\linewidth, height=6.2cm,
xlabel={$t$ (\si{\minute})}, ylabel={$[\ce{I2}]$ (\si{\milli\mol\per\liter})},
xmin=0, xmax=20, ymin=0, ymax=11,
grid=major, grid style={popline!50},
legend style={at={(0.97,0.35)}, anchor=east, font=\small, draw=popline},
axis line style={popdark},
]
\addplot[popcurve, popblue, domain=0:20, samples=200] {10*(1-exp(-0.0866*x))};
\addlegendentry{Sans catalyseur}
\addplot[popcurve, poporange, domain=0:20, samples=200] {10*(1-exp(-2.1*x))};
\addlegendentry{Avec \ce{Fe^{3+}}}
\addplot[popmuted, dashed, domain=0:20] {10};
\draw[popblue, dashed] (axis cs:0,0) -- (axis cs:8,5);
\draw[poporange, dashed] (axis cs:0,0) -- (axis cs:0.33,9.5);
\node[popmuted, font=\small, anchor=west] at (axis cs:12,9.2) {même asymptote};
\end{axis}
\end{tikzpicture}
\end{center}
\legende{Concentration en diiode formé en fonction du temps, sans catalyseur (bleu) et en présence d'ions \ce{Fe^{3+}} (orange). La tangente initiale (pointillés) est beaucoup plus pentue avec catalyseur, mais les deux courbes tendent vers la \textbf{même asymptote} : l'état final est inchangé.}
\end{popfigure}

\courssub{La catalyse hétérogène : le platine et la synthèse de l'eau}

\begin{definition}[Catalyse hétérogène]
La catalyse est \cle{hétérogène} lorsque le catalyseur et les réactifs se trouvent dans des \textbf{phases différentes} : le plus souvent, un catalyseur \textbf{solide} et des réactifs \textbf{gazeux} ou \textbf{liquides}. La réaction se produit à la \emph{surface} du solide.
\end{definition}

\begin{experience}[Le platine catalyse la synthèse de l'eau]
Le mélange dihydrogène--dioxygène est \emph{cinétiquement bloqué} : à température ambiante, la réaction
\[ \ce{2 H2 (g) + O2 (g) -> 2 H2O (l)} \quad \text{(très lente sans catalyseur ; explosive avec Pt)} \]
est pratiquement inerte : un mélange \ce{H2}/\ce{O2} peut rester des années sans réagir. Approchons un fil ou une mousse de \textbf{platine} (solide) : la réaction s'emballe immédiatement, de façon spectaculaire, et le platine est retrouvé intact après coup. Les molécules \ce{H2} et \ce{O2} se fixent à la surface du platine, y sont « activées » puis réagissent ; les molécules d'eau se désorbent, libérant la surface pour de nouvelles molécules.
\end{experience}

\begin{savaistu}[Le pot catalytique : de la chimie au quotidien]
Les pots catalytiques des voitures et motos contiennent des métaux précieux (\textbf{platine}, palladium, rhodium) déposés sur une céramique en nid d'abeille. Ils accélèrent la transformation des gaz toxiques d'échappement — monoxyde de carbone \ce{CO}, oxydes d'azote, hydrocarbures imbrûlés — en gaz moins nocifs (\ce{CO2}, \ce{N2}, \ce{H2O}). C'est un exemple remarquable de \textbf{catalyse hétérogène} : catalyseur solide, réactifs gazeux. Le platine est aussi utilisé industriellement dans le raffinage (reformage), par exemple à la SONARA de Limbé.
\end{savaistu}

\courssub{La catalyse enzymatique}

\begin{definition}[Catalyse enzymatique]
Les \cle{enzymes} sont des catalyseurs \textbf{biologiques} (protéines) d'une \textbf{efficacité} et d'une \textbf{spécificité} extraordinaires : chaque enzyme n'accélère qu'une seule réaction (ou un type de réaction) bien précis, dans des conditions douces (température proche de \SI{37}{\celsius}, pH contrôlé).
\end{definition}

\begin{exemplebox}[Des enzymes partout autour de nous]
\begin{itemize}
\item \textbf{Fermentation} : les levures produisent des enzymes qui transforment les sucres en éthanol — c'est le principe de la fabrication du vin de palme et du bili-bili.
\item \textbf{Digestion} : l'amylase de la salive découpe l'amidon ; la pepsine de l'estomac découpe les protéines.
\item \textbf{Lessives « enzymatiques »} : elles contiennent des enzymes qui dégradent les taches de graisse ou de protéines à basse température.
\end{itemize}
\end{exemplebox}

\begin{attention}[Les enzymes sont sensibles]
Une enzyme perd son activité si la température est trop élevée (on dit qu'elle est \emph{dénaturée}) ou si le pH s'écarte de sa zone optimale. C'est pourquoi on ne fait pas bouillir un levain que l'on veut garder actif.
\end{attention}

\begin{popfigure}
\begin{center}
\begin{tikzpicture}[font=\small]
% Vignette 1 : homogène
\begin{scope}
\draw[fill=popblueL, rounded corners] (-1.6,-1.3) rectangle (1.6,1.3);
\fill[popblue] (-0.7,0.1) circle (0.12); \fill[popblue] (0.5,-0.4) circle (0.12);
\fill[poporange] (0.1,0.5) circle (0.12); \fill[poporange] (-0.9,-0.6) circle (0.12);
\fill[popgreen] (0.9,0.6) circle (0.09); \fill[popgreen] (-0.2,-0.9) circle (0.09);
\node[align=center] at (0,-1.75) {\textbf{Catalyse homogène}\\ une seule phase\\ \ce{Fe^{3+}} + \ce{I-}/\ce{S2O8^{2-}}};
\end{scope}
% Vignette 2 : hétérogène
\begin{scope}[xshift=5.2cm]
\draw[fill=poporangeL, rounded corners] (-1.6,-1.3) rectangle (1.6,1.3);
\fill[popmuted!70] (-1.6,-1.3) rectangle (1.6,-0.55);
\node[popdark, font=\footnotesize] at (0,-0.95) {platine solide};
\fill[popblue] (-0.7,0.3) circle (0.12); \fill[popblue] (0.6,0.7) circle (0.12);
\fill[poporange] (0.0,0.1) circle (0.12); \fill[poporange] (1.0,0.2) circle (0.12);
\draw[->, popdark, thick] (-0.5,0.2) -- (-0.5,-0.45);
\draw[->, popdark, thick] (0.4,0.5) -- (0.4,-0.45);
\node[align=center] at (0,-1.75) {\textbf{Catalyse hétérogène}\\ deux phases\\ Pt (s) + \ce{H2}/\ce{O2} (g)};
\end{scope}
% Vignette 3 : enzymatique
\begin{scope}[xshift=10.4cm]
\draw[fill=popgreenL, rounded corners] (-1.6,-1.3) rectangle (1.6,1.3);
\draw[popgreen, very thick, rounded corners=2pt] (-0.9,-0.2) .. controls (-0.4,0.8) and (0.2,0.8) .. (0.7,0.1) .. controls (0.9,-0.3) and (0.5,-0.7) .. (0.0,-0.5) .. controls (-0.5,-0.3) .. (-0.9,-0.2);
\fill[poporange] (-0.1,0.15) circle (0.12);
\draw[->, popdark, thick] (0.3,0.9) -- (0.05,0.35);
\node[align=center] at (0,-1.75) {\textbf{Catalyse enzymatique}\\ enzyme spécifique\\ fermentation, digestion};
\end{scope}
\end{tikzpicture}
\end{center}
\legende{Les trois types de catalyse. Homogène : catalyseur et réactifs dans la même phase. Hétérogène : la réaction a lieu à la surface du catalyseur solide. Enzymatique : le substrat (orange) se fixe sur le site actif de l'enzyme (vert), comme une clé dans une serrure.}
\end{popfigure}

\courssub{Influence de la surface de contact}

\textbf{Intuition.} Lorsqu'un réactif est solide, seules les entités situées à sa \emph{surface} peuvent entrer en contact avec les autres réactifs. En divisant le solide, on multiplie cette surface accessible.

\begin{propriete}[Surface de contact]
Plus la \textbf{surface de contact} entre les réactifs est grande, plus la réaction est \textbf{rapide}. Un solide \textbf{divisé} (poudre, grenaille, copeaux) réagit beaucoup plus vite qu'un solide \textbf{en un bloc} de même masse.
\end{propriete}

\begin{exemplebox}[Zinc en poudre contre zinc en grenaille]
On verse de l'acide chlorhydrique en excès sur deux échantillons de \SI{1}{g} de zinc : l'un en grenaille (gros grains), l'autre en poudre fine. La réaction
\[ \ce{Zn (s) + 2 H3O+ -> Zn^{2+} + H2 (g) + 2 H2O} \]
dégage du dihydrogène. Avec la poudre, le dégagement est vif et s'achève en moins d'une minute ; avec la grenaille, il dure plusieurs minutes. Même masse, même produit final : seule la vitesse change. C'est aussi pourquoi la poussière de farine ou de charbon peut provoquer des explosions dans les silos : une surface énorme, une réaction quasi instantanée.
\end{exemplebox}

\courssub{Exercice résolu de synthèse}

\begin{exoresolu}[Catalyseur, concentration et température]
On étudie la réaction \ce{S2O8^{2-} + 2 I- -> I2 + 2 SO4^{2-}}. Quatre expériences sont réalisées avec les mêmes quantités de réactifs ; on mesure le temps de demi-réaction $t_{1/2}$ :
\begin{center}
\begin{tabularx}{\linewidth}{|l|X|X|X|}
\hline
\rowcolor{popblueL} Exp. & Température & Catalyseur & $t_{1/2}$ \\
\hline
1 & \SI{20}{\celsius} & aucun & \SI{8}{min} \\
\hline
2 & \SI{20}{\celsius} & \ce{Fe^{3+}} & \SI{20}{s} \\
\hline
3 & \SI{40}{\celsius} & aucun & \SI{2}{min} \\
\hline
4 & \SI{20}{\celsius} & \ce{Fe^{3+}} en grande quantité & \SI{20}{s} \\
\hline
\end{tabularx}
\end{center}
1. Comparer les expériences 1 et 2 ; conclure. \quad 2. Comparer 1 et 3 ; conclure. \quad 3. La quantité de diiode obtenue au bout d'un temps très long diffère-t-elle entre les quatre expériences ? \quad 4. Que montre la comparaison des expériences 2 et 4 ?
\tcblower
\textbf{Solution détaillée.}
\begin{enumerate}
\item À température identique, $t_{1/2}$ passe de \SI{8}{min} à \SI{20}{s} en présence de \ce{Fe^{3+}} : la réaction est environ 24 fois plus rapide. Les ions \ce{Fe^{3+}} jouent le rôle de \textbf{catalyseur} (catalyse homogène, car tout est en solution).
\item Sans catalyseur, passer de \SI{20}{\celsius} à \SI{40}{\celsius} divise $t_{1/2}$ par 4 : la \textbf{température} est un facteur cinétique ; la vitesse augmente avec la température.
\item \textbf{Non.} Les quantités initiales de réactifs sont identiques et la réaction est totale : l'état final est le même dans les quatre cas. Température et catalyseur modifient la \emph{vitesse}, pas la quantité finale de diiode (l'asymptote des courbes $[\ce{I2}] = f(t)$ est la même).
\item Multiplier la quantité de catalyseur ne change pas ici le résultat de façon notable : quelques gouttes de \ce{Fe^{3+}} suffisent, car le catalyseur est \textbf{régénéré} en permanence. Il agit en \textbf{faible quantité}.
\end{enumerate}
\end{exoresolu}

\begin{aretenir}[L'essentiel de la section]
\begin{itemize}
\item Un \textbf{facteur cinétique} est une grandeur qui modifie la vitesse d'une réaction : concentration des réactifs, température, catalyseur, surface de contact.
\item La vitesse \textbf{augmente} avec la concentration et avec la température ; elle augmente aussi quand le solide est plus divisé.
\item Un \textbf{catalyseur} accélère une réaction \emph{possible mais lente}, sans être consommé (il est régénéré) et \textbf{sans modifier l'état final}. Il est sélectif et efficace en faible quantité.
\item Trois types de catalyse : \textbf{homogène} (même phase : \ce{Fe^{3+}} sur \ce{I-}/\ce{S2O8^{2-}}), \textbf{hétérogène} (phases différentes : Pt sur \ce{H2}/\ce{O2}), \textbf{enzymatique} (enzymes biologiques, très efficaces et spécifiques).
\item Piège majeur : un catalyseur ne rend jamais possible une réaction impossible.
\end{itemize}
\end{aretenir}

Nous disposons désormais de tous les outils pour aborder la dernière partie de la leçon : l'\textbf{étude cinétique de l'estérification} entre l'acide éthanoïque et l'éthanol, où nous verrons comment la température et l'acide sulfurique permettent d'atteindre plus vite — mais sans la déplacer — la fameuse asymptote de l'équilibre.

\coursec{Étude cinétique de l'estérification}

Nous venons de voir que la température, la concentration et la présence d'un catalyseur modifient la \cle{vitesse} d'une réaction chimique. Mais une question reste en suspens : ces facteurs modifient-ils aussi la \emph{quantité finale} de produit obtenu ? Pour y répondre, rien de tel qu'une réaction lente, limitée, et facile à suivre au laboratoire : l'\cle{estérification} de l'acide éthanoïque par l'éthanol. C'est la réaction privilégiée par les concepteurs du Baccalauréat pour cette étude — et tu vas comprendre pourquoi.

\courssub{Rappel : une réaction lente et limitée}

L'acide éthanoïque (l'acide du vinaigre) réagit avec l'éthanol pour former l'\cle{éthanoate d'éthyle}, un ester à l'odeur agréable de fruits mûrs, et de l'eau :

\begin{center}
\chemfig{CH_3-C(=[1]O)-[7]OH} \;+\; \chemfig{CH_3-CH_2-OH} \; $\rightleftharpoons$ \; \chemfig{CH_3-C(=[1]O)-[7]O-CH_2-CH_3} \;+\; \chemfig{H_2O}
\end{center}

\[ \ce{CH3COOH + C2H5OH <=> CH3COOC2H5 + H2O} \]

Cette équation se lit : acide éthanoïque + éthanol $\rightleftharpoons$ éthanoate d'éthyle + eau. La double flèche $\rightleftharpoons$ est essentielle : elle signifie que la réaction inverse (l'\emph{hydrolyse} de l'ester) a lieu en même temps que la réaction directe.

\begin{aretenir}[Trois caractéristiques de l'estérification]
\begin{itemize}
\item \textbf{Lente} : à température ambiante, sans catalyseur, il faut plusieurs jours pour approcher l'état final ;
\item \textbf{Athermique} (ou presque) : elle ne dégage pratiquement pas de chaleur ($\Delta_r H \approx 0$) ;
\item \textbf{Limitée} : elle conduit à un \cle{équilibre chimique} — la réaction semble s'arrêter alors qu'il reste des réactifs.
\end{itemize}
\end{aretenir}

\courssub{Protocole de suivi cinétique}

Comment suivre, heure par heure, l'avancement d'une réaction qui dure plusieurs jours ? On utilise la méthode des \cle{prélèvements} avec \cle{trempe}, déjà rencontrée dans la section sur le suivi temporel. On réalise classiquement un suivi par \emph{méthode des tubes témoins}.

\begin{experience}[Suivi de l'estérification par dosage de l'acide restant]
\textbf{Dispositif et protocole :}
\begin{enumerate}
\item On prépare un mélange \cle{équimolaire} : $n_0 = \SI{0,10}{mol}$ d'acide éthanoïque et $n_0 = \SI{0,10}{mol}$ d'éthanol, plus quelques gouttes d'acide sulfurique concentré. Le tout est maintenu à température constante (bain-marie thermostaté).
\item On répartit le mélange dans une série de tubes à essai identiques, bouchés (tubes témoins).
\item À des dates $t$ fixées (par exemple toutes les 10 minutes au début, puis toutes les heures), on plonge un tube dans l'eau glacée : c'est la \textbf{trempe}, qui bloque (quasi arrête) la réaction en abaissant brutalement la température.
\item On dose alors l'\emph{acide éthanoïque restant} du tube par une solution de soude de concentration connue, en présence de phénolphtaléine :
\[ \ce{CH3COOH + HO- -> CH3COO- + H2O} \]
\end{enumerate}
\textbf{Exploitation :} si $n_{\text{acide}}(t)$ est la quantité d'acide restant à la date $t$, alors l'avancement est $x(t) = n_0 - n_{\text{acide}}(t)$, et la quantité d'ester formé vaut $n_{\text{ester}}(t) = x(t)$.
\end{experience}

\begin{attention}[Piège du dosage]
La soude dose \emph{tous} les acides présents, y compris les ions \ce{H3O+} du catalyseur sulfurique ! Comme la quantité de catalyseur est constante, on la détermine une fois pour toutes en dosant un tube à $t = 0$ (témoin), et on soustrait ce volume de soude de tous les résultats ultérieurs. Oublier cette correction est une erreur classique qui fausse systématiquement toutes les valeurs de $n_{\text{acide}}$.
\end{attention}

\courssub{Tracé et lecture des courbes}

À partir des dosages, on trace deux courbes complémentaires :

\begin{itemize}
\item $n_{\text{ester}} = f(t)$ : part de zéro, croît rapidement au début, puis de plus en plus lentement ;
\item $n_{\text{acide}} = g(t)$ : part de $n_0$, décroît de façon symétrique.
\end{itemize}

À tout instant, la conservation de la matière impose :
\[ n_{\text{acide}}(t) + n_{\text{ester}}(t) = n_0 \]

L'observation capitale est la suivante : la courbe $n_{\text{ester}} = f(t)$ ne monte pas indéfiniment — elle se rapproche d'une \cle{asymptote horizontale}.

\begin{definition}[Asymptote horizontale et caractère limité]
L'\cle{asymptote horizontale} de la courbe $n_{\text{ester}} = f(t)$ traduit le caractère \textbf{limité} de la réaction : au bout d'un temps suffisamment long, la quantité d'ester n'évolue plus car le système a atteint un \textbf{état d'équilibre chimique}. La réaction directe et la réaction inverse se poursuivent alors à la même vitesse (équilibre \emph{dynamique}), mais les quantités macroscopiques restent constantes. L'avancement final $x_f$ est inférieur à l'avancement maximal $x_{\max} = n_0$.
\end{definition}

\begin{exemplebox}[La fameuse limite des deux tiers]
Pour un mélange équimolaire acide éthanoïque + éthanol, l'expérience montre que la limite correspond à environ les \textbf{deux tiers} de la conversion totale :
\[ x_f \approx \frac{2}{3}\, n_0 \qquad \text{soit un rendement } \eta = \frac{x_f}{x_{\max}} \approx \SI{67}{\%} \]
Avec $n_0 = \SI{0,10}{mol}$ de chaque réactif, on obtient à l'équilibre environ $\SI{0,067}{mol}$ d'ester et d'eau, et il reste environ $\SI{0,033}{mol}$ d'acide et d'alcool. Ce résultat, indépendant de la température et du catalyseur, est une signature de l'équilibre de cette estérification (constante d'équilibre $K \approx 4$ à \SI{25}{\celsius}).
\end{exemplebox}

\courssub{Influence des facteurs cinétiques sur la courbe}

Reprenons maintenant la même expérience dans trois conditions différentes et comparons les courbes obtenues.

\begin{popfigure}
\begin{center}
\begin{tikzpicture}
\begin{axis}[
  width=0.92\linewidth, height=7cm,
  xlabel={$t$ (min)}, ylabel={$n_{\text{ester}}$ (mmol)},
  xmin=0, xmax=120, ymin=0, ymax=110,
  xtick={0,20,40,60,80,100,120},
  ytick={0,20,40,60,67,80,100},
  yticklabels={0,20,40,60,67,80,100},
  grid=both, grid style={popline!40},
  legend style={at={(0.97,0.35)}, anchor=east, font=\small, draw=popline},
  axis line style={popdark},
]
% Asymptote
\addplot[popmuted, dashed, thick, domain=0:120] {67};
% 20 °C : montée lente
\addplot[popcurve, popblue, domain=0:120, samples=200] {67*(1-exp(-0.022*x))};
\addlegendentry{\SI{20}{\celsius}, sans catalyseur}
% 60 °C : montée rapide
\addplot[popcurve, poporange, domain=0:120, samples=200] {67*(1-exp(-0.075*x))};
\addlegendentry{\SI{60}{\celsius}, sans catalyseur}
% H2SO4 à 20 °C : montée très rapide
\addplot[popcurve, popgreen, domain=0:120, samples=200] {67*(1-exp(-0.14*x))};
\addlegendentry{\SI{20}{\celsius} + \ce{H2SO4}}
% t1/2 repères
\addplot[poppink, dashed, thin] coordinates {(0,33.5) (31.5,33.5)};
\addplot[poppink, dashed, thin] coordinates {(31.5,0) (31.5,33.5)};
\node[font=\scriptsize, poppink, anchor=west] at (axis cs:33,26) {$t_{1/2}$ (\SI{20}{\celsius})};
\node[font=\scriptsize, popmuted, anchor=south west] at (axis cs:2,68) {asymptote : $x_f \approx \SI{67}{mmol}$};
\end{axis}
\end{tikzpicture}
\end{center}
\legende{Quantité d'ester formée en fonction du temps pour un mélange équimolaire ($n_0 = \SI{0,10}{mol}$). Les trois courbes convergent vers la \emph{même} asymptote : température et catalyseur changent la vitesse, pas la limite.}
\end{popfigure}

\textbf{Effet de la température.} En passant de \SI{20}{\celsius} à \SI{60}{\celsius}, la courbe monte beaucoup plus vite : la pente initiale (donc la vitesse initiale) est nettement plus grande, et l'asymptote est atteinte plus rapidement. Mais — point fondamental — \textbf{la valeur de l'asymptote ne change pas}. La température est un facteur purement cinétique ici : elle accélère à la fois l'estérification et l'hydrolyse, donc l'équilibre est atteint plus vite mais n'est pas déplacé de façon notable (la réaction étant quasi athermique, $\Delta_r H \approx 0$).

\textbf{Effet du catalyseur.} Quelques gouttes d'acide sulfurique concentré apportent des ions \ce{H3O+} qui jouent le rôle de \cle{catalyseur} (catalyse homogène). La courbe monte encore plus vite, mais rejoint exactement la même asymptote. Le catalyseur est retrouvé intact en fin de réaction.

\textbf{Effet des concentrations.} Si l'on introduit un \textbf{excès} d'un réactif (par exemple \SI{0,30}{mol} d'éthanol pour \SI{0,10}{mol} d'acide), deux choses changent : la réaction est plus rapide (facteur cinétique) \emph{et} la limite est déplacée — le rendement par rapport à l'acide peut dépasser \SI{90}{\%}. C'est le seul des trois facteurs qui modifie l'asymptote, car il déplace l'équilibre chimique (loi de Le Chatelier, que tu retrouveras dans l'étude des équilibres).

\begin{attention}[L'erreur qui coûte des points au Bac]
Beaucoup d'élèves écrivent : « le catalyseur augmente le rendement ». C'est \textbf{faux}. Le catalyseur (comme la température pour cette réaction) n'augmente que la \textbf{vitesse} : il permet d'atteindre plus vite l'équilibre, sans modifier la composition du mélange à l'équilibre. Seul un excès de réactif (ou l'élimination d'un produit) augmente le rendement. Retiens : \emph{vitesse} $\neq$ \emph{limite}.
\end{attention}

\courssub{Temps de demi-réaction et comparaison des situations}

\begin{methode}[Comparer deux conditions expérimentales]
Pour comparer rigoureusement deux courbes cinétiques :
\begin{enumerate}
\item Comparer les \textbf{pentes initiales} (tangentes à l'origine) : la plus pentue correspond à la vitesse initiale la plus grande ;
\item Déterminer les \textbf{temps de demi-réaction} $t_{1/2}$ : date à laquelle $n_{\text{ester}} = x_f/2$ ; le plus petit $t_{1/2}$ indique la réaction la plus rapide ;
\item Vérifier si les \textbf{asymptotes coïncident} : si oui, le facteur étudié est purement cinétique ; si non, l'équilibre lui-même a été déplacé.
\end{enumerate}
\end{methode}

\begin{exoresolu}[Temps de demi-réaction de l'estérification]
Sur la courbe à \SI{20}{\celsius} sans catalyseur, l'asymptote vaut $x_f = \SI{67}{mmol}$. Déterminer graphiquement $t_{1/2}$ dans les trois situations de la figure et classer les réactions par rapidité.
\tcblower
Le temps de demi-réaction correspond à $n_{\text{ester}} = x_f/2 = \SI{33,5}{mmol}$. On trace l'horizontale d'ordonnée \SI{33,5}{mmol} ; son intersection avec chaque courbe donne, par projection sur l'axe des temps :
\begin{center}
\begin{tabularx}{\linewidth}{|l|X|X|X|}
\hline
\rowcolor{popblueL} \textbf{Condition} & \SI{20}{\celsius} seul & \SI{60}{\celsius} & \SI{20}{\celsius} + \ce{H2SO4} \\
\hline
$t_{1/2}$ (lecture graphique) & $\approx \SI{32}{min}$ & $\approx \SI{9}{min}$ & $\approx \SI{5}{min}$ \\
\hline
\end{tabularx}
\end{center}
Classement par rapidité croissante : \SI{20}{\celsius} $<$ \SI{60}{\celsius} $<$ \SI{20}{\celsius} + \ce{H2SO4}. Dans les trois cas, l'asymptote reste $x_f = \SI{67}{mmol}$ : température et catalyseur accélèrent la réaction sans modifier sa limite.
\end{exoresolu}

\courssub{Bilan de la leçon}

\begin{aretenir}[Ce qu'il faut retenir de l'étude cinétique de l'estérification]
\begin{itemize}
\item L'estérification $\ce{CH3COOH + C2H5OH <=> CH3COOC2H5 + H2O}$ est \textbf{lente} et \textbf{limitée} : la courbe $n_{\text{ester}} = f(t)$ tend vers une asymptote horizontale ($\approx 2/3$ de conversion pour un mélange équimolaire, soit $\eta \approx \SI{67}{\%}$).
\item La \textbf{température} et le \textbf{catalyseur} (\ce{H3O+} apporté par \ce{H2SO4}) font atteindre l'asymptote plus rapidement \emph{sans changer sa valeur}.
\item Un \textbf{excès de réactif} accélère la réaction \emph{et} déplace la limite : c'est le seul facteur qui améliore le rendement.
\item Le suivi s'effectue par prélèvements (tubes témoins), trempe et dosage de l'acide restant par la soude (correction du catalyseur).
\end{itemize}
\end{aretenir}

\begin{savaistu}[L'acide sulfurique, catalyseur \emph{et} déshydratant]
Dans l'industrie — par exemple pour produire les arômes de fruits utilisés dans les boissons et les yaourts — on réalise les estérifications en présence d'acide sulfurique \emph{concentré}. Celui-ci joue un double rôle : ses ions \ce{H3O+} catalysent la réaction (effet cinétique), et son formidable pouvoir \textbf{déshydratant} absorbe l'eau formée. Or, retirer un produit déplace l'équilibre dans le sens direct : le rendement s'en trouve amélioré (effet thermodynamique). Un seul réactif, deux effets distincts. Les savonneries artisanales camerounaises exploitent d'ailleurs la réaction inverse — l'hydrolyse basique des esters de l'huile de palme — pour fabriquer le savon : c'est la saponification, elle, totale et irréversible !
\end{savaistu}

\newpage

\coursec{Travaux pratiques}

\coursec{1. Réactions rapides, lentes et très lentes}

\begin{definition}[Vitesse d'une réaction chimique]
La \cle{vitesse d'une réaction chimique} traduit la rapidité avec laquelle les réactifs se transforment en produits. Elle peut varier sur plusieurs ordres de grandeur.
\end{definition}

\begin{propriete}[Classification selon la vitesse]
On distingue trois grandes catégories :
\begin{itemize}
\item \textbf{Réactions rapides (quasi instantanées) :} la transformation est quasi immédiate à l'échelle humaine (temps caractéristique $< \SI{1}{s}$). Exemples : réaction de précipitation \ce{Ag+ + Cl- -> AgCl(s)}, neutralisation acide-base \ce{H3O+ + OH- -> 2 H2O}, réaction redox \ce{MnO4- + 5 Fe^2+ + 8 H+ -> Mn^2+ + 5 Fe^3+ + 4 H2O}.
\item \textbf{Réactions lentes :} l'évolution se suit sur plusieurs minutes, heures ou jours. Exemples : réaction \ce{2 I- + S2O8^2- -> I2 + 2 SO4^2-} (modèle de TP), estérification \ce{CH3COOH + C2H5OH <=> CH3COOC2H5 + H2O}, formation de la rouille \ce{4 Fe + 3 O2 + 6 H2O -> 4 Fe(OH)3}.
\item \textbf{Réactions très lentes :} imperceptibles sans catalyseur ou conditions extrêmes. Exemple : synthèse de l'eau \ce{2 H2 + O2 -> 2 H2O} à température ambiante sans catalyseur.
\end{itemize}
\end{propriete}

\begin{savaistu}[Au Cameroun]
La corrosion des armatures en fer des bâtiments (béton armé) est une réaction lente d'oxydation du fer par l'oxygène de l'air en présence d'humidité. La rapidité de cette réaction conditionne la durée de vie des ouvrages d'art (ponts, barrages). La protection par peinture, galvanisation ou acier inoxydable vise à ralentir cette cinétique.
\end{savaistu}

\coursec{2. Vitesses de formation et de disparition}

\courssub{Définitions}

Soit une réaction modélisée par l'équation bilan :
\[ \ce{a A + b B -> c C + d D} \]

\begin{definition}[Vitesse de formation d'un produit]
La \cle{vitesse moyenne de formation} du produit C entre $t_1$ et $t_2$ est :
\[ v_{\text{moy}}(\ce{C}) = \frac{[\ce{C}]_{t_2} - [\ce{C}]_{t_1}}{t_2 - t_1} \quad (\text{en \SI{}{mol.L^{-1}.s^{-1}} \text{ ou } \SI{}{mol.L^{-1}.min^{-1}}}) \]
La \cle{vitesse instantanée de formation} à l'instant $t$ est la limite de la vitesse moyenne quand $\Delta t \to 0$ :
\[ v(\ce{C}) = \frac{d[\ce{C}]}{dt} \]
Graphiquement, c'est la \textbf{pente de la tangente} à la courbe $[\ce{C}] = f(t)$ au point d'abscisse $t$.
\end{definition}

\begin{definition}[Vitesse de disparition d'un réactif]
De même pour le réactif A :
\[ v_{\text{moy}}(\ce{A}) = -\frac{[\ce{A}]_{t_2} - [\ce{A}]_{t_1}}{t_2 - t_1} \qquad v(\ce{A}) = -\frac{d[\ce{A}]}{dt} \]
Le signe moins assure une grandeur positive (la concentration d'un réactif diminue).
\end{definition}

\begin{methode}[Détermination graphique de la vitesse instantanée]
\begin{enumerate}
\item Tracer la courbe $[\ce{Espece}] = f(t)$ en choisissant des échelles adaptées (points bien étalés).
\item Repérer l'instant $t$ sur l'axe des abscisses.
\item Tracer la \textbf{tangente} à la courbe en ce point (règle, compas ou logiciel).
\item Choisir \textbf{deux points éloignés} A et B sur cette tangente (pas sur la courbe !).
\item Calculer la pente : $v = \dfrac{y_B - y_A}{x_B - x_A}$.
\item Exprimer le résultat avec son unité.
\end{enumerate}
\end{methode}

\begin{exemplebox}[Exemple : formation du diiode]
On suit la réaction \ce{S2O8^2- + 2 I- -> 2 SO4^2- + I2}. La courbe $[\ce{I2}] = f(t)$ est croissante, de pente décroissante.
\begin{itemize}
\item Entre $t=0$ et $t=\SI{8}{min}$, $[\ce{I2}]$ passe de 0 à \SI{3,7}{mmol.L^{-1}} : $v_{\text{moy}} = \frac{3,7}{8} \approx \SI{0,46}{mmol.L^{-1}.min^{-1}}$.
\item À $t=\SI{10}{min}$, la tangente a pour pente $\approx \SI{0,28}{mmol.L^{-1}.min^{-1}}$. La vitesse diminue car les concentrations des réactifs baissent.
\end{itemize}
\end{exemplebox}

\begin{attention}[Pièges classiques]
\begin{itemize}
\item Confondre vitesse moyenne (pente d'une \textbf{corde}) et vitesse instantanée (pente de la \textbf{tangente}).
\item Oublier le signe moins pour la vitesse de disparition (la pente de la courbe $[\ce{A}]=f(t)$ est négative, la vitesse est positive).
\item Lire la pente sur la courbe elle-même au lieu de la tangente.
\item Ne pas convertir les unités (min en s, mmol en mol) si l'énoncé l'impose.
\end{itemize}
\end{attention}

\coursec{3. Vitesse volumique, avancement et temps de demi-réaction}

\courssub{Vitesse volumique de réaction}

\begin{definition}[Vitesse volumique $v$]
Pour une réaction \ce{a A + b B -> c C + d D}, la \cle{vitesse volumique de réaction} (ou simplement vitesse de réaction) est définie par :
\[ v = \frac{1}{c}\frac{d[\ce{C}]}{dt} = \frac{1}{d}\frac{d[\ce{D}]}{dt} = -\frac{1}{a}\frac{d[\ce{A}]}{dt} = -\frac{1}{b}\frac{d[\ce{B}]}{dt} \]
Elle est unique pour une réaction donnée à un instant donné (unité : \SI{}{mol.L^{-1}.s^{-1}}).
\end{definition}

\begin{propriete}[Lien avec l'avancement volumique $x$]
On définit l'\cle{avancement volumique} $x$ (en \SI{}{mol.L^{-1}}) par :
\[ x = \frac{[\ce{C}] - [\ce{C}]_0}{c} = \frac{[\ce{A}]_0 - [\ce{A}]}{a} \]
Alors $v = \frac{dx}{dt}$. Le tableau d'avancement s'écrit en concentrations :
\begin{center}
\begin{tabularx}{\linewidth}{|c|c|c|c|c|}
\hline
\rowcolor{popblue!60} Espèce & A & B & C & D \\
\hline
État initial & $[\ce{A}]_0$ & $[\ce{B}]_0$ & $[\ce{C}]_0$ & $[\ce{D}]_0$ \\
\hline
État à l'instant $t$ & $[\ce{A}]_0 - a x$ & $[\ce{B}]_0 - b x$ & $[\ce{C}]_0 + c x$ & $[\ce{D}]_0 + d x$ \\
\hline
État final & $[\ce{A}]_0 - a x_f$ & $[\ce{B}]_0 - b x_f$ & $[\ce{C}]_0 + c x_f$ & $[\ce{D}]_0 + d x_f$ \\
\hline
\end{tabularx}
\end{center}
\end{propriete}

\courssub{Temps de demi-réaction}

\begin{definition}[Temps de demi-réaction $t_{1/2}$]
C'est le temps nécessaire pour que la concentration d'un réactif (souvent le réactif limitant) soit divisée par 2, ou pour que l'avancement atteigne la moitié de sa valeur finale $x_f$ :
\[ t_{1/2} \text{ tel que } x(t_{1/2}) = \frac{x_f}{2} \quad \text{ou} \quad [\ce{A}](t_{1/2}) = \frac{[\ce{A}]_0}{2} \text{ (si } [\ce{A}]_f=0\text{)} \]
\end{definition}

\begin{aretenir}[Propriétés du $t_{1/2}$]
\begin{itemize}
\item Pour une réaction d'ordre 1 par rapport à un réactif : $t_{1/2}$ est \textbf{constant} (indépendant de la concentration initiale). C'est le cas de la radioactivité et de nombreuses réactions complexes en pseudo-ordre 1.
\item Pour les autres ordres : $t_{1/2}$ dépend de la concentration initiale.
\item Graphiquement : on repère $[\ce{A}]_0/2$ (ou $x_f/2$) sur l'ordonnée et on reporte sur l'abscisse via la courbe.
\end{itemize}
\end{aretenir}

\coursec{4. Suivi temporel d'une réaction}

\courssub{Méthodes de suivi}

\begin{propriete}[Deux grandes familles]
\begin{itemize}
\item \textbf{Méthode des prélèvements avec trempe (quotientage) :} on prélève un aliquot du mélange réactionnel à des dates connues et on arrête brutalement la réaction (\cle{trempe}) : refroidissement brutal (eau glacée), dilution massive, ajout d'un réactif qui consomme un réactif, inactivation enzymatique. Puis on dose l'espèce suivie (titrage, spectrophotométrie sur le prélèvement). \emph{Exemple type : réaction \ce{I-}/\ce{S2O8^2-} dosée au thiosulfate.}
\item \textbf{Méthodes physiques \emph{in situ} (notion) :} suivi continu sans prélever.
\begin{itemize}
\item \textbf{Colorimétrie / Spectrophotométrie :} loi de Beer-Lambert $A = \varepsilon \ell c$. On suit l'absorbance $A$ à une longueur d'onde où seul un produit ou réactif absorbe. $A \propto c$.
\item \textbf{Conductimétrie :} on suit la conductivité $\sigma$ du mélange. $\sigma = \sum \lambda_i [\ce{X_i}]$. Utile si le nombre d'ions ou leur mobilité change (ex: \ce{H+} remplacé par \ce{Na+}).
\item \textbf{Manométrie :} pour les réactions gazeuses avec changement de nombre de moles gazeuses. $P \propto n_{\text{gaz}}$.
\end{itemize}
\end{itemize}
\end{propriete}

\begin{experience}[TP : Suivi cinétique de la réaction \protect\ce{I-} / \protect\ce{S2O8^2-} par prélèvements et trempe]
\courssub{Objectifs}
\begin{itemize}
\item Réaliser un suivi temporel par la méthode des prélèvements avec trempe.
\item Doser le diiode formé par une solution étalonnée de thiosulfate de sodium.
\item Tracer $[\ce{I2}] = f(t)$, déterminer $v_{\text{moy}}$, $v_{\text{inst}}$ (tangente) et $t_{1/2}$.
\item Mettre en évidence le rôle catalytique des ions \ce{Fe^3+}.
\item Respecter les consignes de sécurité (produits oxydants, diiode).
\end{itemize}

\courssub{Rappel théorique}
La réaction est lente, totale : 
\[ \ce{S2O8^2- + 2 I- -> 2 SO4^2- + I2} \quad (\text{équation bilan}) \]
Le dosage du diiode utilise la réaction rapide et totale :
\[ \ce{I2 + 2 S2O3^2- -> 2 I- + S4O6^2-} \]

\courssub{Matériel et produits}
\begin{center}
\begin{tabularx}{\linewidth}{|l|X|}
\hline
\rowcolor{popblue!60} \textbf{Matériel} & \textbf{Produits} \\
\hline
\rowcolor{popblueL} 1 burette graduée \SI{25}{mL} (classe B) & \ce{KI} \SI{0,20}{mol.L^{-1}} \\
\hline
2 pipettes jaugées \SI{10}{mL} et \SI{20}{mL} (classe A) & \ce{K2S2O8} \SI{0,10}{mol.L^{-1}} \\
\hline
\rowcolor{popblueL} 1 erlenmeyer \SI{250}{mL} & \ce{Na2S2O3} \SI{0,010}{mol.L^{-1}} (titrée) \\
\hline
7 béchers \SI{100}{mL} & \ce{FeCl3} \SI{0,10}{mol.L^{-1}} \\
\hline
\rowcolor{popblueL} 1 agitateur en verre, 1 chronomètre, 1 thermomètre & Empois d'amidon (ou alternative locale) \\
\hline
Bac à glace (glace pilée + eau) & Eau distillée \\
\hline
\end{tabularx}
\end{center}
\end{experience}

\begin{savaistu}[Alternatives locales - Empois d'amidon]
L'empois d'amidon commercial peut être remplacé par une décoction de farine de maïs ou de manioc : faire bouillir \SI{5}{g} de farine dans \SI{50}{mL} d'eau, filtrer sur coton. Le filtrat trouble donne le virage bleu nuit caractéristique avec le diiode (complexe amidon-iode). Le jus de bissap (hibiscus) ne convient pas (indicateur pH, pas complexe iode).
\end{savaistu}

\begin{attention}[Consignes strictes]
\begin{itemize}
\item \textbf{Oxydant / Corrosif :} \ce{K2S2O8} et \ce{FeCl3} irritent peau et yeux. \textbf{Blouse, lunettes, gants} obligatoires.
\item \textbf{Interdiction} de pipeter à la bouche (propipette obligatoire).
\item \ce{I2} tache peau et vêtements (brun tenace). Manipuler au-dessus des béchers.
\item En cas de projection : rinçage \SI{15}{min} à l'eau, signaler à l'enseignant.
\end{itemize}
\end{attention}

\begin{popfigure}
\begin{center}
\begin{tikzpicture}[scale=0.9, every node/.style={font=\small}]
% Couleurs
\colorlet{liq}{poporange!40}
\colorlet{ice}{popblue!20}
\colorlet{titr}{popgold!40}
% Erlenmeyer réaction
\begin{scope}[shift={(0,0)}]
\draw[thick,popdark] (-1.2,0) -- (-0.4,2.5) -- (-0.4,3.2);
\draw[thick,popdark] (1.2,0) -- (0.4,2.5) -- (0.4,3.2);
\draw[thick,popdark] (-1.2,0) -- (1.2,0);
\draw[thick,popdark] (-0.6,3.2) -- (0.6,3.2);
\fill[liq] (-1.06,0.1) -- (-0.48,2.3) -- (0.48,2.3) -- (1.06,0.1) -- cycle;
\draw[popdark] (-0.48,2.3) -- (0.48,2.3);
\node[popdark,align=center] at (0,-0.7) {Erlenmeyer \SI{250}{mL}};
\node[popdark,align=center] at (0,-1.2) {Mélange \ce{KI} + \ce{K2S2O8}};
\node[popdark,align=center] at (0,-1.6) {$T \approx \SI{20}{\celsius}$};
\draw[-{Stealth[length=3mm]},thick,popblue] (1.5,2.5) -- (2.5,2.0);
\node[popblue,align=center,anchor=west] at (2.6,2.0) {Pipette \SI{10}{mL} : prélèvement\\ à la date $t$};
\end{scope}
% Bécher trempe
\begin{scope}[shift={(6.5,0)}]
\draw[thick,popdark] (-1.1,0) -- (-1.1,2.2) -- (1.1,2.2) -- (1.1,0) -- cycle;
\fill[ice] (-1.05,0.05) rectangle (1.05,1.5);
\draw[popdark,dashed] (-1.05,1.5) -- (1.05,1.5);
\draw[popblue!60,thick] (-0.7,1.8) -- (-0.4,1.55) -- (-0.15,1.85);
\draw[popblue!60,thick] (0.1,1.9) -- (0.35,1.6) -- (0.6,1.9);
\node[popdark,align=center] at (0,-0.7) {Bécher trempe n°i};
\node[popdark,align=center] at (0,-1.2) {Eau glacée $\approx \SI{5}{\celsius}$};
\node[popblue,align=center] at (0,0.8) {Prélèvement figé};
\end{scope}
% Burette + Erlenmeyer dosage
\begin{scope}[shift={(12.5,0)}]
% Burette
\draw[thick,popdark] (-0.15,0.5) -- (-0.15,3.8) -- (0.15,3.8) -- (0.15,0.5);
\fill[titr] (-0.13,0.55) rectangle (0.13,3.0);
\draw[thick,popdark] (-0.15,0.5) -- (0,0.1) -- (0.15,0.5);
\draw[thick,popdark] (-0.4,0.2) -- (0.4,0.2);
\draw[thick,popdark] (0.4,0.2) -- (0.4,0.45);
\draw[thick,popdark] (-1.8,3.8) -- (-0.15,3.8);
\draw[thick,popdark] (-1.8,0.0) -- (-1.8,3.8);
\draw[thick,popdark] (-2.3,0.0) -- (1.0,0.0);
\foreach \y in {1.0,1.5,2.0,2.5,3.0} \draw[popdark] (0.15,\y) -- (0.35,\y);
\node[popgold!80!black,anchor=east] at (-1.9,3.9) {\ce{Na2S2O3} \SI{0,010}{M}};
% Erlenmeyer sous burette
\draw[thick,popdark] (-0.8,-1.8) -- (-0.25,-0.2) -- (-0.25,0.1);
\draw[thick,popdark] (0.8,-1.8) -- (0.25,-0.2) -- (0.25,0.1);
\draw[thick,popdark] (-0.8,-1.8) -- (0.8,-1.8);
\fill[liq] (-0.7,-1.75) -- (-0.33,-0.6) -- (0.33,-0.6) -- (0.7,-1.75) -- cycle;
\node[popdark,align=center] at (0,-2.4) {Dosage du \ce{I2}};
\node[popdark,align=center] at (0,-2.9) {par \ce{S2O3^2-}};
\end{scope}
\end{tikzpicture}
\end{center}
\legende{Dispositif : 1. Mélange réactionnel thermostaté. 2. Prélèvement à la pipette jaugée. 3. Trempe immédiate dans eau glacée (arrêt cinétique). 4. Dosage du \ce{I2} figé au thiosulfate à la burette.}
\end{popfigure}

\courssub{Protocole détaillé}

\textbf{Partie A — Suivi cinétique sans catalyseur}
\begin{enumerate}
\item Préparer 7 béchers (\SI{100}{mL}) avec \SI{\approx 40}{mL} d'eau glacée (\SI{\approx 5}{\celsius}). Numéroter 1 à 7.
\item Dans l'erlenmeyer \SI{250}{mL}, introduire \SI{50,0}{mL} de \ce{KI} \SI{0,20}{mol.L^{-1}} (pipette jaugée \SI{20}{mL} $\times$ 2 + \SI{10}{mL} ou pipette \SI{50}{mL} si dispo). Noter $T_i \approx \SI{20}{\celsius}$.
\item Remplir la burette de \ce{Na2S2O3} \SI{0,010}{mol.L^{-1}}, faire le zéro (ménisque bas).
\item \textbf{Instant $t=0$ :} Verser rapidement \SI{50,0}{mL} de \ce{K2S2O8} \SI{0,10}{mol.L^{-1}} dans l'erlenmeyer. Homogénéiser (agitateur). \textbf{Démarrer le chronomètre.}
\item À $t = \SI{2}{min}$, prélever \SI{10,0}{mL} (pipette jaugée) du mélange, verser \textbf{immédiatement} dans le bécher n°1 (trempe). La réaction est quasi arrêtée.
\item Dosage : Titrer le contenu du bécher n°1 par le thiosulfate jusqu'à décoloration de la teinte jaune/brune. \textbf{En fin de dosage} (teinte jaune pâle), ajouter \SI{2}{mL} d'empois d'amidon : le bleu nuit apparaît. Continuer le dosage goutte à goutte jusqu'à disparition du bleu (incolore). Noter $V_{E,1}$.
\item Reproduire les étapes 5-6 aux dates : \SI{4}{min}, \SI{6}{min}, \SI{8}{min}, \SI{12}{min}, \SI{16}{min}, \SI{20}{min} (béchers 2 à 7).
\end{enumerate}

\textbf{Partie B — Rôle catalytique des ions \ce{Fe^3+}}
\begin{enumerate}
\item Rincer l'erlenmeyer. Y mettre \SI{50,0}{mL} \ce{KI} \SI{0,20}{M} + \textbf{5 gouttes} de \ce{FeCl3} \SI{0,10}{M}.
\item Verser \SI{50,0}{mL} \ce{K2S2O8} \SI{0,10}{M}, démarrer le chrono.
\item Observer l'apparition de la coloration brune (comparer à Partie A). Si temps dispo, doser à $t=\SI{2}{min}$ et \SI{4}{min}.
\end{enumerate}

\courssub{Calculs et exploitation des résultats}

\begin{propriete}[Calcul de la concentration en diiode]
Au point d'équivalence du dosage : $n(\ce{I2}) = \frac{1}{2} n(\ce{S2O3^2-})$.
\[ n(\ce{S2O3^2-}) = C_E \times V_E \quad (C_E \text{ en mol/L}, V_E \text{ en L}) \]
Le prélèvement de volume $V_{\text{prél}} = \SI{10,0}{mL}$ contient $n(\ce{I2})$. Sa concentration dans le mélange réactionnel (avant trempe) est :
\[ [\ce{I2}] = \frac{n(\ce{I2})}{V_{\text{prél}}} = \frac{C_E \times V_E}{2 \times V_{\text{prél}}} \]
Avec $C_E = \SI{0,010}{mol.L^{-1}}$, $V_{\text{prél}} = \SI{10,0}{mL} = \SI{0,0100}{L}$ :
\[ [\ce{I2}] (\text{en mol/L}) = \frac{0,010 \times V_E(\text{mL})}{2 \times 10,0} = 5,0 \times 10^{-4} \times V_E(\text{mL}) \]
Soit en \SI{}{mmol.L^{-1}} : $[\ce{I2}] = 0,5 \times V_E(\text{mL})$.
\end{propriete}

\begin{center}
\begin{tabularx}{\linewidth}{|c|c|c|X|}
\hline
\rowcolor{popgreen!60} $t$ (min) & $V_E$ (mL) & $[\ce{I2}]$ (mmol.L$^{-1}$) & Aspect \\
\hline
\rowcolor{popgreenL} 0 & 0,0 & 0,0 & Incolore \\
\hline
2 & 2,4 & 1,2 & Jaune très pâle \\
\hline
\rowcolor{popgreenL} 4 & 4,4 & 2,2 & Jaune \\
\hline
6 & 6,0 & 3,0 & Orange \\
\hline
\rowcolor{popgreenL} 8 & 7,4 & 3,7 & Brun clair \\
\hline
12 & 9,4 & 4,7 & Brun \\
\hline
\rowcolor{popgreenL} 16 & 10,8 & 5,4 & Brun foncé \\
\hline
20 & 11,8 & 5,9 & Brun très foncé \\
\hline
\end{tabularx}
\end{center}

\begin{aretenir}[Interprétation des concentrations initiales et limite théorique]
Dans le mélange initial (\SI{100}{mL} total) :
$[\ce{S2O8^2-}]_0 = \SI{0,10}{mol.L^{-1}} \times \frac{50}{100} = \SI{0,050}{mol.L^{-1}} = \SI{50}{mmol.L^{-1}}$.
$[\ce{I-}]_0 = \SI{0,20}{mol.L^{-1}} \times \frac{50}{100} = \SI{0,10}{mol.L^{-1}} = \SI{100}{mmol.L^{-1}}$.
$\ce{S2O8^2-}$ est limitant. L'avancement final maximal $x_f = \SI{50}{mmol.L^{-1}}$.
Donc $[\ce{I2}]_{\infty, \text{théo}} = \SI{50}{mmol.L^{-1}}$.
Le tableau ne couvre que les \textbf{premiers pourcents} de la réaction (très lente à \SI{20}{\celsius} sans catalyseur). L'asymptote \SI{50}{mmol.L^{-1}} n'est pas atteinte en 20 min.
\end{aretenir}

\courssub{Exploitation graphique et numérique}
\begin{exoresolu}[Exploitation complète des résultats]
\textbf{Énoncé.} À partir du tableau : (1) Tracer $[\ce{I2}] = f(t)$. (2) $v_{\text{moy}}$ de formation de \ce{I2} entre 0 et \SI{8}{min}. (3) $v_{\text{inst}}$ à $t=\SI{10}{min}$. (4) Estimer $t_{1/2}$ (réactif limitant). (5) Conclure sur \ce{Fe^3+}.
\tcblower
\textbf{Solution.}

\textbf{(1) Courbe.} Abscisses : $t$ (\SI{1}{cm} pour \SI{2}{min}). Ordonnées : $[\ce{I2}]$ (\SI{1}{cm} pour \SI{1}{mmol.L^{-1}}). Points alignés sur une courbe concave, pente décroissante.

\textbf{(2) Vitesse moyenne $0 \to \SI{8}{min}$ :}
\[ v_{\text{moy}} = \frac{[\ce{I2}]_8 - [\ce{I2}]_0}{8 - 0} = \frac{3,7 - 0}{8} = \SI{0,46}{mmol.L^{-1}.min^{-1}} \]

\textbf{(3) Vitesse instantanée à $t=\SI{10}{min}$ :}
Tracer la tangente en $t=10$. Choisir deux points sur la tangente : ex. A$(4 ; 1,0)$ et B$(16 ; 4,3)$.
\[ v(10) = \frac{4,3 - 1,0}{16 - 4} = \frac{3,3}{12} = \SI{0,275}{mmol.L^{-1}.min^{-1}} \approx \SI{0,28}{mmol.L^{-1}.min^{-1}} \]
$v(10) < v_{\text{moy}}(0\to8)$ : la vitesse \textbf{décroît} avec le temps (consommation des réactifs).

\textbf{(4) Temps de demi-réaction $t_{1/2}$ :}
Réactif limitant : \ce{S2O8^2-}. $[\ce{S2O8^2-}]_0 = \SI{50}{mmol.L^{-1}}$.
$t_{1/2}$ correspond à $[\ce{S2O8^2-}] = \SI{25}{mmol.L^{-1}}$, soit $x = \SI{25}{mmol.L^{-1}}$, donc $[\ce{I2}] = \SI{25}{mmol.L^{-1}}$.
\textbf{Attention :} Notre courbe s'arrête à \SI{5,9}{mmol.L^{-1}} ($t=20$ min). On ne peut \textbf{pas} lire $t_{1/2}$ sur cette courbe (extrapolation hasardeuse).
\textit{Si l'on utilisait à tort l'asymptote apparente du tableau (\SI{6,7}{mmol.L^{-1}}), on trouverait $t_{1/2} \approx \SI{7}{min}$ pour $[\ce{I2}]=\SI{3,35}{mmol.L^{-1}}$, ce qui n'est pas le vrai $t_{1/2}$ de la réaction totale.}

\textbf{(5) Rôle des ions \ce{Fe^3+} :}
Avec \ce{Fe^3+}, la coloration brune apparaît \textbf{beaucoup plus vite}. À date égale, $V_E$ est plus grand $\Rightarrow$ $[\ce{I2}]$ plus grand $\Rightarrow$ vitesse plus grande.
Les ions \ce{Fe^3+} ne figurent pas dans le bilan, ne sont pas consommés (régénérés) : ce sont des \textbf{catalyseurs homogènes}.
Mécanisme (notion) :
\begin{align*}
\ce{S2O8^2- + 2 Fe^2+ &-> 2 SO4^2- + 2 Fe^3+} \quad (\text{rapide}) \\
\ce{2 Fe^3+ + 2 I- &-> 2 Fe^2+ + I2} \quad (\text{rapide})
\end{align*}
Bilan : \ce{S2O8^2- + 2 I- -> 2 SO4^2- + I2}. Le couple \ce{Fe^3+}/\ce{Fe^2+} abaisse l'énergie d'activation.
\end{exoresolu}

\coursec{5. Facteurs cinétiques et catalyse}

\courssub{Influence de la concentration des réactifs}
\begin{propriete}[Loi de vitesse (notion)]
Pour une réaction élémentaire ou modélisée : $v = k [\ce{A}]^\alpha [\ce{B}]^\beta$.
$\alpha, \beta$ = ordres partiels (entiers $\ge 0$ le plus souvent). $k$ = constante de vitesse (dépend de $T$).
Augmenter la concentration des réactifs augmente la fréquence des chocs efficaces $\Rightarrow$ augmente $v$.
\end{propriete}

\courssub{Influence de la température}
\begin{propriete}[Règle empirique et énergie d'activation]
Une augmentation de \SI{10}{\celsius} multiplie la vitesse par un facteur 2 à 3 (règle de van't Hoff).
Origine microscopique : répartition de Maxwell-Boltzmann. Seules les molécules d'énergie $\ge E_a$ (énergie d'activation) réagissent. Augmenter $T$ augmente la fraction de molécules "réactives".
\end{propriete}

\begin{popfigure}
\begin{center}
\begin{tikzpicture}[scale=0.85, domain=0:5, samples=200]
\begin{axis}[
width=\linewidth, height=7cm,
axis lines=middle, xlabel={Énergie cinétique $E_c$}, ylabel={Nombre de molécules $N$},
xtick=\empty, ytick=\empty,
xmin=0, xmax=5.5, ymin=0, ymax=2.5,
clip=false
]
% Courbe T1
\addplot[popblue, thick, domain=0:5] {2*x*exp(-x)};
\addplot[popblue, dashed, thick, domain=2.5:5] {2*x*exp(-x)};
\node[popblue, anchor=south] at (axis cs:1.5,1.8) {$T_1$};
% Courbe T2 > T1
\addplot[poporange, thick, domain=0:5] {1.5*x*exp(-0.7*x)};
\addplot[poporange, dashed, thick, domain=2.5:5] {1.5*x*exp(-0.7*x)};
\node[poporange, anchor=south] at (axis cs:2.2,1.6) {$T_2 > T_1$};
% Zone Ea
\draw[poppurple, thick] (axis cs:2.5,0) -- (axis cs:2.5,2.2);
\node[poppurple, anchor=west] at (axis cs:2.5,2.3) {$E_a$};
\fill[poppurple!20] (axis cs:2.5,0) -- (axis cs:5,0) -- (axis cs:5, {1.5*5*exp(-0.7*5)}) -- (axis cs:2.5, {1.5*2.5*exp(-0.7*2.5)}) -- cycle;
\fill[poppurple!10] (axis cs:2.5,0) -- (axis cs:5,0) -- (axis cs:5, {2*5*exp(-5)}) -- (axis cs:2.5, {2*2.5*exp(-2.5)}) -- cycle;
\node[poppurple, align=center] at (axis cs:4.0,0.8) {Fraction réactive\\ plus grande à $T_2$};
\end{axis}
\end{tikzpicture}
\end{center}
\legende{Répartition des énergies cinétiques : à température plus élevée, plus de molécules dépassent l'énergie d'activation $E_a$ (zone hachurée).}
\end{popfigure}

\courssub{La catalyse}

\begin{definition}[Catalyseur]
Un \cle{catalyseur} est une espèce chimique qui augmente la vitesse d'une réaction sans apparaître dans son équation bilan. Il n'est pas consommé (régénéré). Il abaisse l'énergie d'activation $E_a$ en fournissant un nouveau mécanisme réactionnel.
\end{definition}

\begin{propriete}[Types de catalyse]
\begin{itemize}
\item \textbf{Catalyse homogène :} catalyseur dans la même phase que les réactifs.
Exemple : ions \ce{Fe^3+} (aq) pour \ce{I-}/\ce{S2O8^2-}. Couple \ce{Fe^3+}/\ce{Fe^2+}.
\item \textbf{Catalyse hétérogène :} catalyseur dans une phase différente (souvent solide).
Exemple : Platine (Pt) pour la synthèse de l'eau \ce{2 H2(g) + O2(g) ->[Pt] 2 H2O(l)}. Adsorption des réactifs sur la surface du métal.
\item \textbf{Catalyse enzymatique :} catalyseur = enzyme (protéine). Spécificité extrême, conditions douces (T, pH). Exemple : invertase (saccharose $\to$ glucose + fructose), amylase (amidon $\to$ maltose). Importance : digestion, brassage (bière, bili-bili), pharmacopée.
\end{itemize}
\end{propriete}

\begin{popfigure}
\begin{center}
\begin{tikzpicture}[scale=0.9]
\begin{axis}[
width=\linewidth, height=7cm,
axis lines=middle,
xlabel={Coordonnée de réaction}, ylabel={Énergie potentielle $E_p$},
xtick=\empty, ytick=\empty,
xmin=-0.5, xmax=5.5, ymin=-0.5, ymax=4.5,
]
% Réaction sans catalyseur
\draw[popdark, thick, samples=100, domain=0:5] plot (\x, {3*exp(-(\x-1.5)^2/0.5) + 0.5*exp(-(\x-3.5)^2/0.5)});
\draw[popdark, dashed] (0,1) -- (0.5,1) node[left] {Réactifs};
\draw[popdark, dashed] (5,1.5) -- (4.5,1.5) node[right] {Produits};
\node[popdark] at (1.5, 3.8) {$E_a$ (sans cat.)};
\draw[<->, popdark] (1.5,1) -- (1.5,3.5);
% Réaction avec catalyseur
\draw[popgreen!70!black, thick, samples=100, domain=0:5] plot (\x, {1.5*exp(-(\x-1.5)^2/0.8) + 0.3*exp(-(\x-3.5)^2/0.5) + 0.5});
\node[popgreen!70!black] at (1.5, 2.2) {$E_a'$ (avec cat.)};
\draw[<->, popgreen!70!black] (1.5,1) -- (1.5,1.8);
% Intermédiaire
\draw[poppurple, dashed] (3.5, 1.2) -- (3.5, 1.8);
\node[poppurple, anchor=north] at (3.5, 1.8) {Intermédiaire};
\end{axis}
\end{tikzpicture}
\end{center}
\legende{Profil énergétique : le catalyseur abaisse l'énergie d'activation ($E_a' < E_a$) en stabilisant un intermédiaire réactionnel. L'enthalpie de réaction $\Delta_r H$ est inchangée.}
\end{popfigure}

\courssub{TP : Mise en évidence de la catalyse par \ce{Fe^3+} (Partie B du TP précédent)}
\begin{exoresolu}[Analyse de l'effet catalytique]
\textbf{Énoncé.} On compare les courbes $[\ce{I2}] = f(t)$ avec et sans \ce{Fe^3+}. La courbe avec \ce{Fe^3+} est plus pentue au départ et atteint plus vite des concentrations élevées.
\tcblower
\textbf{Solution.}
\begin{itemize}
\item \textbf{Observation :} Apparition de la couleur brune quasi immédiate avec \ce{Fe^3+} (quelques secondes) vs plusieurs minutes sans.
\item \textbf{Interprétation :} À $t$ donné, $[\ce{I2}]_{\text{avec cat.}} > [\ce{I2}]_{\text{sans cat.}} \Rightarrow v_{\text{avec cat.}} > v_{\text{sans cat.}}$.
\item \textbf{Nature :} \ce{Fe^3+} est un catalyseur \textbf{homogène} (même phase aqueuse). Le cycle catalytique implique le couple redox \ce{Fe^3+}/\ce{Fe^2+}.
\item \textbf{Application industrielle :} Ce principe est utilisé dans certaines synthèses d'oxydants ou le traitement des eaux.
\end{itemize}
\end{exoresolu}

\coursec{6. Étude cinétique de l'estérification}

\courssub{Réaction modèle}
L'estérification de l'acide éthanoïque par l'éthanol est une réaction \textbf{lente, limitée} (équilibre), catalysée par les ions \ce{H3O+} (catalyse acide homogène).
\[ \ce{CH3COOH_{(aq)} + C2H5OH_{(aq)} <=>[H2SO4] CH3COOC2H5_{(aq)} + H2O_{(l)}} \]
\begin{center}
\chemfig{CH_3-C(=[2]O)-[6]OH} + \chemfig{CH_3-CH_2-[6]OH} <=> \chemfig{CH_3-C(=[2]O)-[6]O-CH_2-CH_3} + \chemfig{H_2O}
\end{center}
Acide éthanoïque + Éthanol $\rightleftharpoons$ Éthanoate d'éthyle + Eau.

\courssub{Suivi cinétique : méthode des prélèvements (quotientage)}
\begin{methode}[Protocole type]
\begin{enumerate}
\item Mélanger volumes connus d'acide éthanoïque et d'éthanol (souvent excès d'éthanol pour pseudo-ordre 1). Ajouter catalyseur (\ce{H2SO4} concentré, quelques gouttes).
\item À $t=0$, prélever un aliquot, tremper dans eau glacée + eau distillée (dilution + arrêt).
\item Dosage de l'acide résiduel par une solution de base étalonnée (\ce{NaOH}) au phénolphtaléine.
\item Répéter à intervalles réguliers jusqu'à constance du volume de base (équilibre).
\end{enumerate}
\end{methode}

\begin{propriete}[Lien dosage / avancement]
Soit $V_E$ le volume de \ce{NaOH} (concentration $C_B$) versé pour neutraliser l'acide résiduel dans le prélèvement $V_{\text{prél}}$.
$n(\ce{CH3COOH})_{\text{résiduel}} = C_B V_E$.
Concentration acide dans le mélange : $[\ce{CH3COOH}] = \frac{C_B V_E}{V_{\text{prél}}}$.
Avancement $x = [\ce{CH3COOH}]_0 - [\ce{CH3COOH}]$.
À l'équilibre ($t \to \infty$) : $V_E = V_{E,\infty}$, $x = x_f$.
\end{propriete}

\courssub{Courbes cinétiques et facteurs}

\begin{popfigure}
\begin{center}
\begin{tikzpicture}[scale=0.9]
\begin{axis}[
width=\linewidth, height=8cm,
axis lines=middle,
xlabel={Temps $t$}, ylabel={$x$ (mol.L$^{-1}$) ou $[\ce{Produit}]$},
xtick=\empty, ytick=\empty,
xmin=0, xmax=10, ymin=0, ymax=4.5,
]
% Courbe référence (T1, sans cat)
\addplot[popblue, thick, domain=0:10, samples=200] {3*(1-exp(-0.5*x))};
\node[popblue, anchor=south west] at (axis cs:8, 2.8) {$T_1$, sans cat.};
% Courbe T2 > T1
\addplot[poporange, thick, domain=0:10, samples=200] {3*(1-exp(-1.2*x))};
\node[poporange, anchor=south west] at (axis cs:5, 2.9) {$T_2 > T_1$, sans cat.};
% Courbe avec catalyseur (même T1)
\addplot[popgreen!70!black, thick, domain=0:10, samples=200] {3*(1-exp(-1.5*x))};
\node[popgreen!70!black, anchor=north west] at (axis cs:2, 3.2) {$T_1$, avec \ce{H2SO4}};
% Asymptote
\draw[popmuted, dashed, thick] (0,3) -- (10,3);
\node[popmuted, anchor=south] at (10,3) {$x_f$ (inchangé)};
\end{axis}
\end{tikzpicture}
\end{center}
\legende{Courbes d'avancement $x(t)$ pour l'estérification. L'asymptote $x_f$ (équilibre) est la même. La température et le catalyseur accélèrent l'atteinte de l'équilibre sans le déplacer.}
\end{popfigure}

\begin{aretenir}[Enseignements clés de l'estérification]
\begin{itemize}
\item Réaction \textbf{limitée} : courbe $x(t)$ tend vers une asymptote horizontale $x_f$ (équilibre chimique).
\item \textbf{Facteurs cinétiques :} Augmenter $T$ ou ajouter catalyseur (\ce{H2SO4}) $\Rightarrow$ pente initiale plus forte, asymptote atteinte \textbf{plus vite}.
\item \textbf{Facteur thermodynamique :} Ni $T$ (si $\Delta_r H \neq 0$, léger déplacement), ni catalyseur ne modifient la position de l'équilibre ($x_f$ constant à $T$ constant). Le catalyseur accélère \textbf{les deux sens} (direct et inverse) de même facteur.
\end{itemize}
\end{aretenir}

\begin{savaistu}[Esters au Cameroun]
Les esters sont responsables des arômes fruités. L'éthanoate d'éthyle sent la poire/la colle. Dans le bili-bili (bière de mil) ou le vin de palme, les esters formés pendant la fermentation (ex: acétate d'éthyle, acétate d'isoamyle) donnent les notes fruitées caractéristiques. La maîtrise de la température de fermentation contrôle la cinétique de formation de ces esters.
\end{savaistu}

\coursec{Exercices d'entraînement (Bac blanc)}

\begin{exercice}[Exercice 1 : Vitesse et avancement]
On étudie la réaction : \ce{N2(g) + 3 H2(g) -> 2 NH3(g)}.
À $t=0$, on a $n_0(\ce{N2}) = \SI{2,0}{mol}$, $n_0(\ce{H2}) = \SI{6,0}{mol}$ dans un réacteur de volume $V = \SI{10,0}{L}$.
À $t = \SI{10}{min}$, on mesure $n(\ce{NH3}) = \SI{0,80}{mol}$.
\begin{enumerate}
\item Calculer l'avancement $x$ (en mol) à $t=\SI{10}{min}$.
\item En déduire $n(\ce{N2})$ et $n(\ce{H2})$ à cet instant.
\item Calculer la vitesse volumique moyenne $v_{\text{moy}}$ entre 0 et \SI{10}{min}.
\item Exprimer les vitesses de formation de \ce{NH3} et de disparition de \ce{H2} en \SI{}{mol.L^{-1}.min^{-1}}.
\end{enumerate}
\end{exercice}

\begin{corrige}[Exercice 1]
\begin{enumerate}
\item $x = \frac{n(\ce{NH3}) - 0}{2} = \frac{0,80}{2} = \SI{0,40}{mol}$.
\item $n(\ce{N2}) = 2,0 - 0,40 = \SI{1,60}{mol}$ ; $n(\ce{H2}) = 6,0 - 3\times0,40 = \SI{4,80}{mol}$.
\item $v_{\text{moy}} = \frac{x}{V \Delta t} = \frac{0,40}{10,0 \times 10} = \SI{4,0 \times 10^{-3}}{mol.L^{-1}.min^{-1}}$.
\item $v_{\text{moy}}(\ce{NH3}) = 2 v_{\text{moy}} = \SI{8,0 \times 10^{-3}}{mol.L^{-1}.min^{-1}}$.
$v_{\text{moy}}(\ce{H2}) = 3 v_{\text{moy}} = \SI{1,2 \times 10^{-2}}{mol.L^{-1}.min^{-1}}$.
\end{enumerate}
\end{corrige}

\begin{exercice}[Exercice 2 : Suivi spectrophotométrique]
On suit la réaction \ce{A + B -> C} par spectrophotométrie à $\lambda = \SI{450}{nm}$. Seule l'espèce C absorbe ($\varepsilon = \SI{2000}{L.mol^{-1}.cm^{-1}}$, $\ell = \SI{1,00}{cm}$).
On relève : $t=0, A=0,00$ ; $t=\SI{5}{min}, A=0,40$ ; $t=\SI{10}{min}, A=0,64$ ; $t=\SI{20}{min}, A=0,80$ ; $t=\SI{30}{min}, A=0,88$.
\begin{enumerate}
\item Calculer $[\ce{C}]$ à chaque instant.
\item Tracer $[\ce{C}] = f(t)$ (qualitatif). La réaction est-elle totale ou limitée ?
\item Déterminer $v_{\text{moy}}(\ce{C})$ entre 0 et \SI{5}{min} et $v_{\text{inst}}(\ce{C})$ à $t=\SI{10}{min}$ (tangente).
\end{enumerate}
\end{exercice}

\begin{corrige}[Exercice 2]
\begin{enumerate}
\item Loi de Beer-Lambert : $[\ce{C}] = \frac{A}{\varepsilon \ell}$.
$t=0: 0$ ; $t=5: \frac{0,40}{2000} = \SI{2,0 \times 10^{-4}}{M}$ ; $t=10: \SI{3,2 \times 10^{-4}}{M}$ ; $t=20: \SI{4,0 \times 10^{-4}}{M}$ ; $t=30: \SI{4,4 \times 10^{-4}}{M}$.
\item Courbe croissante, pente décroissante, tend vers une asymptote $\approx \SI{4,8 \times 10^{-4}}{M}$. Réaction \textbf{limitée} (équilibre).
\item $v_{\text{moy}}(0\to5) = \frac{2,0 - 0}{5} = \SI{4,0 \times 10^{-4}}{M.min^{-1}}$.
Tangente en $t=10$ : pente $\approx \frac{4,4 - 2,0}{20 - 5} = \SI{1,6 \times 10^{-4}}{M.min^{-1}}$ (valeurs approchées graphiques).
\end{enumerate}
\end{corrige}

\begin{exercice}[Exercice 3 : Catalyse et énergie d'activation]
On compare deux réactions identiques : l'une sans catalyseur ($E_a = \SI{75}{kJ.mol^{-1}}$), l'autre avec catalyseur ($E_a' = \SI{50}{kJ.mol^{-1}}$). À $T = \SI{300}{K}$, la vitesse sans catalyseur est $v_0$.
En utilisant la relation d'Arrhenius simplifiée ($k \propto e^{-E_a/RT}$), estimer le rapport $\frac{v_{\text{cat}}}{v_0}$ à \SI{300}{K}. ($R = \SI{8,314}{J.mol^{-1}.K^{-1}}$).
\end{exercice}

\begin{corrige}[Exercice 3]
\[ \frac{v_{\text{cat}}}{v_0} = \frac{k_{\text{cat}}}{k_0} = \exp\left( \frac{E_a - E_a'}{RT} \right) = \exp\left( \frac{(75 - 50) \times 10^3}{8,314 \times 300} \right) \]
\[ = \exp\left( \frac{25000}{2494{,}2} \right) = \exp(10,02) \approx \num{2,2e4} \]
Le catalyseur accélère la réaction d'un facteur \textbf{environ 22 000 fois} à température ambiante. Cela illustre la puissance de la catalyse.
\end{corrige}

\newpage

\coursec{Méthodes et exercices résolus}

\coursec{Cinétique chimique}

\courssub{Réactions rapides, lentes et très lentes}

\begin{definition}[Vitesse d'une réaction chimique]
La \cle{vitesse d'une réaction chimique} caractérise la rapidité avec laquelle un système chimique évolue vers son état final. Elle se mesure par la variation de concentration d'une espèce par unité de temps.
\end{definition}

\begin{aretenir}[Classification des réactions suivant leur vitesse]
\begin{itemize}
\item \textbf{Réactions rapides (quasi instantanées) :} vitesse de l'ordre de $10^{-3}$ à $10^{-6}\,\si{s}$.
      \begin{itemize}
      \item Réactions de précipitation : \ce{Ag+(aq) + Cl-(aq) -> AgCl(s)}.
      \item Réactions acide-base (neutralisation) : \ce{H3O+(aq) + HO-(aq) -> 2 H2O(l)}.
      \item Réactions d'oxydo-réduction en solution (ex. : \ce{MnO4-}/\ce{Fe^2+}).
      \end{itemize}
\item \textbf{Réactions lentes :} vitesse mesurable en minutes, heures ou jours. Elles permettent un suivi cinétique.
      \begin{itemize}
      \item \ce{2 I- + S2O8^2- -> I2 + 2 SO4^2-} (réaction de référence en TP).
      \item Estérification : \ce{CH3COOH + C2H5OH <=> CH3COOC2H5 + H2O}.
      \item Formation de la rouille (corrosion du fer) : \ce{4 Fe + 3 O2 + 6 H2O -> 4 Fe(OH)3}.
      \end{itemize}
\item \textbf{Réactions très lentes :} vitesse négligeable à température ambiante sans catalyseur.
      \begin{itemize}
      \item Synthèse de l'eau : \ce{2 H2 + O2 -> 2 H2O} (nécessite étincelle ou platine).
      \item Combustion du diamant : \ce{C(diamant) + O2 -> CO2}.
      \end{itemize}
\end{itemize}
\end{aretenir}

\begin{savaistu}[Cinétique au Cameroun : du terrain au laboratoire]
\begin{itemize}
\item \textbf{Bili-bili et vin de palme :} La fermentation des sucres (mil, maïs, sève de palmier) en éthanol et \ce{CO2} est une \cle{catalyse enzymatique} (levures). La température ambiante tropicale ($\approx \SI{28}{\celsius}$) accélère naturellement la cinétique.
\item \textbf{Huile de palme et savonneries artisanales :} La saponification (\cle{hydrolyse basique} des triglycérides) est lente à froid ; le chauffage (facteur cinétique température) et l'excès de soude (facteur concentration) l'accélèrent.
\item \textbf{SONARA (Limbe) :} Le craquage catalytique du pétrole utilise des zéolites (catalyse hétérogène) pour briser les grosses molécules d'hydrocarbures à haute température.
\item \textbf{CIMENCAM :} La cuisson du clinker (\ce{CaCO3 -> CaO + CO2}) est une réaction solide-gaz lente, contrôlée par la température et la surface de contact (broyage fin).
\item \textbf{Pharmacopée traditionnelle :} La décoction (ébullition prolongée) extrait les principes actifs en augmentant la température et la surface de contact (plantes broyées), deux facteurs cinétiques majeurs.
\end{itemize}
\end{savaistu}

\courssub{Vitesses de formation et de disparition}

\begin{methode}[Construire la courbe [I2]=f(t) par dosage des prélèvements (trempe)]
\begin{enumerate}
\item \textbf{Identifier la réaction de formation du diiode} (couple lent de référence) :
      \[ \ce{S2O8^2- + 2 I- -> 2 SO4^2- + I2} \]
      \textbf{Vérification des charges :} $-2 + 2(-1) = -4$ à gauche ; $2(-2) + 0 = -4$ à droite. \textbf{Équilibrée.}
\item \textbf{Protocole de trempe :} À chaque instant $t_i$, prélever un volume $V_p$ du mélange réactionnel et le verser immédiatement dans un erlenmeyer contenant de l'eau glacée (volume important). La \cle{trempe} abaisse brutalement la température et dilue les réactifs : la réaction est \emph{bloquée}.
\item \textbf{Dosage du diiode formé} par titrage au thiosulfate de sodium :
      \[ \ce{I2 + 2 S2O3^2- -> 2 I- + S4O6^2-} \]
      \textbf{Vérification :} I : $2 \to 2$ ; S : $2 \to 4$ ; O : $6 \to 6$ ; Charges : $0 + 2(-2) = -4 \to 2(-1) + (-2) = -4$. \textbf{Équilibrée.}
      On en déduit $n(\ce{I2})_i = \dfrac{C_{\ce{S2O3^2-}} \times V_{\ce{S2O3^2-}}}{2}$, puis $[\ce{I2}]_i = \dfrac{n(\ce{I2})_i}{V_{\text{total}}}$.
\item \textbf{Reporter les points} $(t_i\,;\,[\ce{I2}]_i)$ sur un graphe : axes gradués régulièrement, grandeurs et unités indiquées, échelle utilisant au moins la moitié de la feuille.
\item \textbf{Tracer une courbe lisse} passant au plus près des points (jamais de segments de droite !) et partant de l'origine si $[\ce{I2}]_0 = 0$.
\item \textbf{Interpréter l'allure} : courbe croissante, pente décroissante (réaction qui ralentit car les réactifs s'épuisent), tendant vers une asymptote horizontale $[\ce{I2}]_{\infty} = \dfrac{x_{\max}}{V}$.
\end{enumerate}
\end{methode}

\begin{exoresolu}[Courbe de formation du diiode]
On mélange $V_1 = \SI{50,0}{mL}$ d'une solution d'iodure de potassium $C_1 = \SI{0,50}{mol.L^{-1}}$ avec $V_2 = \SI{50,0}{mL}$ d'une solution de peroxodisulfate de potassium $C_2 = \SI{0,10}{mol.L^{-1}}$. Le dosage du diiode formé donne :

\begin{center}
\begin{tabularx}{\linewidth}{|l|X|X|X|X|X|X|}
\hline
\rowcolor{popblueL} $t$ (min) & 0 & 5 & 10 & 20 & 30 & 40 \\
\hline $[\ce{I2}]$ (\SI{e-2}{mol.L^{-1}}) & 0 & 2,1 & 3,4 & 4,6 & 5,1 & 5,3 \\
\hline
\end{tabularx}
\end{center}

Tracer la courbe $[\ce{I2}]=f(t)$ et déterminer la concentration finale théorique du diiode.
\tcblower
\textbf{1. Vérification de la cohérence des données.}
Quantités initiales introduites :
\begin{align*}
n_0(\ce{I-}) &= C_1 V_1 = 0,50 \times 50,0\times 10^{-3} = \SI{2,5e-2}{mol}\\
n_0(\ce{S2O8^2-}) &= C_2 V_2 = 0,10 \times 50,0\times 10^{-3} = \SI{5,0e-3}{mol}
\end{align*}
D'après l'équation $\ce{S2O8^2- + 2 I- -> 2 SO4^2- + I2}$, il faut 2 mol de \ce{I-} pour 1 mol de \ce{S2O8^2-}.
$\dfrac{n_0(\ce{I-})}{2} = \SI{1,25e-2}{mol} > n_0(\ce{S2O8^2-}) = \SI{5,0e-3}{mol}$ : le peroxodisulfate est le \cle{réactif limitant}, donc $x_{\max} = \SI{5,0e-3}{mol}$.

\textbf{2. Concentration finale théorique.}
Volume total $V = V_1 + V_2 = \SI{100,0}{mL} = \SI{0,1000}{L}$ :
\[ [\ce{I2}]_{\infty} = \frac{x_{\max}}{V} = \frac{5,0\times 10^{-3}}{0,1000} = \SI{5,0e-2}{mol.L^{-1}} \]
La valeur expérimentale limite ($5,3\times 10^{-2}$ \si{mol.L^{-1}} à $t = \SI{40}{min}$) est proche : cohérent.

\textbf{3. Tracé.}

\begin{popfigure}
\begin{tikzpicture}
\begin{axis}[width=0.9\linewidth,height=6.5cm,
xlabel={$t$ (min)}, ylabel={$[\ce{I2}]\ (\times 10^{-2}\ \si{mol.L^{-1}})$},
xmin=0, xmax=45, ymin=0, ymax=6,
grid=both, grid style={popline!40}, axis lines=left,
label style={font=\small}, tick label style={font=\small}]
% Courbe théorique lissée (modèle exponentiel)
\addplot[popcurve, domain=0:45, samples=200] {5.3*(1-exp(-x/9))};
% Asymptote théorique
\addplot[dashed, poporange, thick, domain=0:45] {5.0};
% Points expérimentaux
\addplot[only marks, mark=*, popdark, mark size=2.5] coordinates {
(0,0) (5,2.1) (10,3.4) (20,4.6) (30,5.1) (40,5.3)
};
\end{axis}
\end{tikzpicture}
\legende{Courbe $[\ce{I2}]=f(t)$ : croissance rapide puis ralentissement vers l'asymptote $[\ce{I2}]_{\infty} = \SI{5,0e-2}{mol.L^{-1}}$ (trait orange pointillé).}
\end{popfigure}

\textbf{Conclusion :} La courbe est croissante, de pente décroissante, et tend vers $[\ce{I2}]_{\infty} = \SI{5,0e-2}{mol.L^{-1}}$, valeur fixée par le réactif limitant \ce{S2O8^2-}.
\end{exoresolu}

\begin{methode}[Déterminer vitesse moyenne et vitesse instantanée à partir d'une courbe $C=f(t)$]
\begin{enumerate}
\item \textbf{Vitesse moyenne de formation} d'un produit $P$ entre $t_1$ et $t_2$ :
      \[ v_{m}(P) = \frac{[P]_2 - [P]_1}{t_2 - t_1} = \frac{\Delta [P]}{\Delta t} \]
      C'est la pente de la \cle{sécante} reliant les deux points de la courbe.
\item \textbf{Vitesse moyenne de disparition} d'un réactif $R$ (valeur positive) :
      \[ v_{m}(R) = -\frac{[R]_2 - [R]_1}{t_2 - t_1} = \frac{[R]_1 - [R]_2}{t_2 - t_1} \]
\item \textbf{Vitesse instantanée} à la date $t$ : pente de la \cle{tangente} à la courbe en ce point.
      \[ v(P,t) = \left(\frac{\mathrm{d}[P]}{\mathrm{d}t}\right)_t \qquad v(R,t) = -\left(\frac{\mathrm{d}[R]}{\mathrm{d}t}\right)_t \]
\item \textbf{Protocole graphique pour la tangente :}
      \begin{enumerate}
      \item Tracer la tangente à la courbe au point d'abscisse $t$ (règle, bien prolongée de part et d'autre).
      \item Choisir deux points \textbf{éloignés} $A$ et $B$ \emph{sur la tangente} (pas sur la courbe !).
      \item Lire leurs coordonnées $(t_A; C_A)$ et $(t_B; C_B)$.
      \item Calculer $v = \dfrac{C_B - C_A}{t_B - t_A}$.
      \end{enumerate}
\item \textbf{Unités et chiffres significatifs :} \si{mol.L^{-1}.min^{-1}} ou \si{mol.L^{-1}.s^{-1}}, 2 chiffres significatifs (lecture graphique).
\item \textbf{Réflexe Bac :} La vitesse instantanée \emph{diminue} au cours du temps car les concentrations des réactifs diminuent (facteur cinétique concentration).
\end{enumerate}
\end{methode}

\begin{exoresolu}[Vitesses de formation du diiode]
On reprend la courbe $[\ce{I2}]=f(t)$ de l'exercice précédent. La tangente à l'origine passe par le point $(t = \SI{10}{min}\,;\,[\ce{I2}] = \SI{4,0e-2}{mol.L^{-1}})$ et la tangente à $t = \SI{20}{min}$ passe par les points $(\SI{10}{min}\,;\,\SI{4,0e-2}{mol.L^{-1}})$ et $(\SI{30}{min}\,;\,\SI{5,2e-2}{mol.L^{-1}})$.
\begin{enumerate}
\item Calculer la vitesse moyenne de formation du diiode entre $t = \SI{5}{min}$ et $t = \SI{20}{min}$.
\item Calculer les vitesses instantanées à $t = 0$ et à $t = \SI{20}{min}$. Conclure.
\end{enumerate}
\tcblower
\textbf{1. Vitesse moyenne entre 5 min et 20 min.}
Lecture sur le tableau : $[\ce{I2}]_5 = \SI{2,1e-2}{mol.L^{-1}}$ et $[\ce{I2}]_{20} = \SI{4,6e-2}{mol.L^{-1}}$.
\[ v_m = \frac{[\ce{I2}]_{20}-[\ce{I2}]_5}{20-5} = \frac{(4,6-2,1)\times 10^{-2}}{15} = \SI{1,7e-3}{mol.L^{-1}.min^{-1}} \]

\textbf{2. Vitesses instantanées.}
À $t = 0$, la tangente passe par l'origine et par $(\SI{10}{min}\,;\,\SI{4,0e-2}{mol.L^{-1}})$ :
\[ v(0) = \frac{4,0\times 10^{-2} - 0}{10 - 0} = \SI{4,0e-3}{mol.L^{-1}.min^{-1}} \]

À $t = \SI{20}{min}$, on utilise les deux points de la tangente :
\[ v(20) = \frac{(5,2-4,0)\times 10^{-2}}{30-10} = \frac{1,2\times 10^{-2}}{20} = \SI{6,0e-4}{mol.L^{-1}.min^{-1}} \]

\textbf{3. Conclusion.}
On constate que $v(20) < v_m < v(0)$ : la vitesse de formation du diiode \textbf{diminue au cours du temps}. Cela s'explique par la diminution des concentrations des réactifs \ce{I-} et \ce{S2O8^2-}, facteur cinétique essentiel : moins il reste de réactifs, plus la réaction est lente.
\end{exoresolu}

\courssub{Vitesse volumique, avancement et temps de demi-réaction}

\begin{definition}[Vitesse volumique de réaction]
Pour une réaction se déroulant dans un volume $V$ constant, la \cle{vitesse volumique de réaction} $v$ est définie par :
\[ v = \frac{1}{V}\,\frac{\mathrm{d}x}{\mathrm{d}t} \]
où $x$ est l'\cle{avancement} (en mol) et $V$ le volume du milieu réactionnel (en L).
Unité : \si{mol.L^{-1}.s^{-1}} (ou \si{mol.L^{-1}.min^{-1}}).
\end{definition}

\begin{propriete}[Lien entre $v$ et variations de concentrations]
Pour une équation bilan $\sum \nu_i \ce{A_i} = 0$ (coefficients $\nu_i$ algébriques : $>0$ pour produits, $<0$ pour réactifs) :
\[ v = \frac{1}{\nu_P}\frac{\mathrm{d}[P]}{\mathrm{d}t} = -\frac{1}{\nu_R}\frac{\mathrm{d}[R]}{\mathrm{d}t} \]
\textbf{Attention aux coefficients stœchiométriques !}
\end{propriete}

\begin{methode}[Exploiter $v = \frac{1}{V}\frac{\mathrm{d}x}{\mathrm{d}t}$ et déterminer $t_{1/2}$]
\begin{enumerate}
\item \textbf{Relier $x$ aux concentrations.} Pour un produit $P$ : $[P] = \dfrac{\nu_P\, x}{V} \Rightarrow \dfrac{\mathrm{d}[P]}{\mathrm{d}t} = \dfrac{\nu_P}{V}\dfrac{\mathrm{d}x}{\mathrm{d}t} = \nu_P\, v$.
\item \textbf{Graphiquement :} $\dfrac{\mathrm{d}x}{\mathrm{d}t}$ est la pente de la tangente à la courbe $x=f(t)$ ; on divise ensuite par $V$.
\item \textbf{Temps de demi-réaction $t_{1/2}$ :} durée au bout de laquelle l'avancement atteint la moitié de sa valeur \emph{finale} :
      \[ x(t_{1/2}) = \frac{x_{f}}{2} \]
      ($x_f = x_{\max}$ si réaction totale ; $x_f < x_{\max}$ si réaction limitée).
\item \textbf{Lecture graphique :} calculer $x_f/2$, repérer sur l'axe des ordonnées, projeter sur la courbe $x=f(t)$ puis sur l'axe des temps.
\item \textbf{Vérifier la cohérence :} $t_{1/2} < t_{\text{fin}}$ et $v(t_{1/2}) < v(0)$.
\end{enumerate}
\end{methode}

\begin{exoresolu}[Décomposition de l'eau oxygénée catalysée]
On étudie la décomposition de l'eau oxygénée en présence d'un catalyseur :
\[ \ce{2 H2O2 -> 2 H2O + O2} \]
Le volume de la solution est $V = \SI{200}{mL}$. Le suivi du dégagement de dioxygène donne la quantité de \ce{H2O2} disparue, d'où la courbe $x = f(t)$. On lit : $x_{\max} = \SI{4,0e-3}{mol}$ ; la tangente à l'origine coupe la droite $x = x_{\max}$ en $t = \SI{8,0}{min}$ ; la courbe atteint $x = \SI{2,0e-3}{mol}$ à $t = \SI{5,5}{min}$.
\begin{enumerate}
\item Déterminer la vitesse volumique initiale de la réaction.
\item Déterminer le temps de demi-réaction.
\item En déduire la concentration initiale en eau oxygénée.
\end{enumerate}
\tcblower
\textbf{1. Vitesse volumique initiale.}
Pente de la tangente à l'origine de $x=f(t)$ (triangle rectangle) :
\[ \left(\frac{\mathrm{d}x}{\mathrm{d}t}\right)_0 = \frac{x_{\max}}{8,0} = \frac{4,0\times 10^{-3}}{8,0} = \SI{5,0e-4}{mol.min^{-1}} \]
D'où, avec $V = \SI{0,200}{L}$ :
\[ v_0 = \frac{1}{V}\left(\frac{\mathrm{d}x}{\mathrm{d}t}\right)_0 = \frac{5,0\times 10^{-4}}{0,200} = \SI{2,5e-3}{mol.L^{-1}.min^{-1}} \]

\textbf{2. Temps de demi-réaction.}
La réaction est totale ($x_f = x_{\max}$), donc :
\[ x(t_{1/2}) = \frac{x_{\max}}{2} = \frac{4,0\times 10^{-3}}{2} = \SI{2,0e-3}{mol} \]
D'après l'énoncé, cette valeur est atteinte à $t = \SI{5,5}{min}$. Donc $\boxed{t_{1/2} = \SI{5,5}{min}}$.

\textbf{3. Concentration initiale en \ce{H2O2}.}
Quand la réaction est terminée, tout le réactif a disparu. Coefficient stœchiométrique 2 devant \ce{H2O2} :
\[ n_0(\ce{H2O2}) = 2\,x_{\max} = 2 \times 4,0\times 10^{-3} = \SI{8,0e-3}{mol} \]
\[ [\ce{H2O2}]_0 = \frac{n_0}{V} = \frac{8,0\times 10^{-3}}{0,200} = \SI{4,0e-2}{mol.L^{-1}} \]

\textbf{Conclusion :} $v_0 = \SI{2,5e-3}{mol.L^{-1}.min^{-1}}$, $t_{1/2} = \SI{5,5}{min}$ et $[\ce{H2O2}]_0 = \SI{4,0e-2}{mol.L^{-1}}$.
\end{exoresolu}

\courssub{Suivi temporel d'une réaction}

\begin{methode}[Choisir et mettre en œuvre une méthode de suivi]
\begin{enumerate}
\item \textbf{Dosage des prélèvements (méthode par trempe) :}
      \begin{itemize}
      \item Prélèvement d'un aliquote à l'instant $t$.
      \item \cle{Trempe} : dilution dans eau glacée (baisse $T$, dilution $\Rightarrow$ réaction bloquée).
      \item Dosage titrimétrique de l'espèce suivie (ex. : acide restant, diiode formé).
      \item \emph{Inconvénient :} discontinu, perturbant (prélèvement modifie $V$ si non compensé).
      \end{itemize}
\item \textbf{Méthodes physiques (suivi continu, non perturbant) :}
      \begin{itemize}
      \item \textbf{Colorimétrie / Spectrophotométrie :} si une espèce est colorée (\ce{I2} brun, \ce{MnO4-} violet, \ce{Fe(SCN)^2+} rouge). Loi de Beer-Lambert : $A = \varepsilon\, l\, C$. On suit $A(t) \propto C(t)$.
      \item \textbf{Conductimétrie :} si la conductivité $\sigma$ varie (disparition/apparition d'ions). $\sigma = \sum \lambda_i C_i$. Rapide, précis, applicable aux réactions ioniques.
      \item \textbf{Manométrie (pression) :} si réaction gazeuse avec changement de nombre de molécules gazeuses ($\Delta n_g \neq 0$). On suit $P(t)$ dans un volume fixe.
      \item \textbf{pH-métrie :} si $\ce{H3O+}$ ou $\ce{HO-}$ apparaît/disparaît.
      \end{itemize}
\item \textbf{Critères de choix :} espèce suivie (colorée ? ionique ? gazeuse ?), rapidité de la réaction, matériel disponible.
\end{enumerate}
\end{methode}

\begin{exemplebox}[Exemples de suivi adaptés]
\begin{itemize}
\item \ce{S2O8^2- + 2 I- -> I2 + 2 SO4^2-} : \textbf{Colorimétrie} (apparition \ce{I2} brun) ou \textbf{Dosage par trempe} (thiosulfate).
\item \ce{CH3COOH + C2H5OH <=> CH3COOC2H5 + H2O} : \textbf{Dosage par trempe} (acide restant par \ce{NaOH}) car aucune espèce n'est colorée ni gazeuse, conductivité peu variable.
\item \ce{2 H2O2 -> 2 H2O + O2} : \textbf{Manométrie} (apparition \ce{O2} gazeux) ou \textbf{Dosage par trempe} (\ce{H2O2} restant par \ce{KMnO4}).
\item \ce{H3O+ + HO- -> 2 H2O} : \textbf{Conductimétrie} (disparition ions forts) ou \textbf{pH-métrie}.
\end{itemize}
\end{exemplebox}

\courssub{Tracer les courbes de formation des produits et de disparition des réactifs}

\begin{methode}[Courbes des réactifs et des produits sur un même graphe]
\begin{enumerate}
\item \textbf{Écrire le tableau d'avancement} (concentrations) pour relier toutes les quantités à $\xi(t) = x(t)/V$.
\item \textbf{Exprimer chaque concentration :}
      \begin{align*}
      \text{Réactif } R\ (\nu_R > 0) &: [R](t) = [R]_0 - \nu_R\,\xi(t) \\
      \text{Produit } P\ (\nu_P > 0) &: [P](t) = \nu_P\,\xi(t)
      \end{align*}
\item \textbf{Identifier les valeurs limites :}
      \begin{itemize}
      \item Courbes des réactifs : partent de $[R]_0$, \emph{décroissent}, tendent vers $[R]_f$ (0 si limitant, $>0$ si en excès).
      \item Courbes des produits : partent de 0, \emph{croissent}, tendent vers $[P]_f = \nu_P\,\xi_f$.
      \item Toutes tendent vers des asymptotes horizontales.
      \end{itemize}
\item \textbf{Respecter la stœchiométrie sur le graphe :} si $\nu_R = 2\,\nu_P$, le réactif disparaît deux fois plus vite que le produit n'apparaît (pente initiale double en valeur absolue).
\item \textbf{Point remarquable :} à $t_{1/2}$, toutes les courbes ont parcouru la moitié de leur variation totale.
\end{enumerate}
\end{methode}

\begin{exoresolu}[Courbes de la réaction \ce{I-}/\ce{S2O8^2-}]
On reprend le mélange de l'exercice 1 : $[\ce{I-}]_0 = \SI{0,25}{mol.L^{-1}}$, $[\ce{S2O8^2-}]_0 = \SI{5,0e-2}{mol.L^{-1}}$ (concentrations après mélange), $V = \SI{100}{mL}$, réaction totale d'avancement maximal $x_{\max} = \SI{5,0e-3}{mol}$.
\begin{enumerate}
\item Exprimer $[\ce{I-}](t)$, $[\ce{S2O8^2-}](t)$ et $[\ce{I2}](t)$ en fonction de $x(t)$.
\item Calculer les concentrations finales et esquisser les trois courbes.
\end{enumerate}
\tcblower
\textbf{1. Tableau d'avancement} (concentrations en \si{mol.L^{-1}}, $\xi = x/V$) :

\begin{center}
\begin{tabularx}{\linewidth}{|l|X|X|X|X|}
\hline
\rowcolor{popblueL} État & \ce{S2O8^2-} & \ce{I-} & \ce{SO4^2-} & \ce{I2} \\
\hline Initial ($t=0$) & $5,0\times 10^{-2}$ & $2,5\times 10^{-1}$ & 0 & 0 \\
\hline En cours ($t$) & $5,0\times 10^{-2}-\xi$ & $2,5\times 10^{-1}-2\xi$ & $2\xi$ & $\xi$ \\
\hline Final & $0$ & $1,5\times 10^{-1}$ & $1,0\times 10^{-1}$ & $5,0\times 10^{-2}$ \\
\hline
\end{tabularx}
\end{center}

avec $\xi(t) = \dfrac{x(t)}{V}$ et $\xi_{\max} = \dfrac{x_{\max}}{V} = \SI{5,0e-2}{mol.L^{-1}}$.

\textbf{2. Concentrations finales.}
Le réactif limitant est \ce{S2O8^2-} :
\begin{align*}
[\ce{S2O8^2-}]_f &= 0 \\
[\ce{I-}]_f &= 0,25 - 2\times 5,0\times 10^{-2} = \SI{0,15}{mol.L^{-1}} \\
[\ce{I2}]_f &= \SI{5,0e-2}{mol.L^{-1}}
\end{align*}

\textbf{3. Esquisse des courbes.}

\begin{popfigure}
\begin{tikzpicture}
\begin{axis}[width=0.9\linewidth,height=7cm,
xlabel={$t$ (min)}, ylabel={Concentration (\si{mol.L^{-1}})},
xmin=0, xmax=45, ymin=0, ymax=0.28,
grid=both, grid style={popline!40}, axis lines=left,
label style={font=\small}, tick label style={font=\small},
legend style={font=\small, at={(0.98,0.98)}, anchor=north east, draw=popline}]
\addplot[popcurve, thick, domain=0:45, samples=200] {0.05*(1-exp(-x/9))};
\addlegendentry{$[\ce{I2}]$}
\addplot[poporange, thick, domain=0:45, samples=200] {0.05*exp(-x/9)};
\addlegendentry{$[\ce{S2O8^2-}]$}
\addplot[poppurple, thick, domain=0:45, samples=200] {0.25 - 0.10*(1-exp(-x/9))};
\addlegendentry{$[\ce{I-}]$}
% Asymptotes
\addplot[dashed, popmuted, domain=0:45] {0.05};
\addplot[dashed, popmuted, domain=0:45] {0.15};
\end{axis}
\end{tikzpicture}
\legende{Évolution des concentrations : \ce{S2O8^2-} s'annule (limitant), \ce{I-} décroît de 0,25 à 0,15 (variation double), \ce{I2} croît vers 0,05. Asymptotes en pointillés.}
\end{popfigure}

\textbf{Conclusion :} $[\ce{S2O8^2-}]$ s'annule, $[\ce{I-}]$ décroît de $\SI{0,25}{}$ à $\SI{0,15}{mol.L^{-1}}$ (variation double de celle de \ce{S2O8^2-}, coefficient 2), et $[\ce{I2}]$ croît vers $\SI{5,0e-2}{mol.L^{-1}}$.
\end{exoresolu}

\courssub{Facteurs cinétiques et catalyse}

\begin{aretenir}[Facteurs influençant la vitesse d'une réaction]
\begin{enumerate}
\item \textbf{Concentration des réactifs :} Plus la concentration est élevée, plus les chocs efficaces sont fréquents $\Rightarrow$ vitesse plus grande. (Loi de vitesse formelle vue en Spé).
\item \textbf{Température :} L'élévation de $T$ augmente l'énergie cinétique des molécules $\Rightarrow$ plus de chocs efficaces (énergie $> E_a$) $\Rightarrow$ vitesse plus grande. Règle empirique : $v$ double pour $\Delta T \approx \SI{10}{\celsius}$.
\item \textbf{Surface de contact (réactions hétérogènes) :} Pour un solide (ex. : \ce{CaCO3}, \ce{Fe}, catalyseur Pt), plus la surface est grande (poudre vs morceaux), plus la vitesse est grande.
\item \textbf{Catalyseur :} Espèce qui accélère la réaction \textbf{sans être consommée} et \textbf{sans modifier l'état final} (même $x_f$, même asymptote).
\end{enumerate}
\end{aretenir}

\begin{definition}[Types de catalyse]
\begin{itemize}
\item \textbf{Catalyse homogène :} Catalyseur et réactifs dans la \emph{même phase} (généralement solution aqueuse).
      Ex. : Ions \ce{Fe^3+} pour \ce{2 I- + S2O8^2- -> I2 + 2 SO4^2-}.
\item \textbf{Catalyse hétérogène :} Catalyseur et réactifs dans des \emph{phases différentes}.
      Ex. : Platine solide (Pt) pour la synthèse de l'eau \ce{2 H2(g) + O2(g) -> 2 H2O(l/g)}.
\item \textbf{Catalyse enzymatique :} Catalyse par une enzyme (protéine) en milieu aqueux, très spécifique, très efficace (ex. : digestion, fermentation bili-bili).
\end{itemize}
\end{definition}

\begin{methode}[Mettre en évidence la catalyse expérimentalement]
\begin{enumerate}
\item \textbf{Réaliser deux expériences identiques} (mêmes réactifs, volumes, concentrations, température), l'une \emph{sans} catalyseur (témoin), l'autre \emph{avec} catalyseur.
\item \textbf{Observer et comparer} : vitesse d'apparition d'une coloration, dégagement gazeux, temps de demi-réaction $t_{1/2}$.
\item \textbf{Conclure :} Si la réaction est plus rapide avec le catalyseur mais que l'état final (asymptote, teinte finale, volume gazeux final) est identique $\Rightarrow$ c'est un catalyseur.
\item \textbf{Interpréter :} Le catalyseur propose un « chemin réactionnel » différent (mécanisme à étapes élémentaires rapides) remplaçant l'étape lente de la réaction non catalysée.
\end{enumerate}
\end{methode}

\begin{exoresolu}[Catalyse de la réaction \ce{I-}/\ce{S2O8^2-} par les ions \ce{Fe^3+}]
Dans trois béchers identiques, on mélange \SI{20}{mL} de solution d'iodure de potassium \SI{0,20}{mol.L^{-1}} et \SI{20}{mL} de solution de peroxodisulfate de potassium \SI{0,05}{mol.L^{-1}}.
\begin{itemize}
\item Bécher 1 : rien n'est ajouté (témoin pur).
\item Bécher 2 : on ajoute 5 gouttes de solution de chlorure de fer(III) \ce{FeCl3}.
\item Bécher 3 : on ajoute 5 gouttes d'eau distillée (témoin dilution).
\end{itemize}
La coloration brune apparaît nettement au bout de \SI{3}{min} dans le bécher 2, de \SI{15}{min} dans les béchers 1 et 3. Au bout d'une heure, les trois béchers ont la même teinte brune finale.
\begin{enumerate}
\item Quel est le rôle des ions \ce{Fe^3+} ? Justifier.
\item Pourquoi la teinte finale est-elle la même dans les trois béchers ?
\item De quel type de catalyse s'agit-il ?
\end{enumerate}
\tcblower
\textbf{1. Rôle des ions \ce{Fe^3+}.}
Les béchers 1 et 3 évoluent à la même vitesse : l'ajout d'eau distillée ne change pas la cinétique (dilution négligeable). Le bécher 2, seul à contenir les ions \ce{Fe^3+}, brunit 5 fois plus vite ($15/3 = 5$). Les ions \ce{Fe^3+} \textbf{accélèrent la réaction} : ils jouent le rôle de \cle{catalyseur}.
Le mécanisme proposé fait intervenir deux étapes rapides :
\begin{align*}
\ce{2 Fe^3+ + 2 I- &-> 2 Fe^2+ + I2} \quad \text{(étape 1)} \\
\ce{2 Fe^2+ + S2O8^2- &-> 2 Fe^3+ + 2 SO4^2-} \quad \text{(étape 2)}
\end{align*}
Bilan : \ce{2 I- + S2O8^2- -> I2 + 2 SO4^2-}. Les ions \ce{Fe^3+} consommés à l'étape 1 sont \textbf{régénérés} à l'étape 2 : ils ne figurent pas dans le bilan global, confirmant qu'un catalyseur n'est pas consommé.

\textbf{2. Teinte finale identique.}
La teinte brune est due au diiode \ce{I2}. La quantité finale de diiode ne dépend que des quantités initiales de réactifs et de l'avancement maximal $x_{\max}$, identiques dans les trois béchers. Le catalyseur \textbf{modifie la vitesse mais pas l'état final} : il ne déplace pas l'équilibre ni la limite de la réaction, il permet seulement de l'atteindre plus vite.

\textbf{3. Type de catalyse.}
Les ions \ce{Fe^3+} sont en solution aqueuse, comme les réactifs \ce{I-} et \ce{S2O8^2-} : catalyseur et réactifs appartiennent à la \textbf{même phase}. Il s'agit donc d'une \cle{catalyse homogène}.

\textbf{Conclusion :} Les ions \ce{Fe^3+} catalysent de manière homogène la réaction entre \ce{I-} et \ce{S2O8^2-} : la réaction est environ 5 fois plus rapide, mais la quantité finale de diiode est inchangée.
\end{exoresolu}

\courssub{Étude cinétique de l'estérification}

\begin{methode}[Étude cinétique de l'estérification : acide éthanoïque + éthanol]
\begin{enumerate}
\item \textbf{Écrire l'équation} de la réaction, lente et limitée :
      \[ \ce{CH3COOH + C2H5OH <=> CH3COOC2H5 + H2O} \]
      \chemfig{CH_3-C(=[1]O)-OH} + \chemfig{CH_3-CH_2-OH} \ce{<=>} \chemfig{CH_3-C(=[1]O)-O-CH_2-CH_3} + \ce{H2O}
\item \textbf{Préparer un mélange équimolaire} (ex. $n_0 = \SI{0,10}{mol}$ de chaque réactif), réparti en plusieurs tubes scellés placés dans un thermostat à température contrôlée.
\item \textbf{Suivre l'avancement} : à des dates $t_i$, sortir un tube, le \cle{tremper} (glace) puis doser l'acide restant par la soude : $n_{\text{acide}}(t) = n_0 - x(t)$, d'où $x(t)$ et $n_{\text{ester}}(t) = x(t)$.
\item \textbf{Tracer $n_{\text{ester}} = f(t)$ (ou $x=f(t)$) :} courbe croissante tendant vers une \emph{asymptote} $x_f < x_{\max}$ car la réaction est \textbf{limitée} par l'hydrolyse inverse (équilibre chimique).
\item \textbf{Comparer l'effet des facteurs cinétiques} sur des courbes superposées :
      \begin{itemize}
      \item \emph{Température plus élevée} ou \emph{catalyseur} (\ce{H2SO4}, ions \ce{H3O+}) : l'asymptote est atteinte \textbf{plus rapidement} mais sa \textbf{valeur est inchangée} ($x_f$ constant).
      \item \emph{Excès d'un réactif} : l'asymptote est \textbf{plus haute} (déplacement de l'équilibre vers les produits) et atteinte plus vite.
      \end{itemize}
\item \textbf{Exploiter :} déterminer $t_{1/2} = t$ tel que $x = x_f/2$ sur chaque courbe et comparer.
\end{enumerate}
\end{methode}

\begin{exoresolu}[Estérification de l'acide éthanoïque par l'éthanol]
On réalise à \SI{60}{\celsius} un mélange équimolaire de $n_0 = \SI{0,20}{mol}$ d'acide éthanoïque et $n_0 = \SI{0,20}{mol}$ d'éthanol, réparti en tubes. Le dosage de l'acide restant donne l'avancement $x(t)$ :

\begin{center}
\begin{tabularx}{\linewidth}{|l|X|X|X|X|X|X|}
\hline
\rowcolor{popblueL} $t$ (h) & 0 & 2 & 5 & 10 & 20 & 40 \\
\hline $x$ (\SI{e-2}{mol}) — sans catalyseur & 0 & 2,2 & 4,6 & 7,6 & 10,8 & 12,6 \\
\hline $x$ (\SI{e-2}{mol}) — avec \ce{H2SO4} & 0 & 6,8 & 10,9 & 12,7 & 13,2 & 13,3 \\
\hline
\end{tabularx}
\end{center}

L'avancement final vaut $x_f = \SI{0,133}{mol}$ dans les deux cas.
\begin{enumerate}
\item Pourquoi $x_f$ est-il inférieur à $x_{\max} = \SI{0,20}{mol}$ ?
\item Déterminer le temps de demi-réaction dans chaque cas.
\item Conclure sur le rôle de l'acide sulfurique.
\end{enumerate}
\tcblower
\textbf{1. Avancement final limité.}
L'estérification est une réaction \textbf{lente et limitée} : simultanément, l'ester formé réagit avec l'eau (hydrolyse, réaction inverse). Il s'établit un \cle{équilibre chimique} :
\[ \ce{CH3COOH + C2H5OH <=> CH3COOC2H5 + H2O} \]
Pour un mélange équimolaire acide/alcool primaire, le rendement limite est d'environ $67\,\%$ : $x_f \approx \frac{2}{3}x_{\max} = \SI{0,133}{mol} < x_{\max} = \SI{0,20}{mol}$. La réaction s'arrête d'évoluer \emph{macroscopiquement} avant d'avoir consommé tout le réactif.

\textbf{2. Temps de demi-réaction.}
Par définition : $x(t_{1/2}) = \dfrac{x_f}{2} = \dfrac{0,133}{2} \approx \SI{6,65e-2}{mol}$.

\emph{Sans catalyseur :} $x = \SI{6,65e-2}{mol}$ est compris entre les valeurs à $t = \SI{5}{h}$ ($4,6\times 10^{-2}$) et $t = \SI{10}{h}$ ($7,6\times 10^{-2}$). Par interpolation linéaire :
\[ t_{1/2} \approx 5 + (10-5)\times\frac{6,65-4,6}{7,6-4,6} = 5 + 5\times\frac{2,05}{3,0} \approx \SI{8,4}{h} \]

\emph{Avec \ce{H2SO4} :} $x = \SI{6,65e-2}{mol}$ est atteint dès $t \approx \SI{2}{h}$ (valeur $6,8\times 10^{-2}$). Donc $t_{1/2} \approx \SI{2}{h}$.

\textbf{3. Rôle de l'acide sulfurique.}
Le temps de demi-réaction passe d'environ \SI{8,4}{h} à \SI{2}{h} : la réaction est environ \textbf{4 fois plus rapide}. Pourtant l'avancement final est identique ($x_f = \SI{0,133}{mol}$) : l'asymptote de la courbe $x=f(t)$ est la même, elle est simplement atteinte plus vite. L'acide sulfurique (source d'ions \ce{H3O+}) est donc un \cle{catalyseur} : il accélère l'estérification \emph{et} l'hydrolyse inverse, sans déplacer l'équilibre.

\textbf{Conclusion :} L'estérification est lente et limitée ($x_f < x_{\max}$) ; le catalyseur \ce{H2SO4} réduit fortement $t_{1/2}$ sans modifier la limite $x_f$, ce qui illustre parfaitement la définition d'un catalyseur.
\end{exoresolu}

\begin{attention}[Les 5 erreurs qui coûtent des points au Bac]
\begin{enumerate}
\item \textbf{Oublier le signe moins pour la vitesse de disparition.} Une vitesse est toujours \emph{positive} : pour un réactif, écrire $v = -\dfrac{\mathrm{d}[R]}{\mathrm{d}t}$ ou $v_m = -\dfrac{\Delta[R]}{\Delta t}$.
\item \textbf{Calculer la pente de la tangente avec des points de la courbe.} Il faut lire les coordonnées de deux points \emph{sur la tangente tracée}, jamais sur la courbe elle-même.
\item \textbf{Négliger les coefficients stœchiométriques.} $v = \dfrac{1}{\nu_P}\dfrac{\mathrm{d}[P]}{\mathrm{d}t}$ : pour $\ce{2 H2O2 -> 2 H2O + O2}$, la vitesse de disparition de \ce{H2O2} vaut le \emph{double} de la vitesse volumique de réaction ($v(\ce{H2O2}) = 2v$).
\item \textbf{Confondre $t_{1/2}$ avec le temps de fin de réaction.} Le temps de demi-réaction correspond à $x = x_f/2$ (la moitié de l'avancement \emph{final}, pas de $x_{\max}$ si la réaction est limitée) ; il est toujours inférieur au temps de fin.
\item \textbf{Dire qu'un catalyseur augmente le rendement.} Faux ! Un catalyseur (ou une élévation de température) fait atteindre l'état final \emph{plus vite}, sans modifier $x_f$ ni l'asymptote. Seul un excès de réactif (ou l'élimination d'un produit) déplace l'équilibre et améliore le rendement de l'estérification.
\end{enumerate}
\end{attention}

\newpage

\coursec{Exercices}

\courssub{Je vérifie mes connaissances}

\begin{exercice}[QCM : réactions rapides et lentes]
Pour chaque question, une seule réponse est exacte. Recopie la lettre correspondante.
\begin{enumerate}
\item La réaction entre les ions \ce{Ag+} et les ions \ce{Cl-} en solution aqueuse est :
\begin{enumerate}
\item une réaction lente ; \item une réaction rapide ; \item une réaction très lente ; \item une réaction impossible.
\end{enumerate}
\item La formation de la rouille sur une barre de fer exposée à l'air humide est :
\begin{enumerate}
\item instantanée ; \item rapide ; \item lente ; \item explosive.
\end{enumerate}
\item La synthèse de l'eau à partir du mélange \ce{H2} + \ce{O2}, en l'absence de catalyseur et à température ambiante, est :
\begin{enumerate}
\item rapide ; \item lente ; \item très lente (quasi nulle) ; \item totale en quelques secondes.
\end{enumerate}
\item Un catalyseur :
\begin{enumerate}
\item augmente la quantité finale de produit formé ; \item augmente la vitesse de réaction sans modifier l'état final ; \item est consommé pendant la réaction ; \item rend totale une réaction limitée.
\end{enumerate}
\end{enumerate}
\end{exercice}

\begin{exercice}[Vrai ou faux justifié]
Pour chaque affirmation, indique si elle est vraie ou fausse, puis justifie brièvement.
\begin{enumerate}
\item La vitesse instantanée de formation d'un produit à la date $t$ est égale au coefficient directeur de la tangente à la courbe $[\text{produit}] = f(t)$ en ce point.
\item La vitesse de disparition d'un réactif est une grandeur négative.
\item Le temps de demi-réaction $t_{1/2}$ est la date à laquelle l'avancement vaut la moitié de l'avancement final.
\item Augmenter la température d'un milieu réactionnel ralentit toujours la réaction.
\item Dans l'estérification de l'acide éthanoïque par l'éthanol, l'ajout d'acide sulfurique concentré augmente la quantité d'ester formée à l'équilibre.
\end{enumerate}
\end{exercice}

\begin{exercice}[Texte à trous]
Recopie et complète le texte suivant avec les mots ou expressions de la liste : \emph{tangente ; volumique ; trempe ; demi-réaction ; homogène ; hétérogène ; colorimétrie ; facteurs cinétiques}.

\medskip
Pour suivre l'évolution d'une réaction lente, on peut prélever des échantillons du milieu réactionnel que l'on soumet à une \ldots\ldots\ldots\ldots (refroidissement brutal et dilution) afin de bloquer la réaction avant le dosage. On peut aussi utiliser des méthodes physiques comme la \ldots\ldots\ldots\ldots lorsqu'une espèce colorée se forme. La vitesse \ldots\ldots\ldots\ldots de réaction est le quotient de la vitesse de réaction par le volume du milieu. La vitesse instantanée se détermine graphiquement grâce à la \ldots\ldots\ldots\ldots à la courbe. Le temps de \ldots\ldots\ldots\ldots est la durée nécessaire pour consommer la moitié du réactif limitant. La catalyse est dite \ldots\ldots\ldots\ldots lorsque le catalyseur et les réactifs sont dans la même phase (exemple : \ce{Fe^3+} sur le couple \ce{I-}/\ce{S2O8^2-}), et \ldots\ldots\ldots\ldots lorsqu'ils sont dans des phases différentes (exemple : le platine sur la synthèse de l'eau). La concentration, la température, le catalyseur et la surface de contact sont des \ldots\ldots\ldots\ldots .
\end{exercice}

\begin{exercice}[Définitions]
Définis en une ou deux phrases chacun des termes suivants, en précisant les unités usuelles quand c'est pertinent :
\begin{enumerate}
\item vitesse moyenne de formation d'un produit ;
\item vitesse instantanée de disparition d'un réactif ;
\item vitesse volumique de réaction ;
\item temps de demi-réaction ;
\item catalyseur.
\end{enumerate}
\end{exercice}

\courssub{Je m'entraîne}

\begin{exercice}[Vitesse moyenne de formation du diiode]
Lors de la réaction entre les ions iodure et les ions peroxodisulfate :
\[ \ce{2I- + S2O8^2- -> I2 + 2SO4^2-} \]
le suivi colorimétrique donne les valeurs suivantes :
\begin{center}
\begin{tabularx}{\linewidth}{|l|*{6}{>{\centering\arraybackslash}X|}}
\hline
\rowcolor{popblueL} $t$ (min) & 0 & 2 & 4 & 6 & 8 & 10 \\
\hline
$[\ce{I2}]$ (\SI{e-3}{mol.L^{-1}}) & 0 & 1,8 & 3,1 & 4,0 & 4,6 & 5,0 \\
\hline
\end{tabularx}
\end{center}
\begin{enumerate}
\item Calcule la vitesse moyenne de formation du diiode entre $t_1 = \SI{2}{min}$ et $t_2 = \SI{8}{min}$.
\item Calcule la vitesse moyenne de formation entre $t = 0$ et $t = \SI{10}{min}$.
\item Compare ces deux vitesses et explique l'évolution de la vitesse au cours du temps.
\end{enumerate}
\end{exercice}

\begin{exercice}[Vitesse volumique et avancement]
On réalise, à température ambiante, un mélange de volume total $V = \SI{100}{mL}$ contenant des ions iodure et des ions peroxodisulfate. À la date $t = \SI{5}{min}$, l'avancement de la réaction \ce{2I- + S2O8^2- -> I2 + 2SO4^2-} vaut $x = \SI{4,0e-4}{mol}$ ; à la date $t' = \SI{15}{min}$, il vaut $x' = \SI{8,0e-4}{mol}$.
\begin{enumerate}
\item Calcule la vitesse volumique moyenne de la réaction entre ces deux dates, en \si{mol.L^{-1}.min^{-1}}.
\item En utilisant les coefficients stœchiométriques, déduis-en les vitesses volumiques moyennes de disparition des ions \ce{I-} et \ce{S2O8^2-} sur le même intervalle.
\item Quelle est la concentration en diiode formé à la date $t' = \SI{15}{min}$ ?
\end{enumerate}
\end{exercice}

\begin{exercice}[Interprétation d'expériences témoins]
On réalise les quatre expériences suivantes :
\begin{itemize}
\item \textbf{Expérience 1 :} on mélange une solution d'iodure de potassium et une solution de peroxodisulfate de potassium ; la coloration brune apparaît lentement.
\item \textbf{Expérience 2 :} on refait le même mélange en ajoutant quelques gouttes de solution de chlorure de fer(III) ; la coloration brune apparaît beaucoup plus vite.
\item \textbf{Expérience 3 :} un mélange de dihydrogène et de dioxygène reste inchangé pendant des mois dans un ballon ; on introduit une mousse de platine : une réaction vive se produit.
\item \textbf{Expérience 4 :} on verse de l'acide chlorhydrique sur un bloc de zinc, puis la même quantité d'acide sur la même masse de zinc en poudre ; le dégagement gazeux est nettement plus rapide dans le second cas.
\end{itemize}
\begin{enumerate}
\item Pour chaque expérience, identifie le facteur cinétique mis en évidence.
\item Précise, dans les expériences 2 et 3, la nature de la catalyse (homogène ou hétérogène) en justifiant.
\item L'ion \ce{Fe^3+} modifie-t-il la quantité finale de diiode formé dans l'expérience 2 ? Justifie.
\end{enumerate}
\end{exercice}

\begin{exercice}[Courbes d'estérification]
On étudie la réaction entre l'acide éthanoïque et l'éthanol :
\[ \ce{CH3COOH + C2H5OH <=> CH3COOC2H5 + H2O} \]
On part dans les trois cas d'un mélange équimolaire ($n_0 = \SI{0,10}{mol}$ de chaque réactif). Trois courbes donnant la quantité d'ester formé en fonction du temps sont obtenues dans les conditions suivantes :
\begin{itemize}
\item courbe (a) : l'asymptote est atteinte au bout de \SI{50}{h} environ ;
\item courbe (b) : la même asymptote est atteinte au bout de \SI{8}{h} environ ;
\item courbe (c) : la même asymptote est atteinte au bout de \SI{15}{h} environ.
\end{itemize}
Conditions expérimentales utilisées : (1) mélange à \SI{20}{\celsius} ; (2) mélange à \SI{60}{\celsius} ; (3) mélange à \SI{20}{\celsius} avec quelques gouttes d'acide sulfurique concentré.
\begin{enumerate}
\item Associe chaque courbe à sa condition expérimentale en justifiant par la pente initiale et la position de l'asymptote.
\item Pourquoi les trois courbes présentent-elles la même asymptote ? Calcule la valeur de la quantité d'ester à l'équilibre sachant que le rendement de l'estérification est de \SI{67}{\%} pour un mélange équimolaire d'acide primaire et d'alcool primaire.
\item Comment pourrait-on augmenter la valeur de l'asymptote sans changer la température ni le catalyseur ?
\end{enumerate}
\end{exercice}

\courssub{Je me prépare au Bac}

\begin{exercice}[Type Bac : cinétique de la réaction iodure--peroxodisulfate]
\textit{(D'après Baccalauréat C, session normale)}

Au laboratoire du lycée, un groupe d'élèves de Terminale C étudie la cinétique de la réaction d'oxydation des ions iodure par les ions peroxodisulfate :
\[ \ce{2I- + S2O8^2- -> I2 + 2SO4^2-} \]
Ils mélangent, à \SI{25}{\celsius}, un volume $V_1 = \SI{50,0}{mL}$ d'une solution d'iodure de potassium de concentration $C_1 = \SI{0,50}{mol.L^{-1}}$ et un volume $V_2 = \SI{50,0}{mL}$ d'une solution de peroxodisulfate de potassium de concentration $C_2 = \SI{0,10}{mol.L^{-1}}$. Le diiode formé, espèce colorée, est dosé par spectrophotométrie. Les résultats sont consignés dans le tableau ci-dessous :
\begin{center}
\begin{tabularx}{\linewidth}{|l|*{9}{>{\centering\arraybackslash}X|}}
\hline
\rowcolor{popblueL} $t$ (min) & 0 & 5 & 10 & 15 & 20 & 30 & 40 & 50 & 60 \\
\hline
$[\ce{I2}]$ (\SI{e-3}{mol.L^{-1}}) & 0 & 9,0 & 15,6 & 20,6 & 24,4 & 29,6 & 32,8 & 34,7 & 35,8 \\
\hline
\end{tabularx}
\end{center}
\begin{enumerate}
\item Calcule les quantités initiales $n_1$ d'ions \ce{I-} et $n_2$ d'ions \ce{S2O8^2-} introduites. Construis le tableau d'avancement et montre que l'ion peroxodisulfate est le réactif limitant. En déduire l'avancement maximal $x_{\max}$.
\item Trace, sur papier millimétré, la courbe $[\ce{I2}] = f(t)$. Échelle : \SI{1}{cm} pour \SI{5}{min} en abscisse ; \SI{1}{cm} pour \SI{4e-3}{mol.L^{-1}} en ordonnée.
\item Définis la vitesse volumique moyenne de formation du diiode, puis calcule sa valeur entre $t_1 = \SI{10}{min}$ et $t_2 = \SI{30}{min}$.
\item Détermine graphiquement la vitesse volumique instantanée de formation du diiode à la date $t = \SI{20}{min}$ (on tracera la tangente sur le graphique). Compare cette valeur à la vitesse moyenne de la question 3 et commente.
\item Déduis de la question 4 les vitesses instantanées de disparition des ions \ce{I-} et \ce{S2O8^2-} à cette même date.
\item Définis le temps de demi-réaction $t_{1/2}$ et détermine sa valeur graphiquement.
\item On refait l'expérience en ajoutant quelques gouttes de solution de chlorure de fer(III). Sans refaire le tracé, représente en pointillés, sur le même graphique, l'allure de la nouvelle courbe. Justifie. Le temps de demi-réaction est-il modifié ? La concentration finale en diiode est-elle modifiée ?
\end{enumerate}
\emph{Données :} masses molaires atomiques en \si{g.mol^{-1}} : $M(\ce{K}) = 39,1$ ; $M(\ce{I}) = 126,9$ ; $M(\ce{S}) = 32,1$ ; $M(\ce{O}) = 16,0$.
\end{exercice}

\begin{exercice}[Type Bac : suivi d'un dégagement gazeux par mesure de pression]
\textit{(D'après Baccalauréat D, session de rattrapage)}

Pour étudier la décomposition de l'eau oxygénée (eau oxygénée du commerce utilisée comme désinfectant dans les centres de santé), on introduit dans un ballon fermé de volume $V = \SI{1,00}{L}$ un volume $V_0 = \SI{20,0}{mL}$ d'une solution d'eau oxygénée de concentration $C_0 = \SI{0,89}{mol.L^{-1}}$, en présence d'un catalyseur. La réaction, supposée totale, est :
\[ \ce{2H2O2(aq) -> 2H2O(l) + O2(g)} \]
On mesure la surpression $\Delta p$ due au dioxygène formé, à la température constante $\theta = \SI{25}{\celsius}$. Le volume gazeux disponible dans le ballon est $V_g = \SI{0,98}{L}$.
\begin{center}
\begin{tabularx}{\linewidth}{|l|*{7}{>{\centering\arraybackslash}X|}}
\hline
\rowcolor{popblueL} $t$ (min) & 0 & 5 & 10 & 15 & 20 & 30 & 40 \\
\hline
$\Delta p$ (hPa) & 0 & 86 & 150 & 199 & 236 & 286 & 316 \\
\hline
\end{tabularx}
\end{center}
\begin{enumerate}
\item Calcule la quantité initiale $n_0$ d'eau oxygénée. En déduire la quantité maximale de dioxygène $n_{\max}(\ce{O2})$ attendue.
\item En assimilant le dioxygène à un gaz parfait, montre que la quantité de dioxygène formé à la date $t$ est donnée par $n(\ce{O2}) = \dfrac{\Delta p \cdot V_g}{R\,T}$. Calcule $n(\ce{O2})$ aux dates $t = \SI{10}{min}$ et $t = \SI{40}{min}$.
\item Exprime l'avancement $x$ de la réaction en fonction de $n(\ce{O2})$, puis la concentration restante en eau oxygénée $[\ce{H2O2}]$ en fonction de $x$ et $V_0$. Calcule $[\ce{H2O2}]$ à $t = \SI{10}{min}$.
\item Trace la courbe $[\ce{H2O2}] = f(t)$ après avoir complété le tableau des concentrations pour toutes les dates. Échelle : \SI{1}{cm} pour \SI{5}{min} ; \SI{1}{cm} pour \SI{0,10}{mol.L^{-1}}.
\item Détermine graphiquement la vitesse volumique instantanée de disparition de l'eau oxygénée à $t = \SI{15}{min}$.
\item Détermine graphiquement le temps de demi-réaction. Vérifie la cohérence avec la valeur de $n_{\max}(\ce{O2})$.
\item Cite deux moyens expérimentaux permettant d'atteindre plus rapidement l'état final, sans modifier les quantités de réactifs.
\end{enumerate}
\emph{Données :} $R = \SI{8,314}{J.mol^{-1}.K^{-1}}$ ; $T(\si{K}) = \theta(\si{\celsius}) + 273$ ; \SI{1}{hPa} = \SI{100}{Pa}. Masses molaires en \si{g.mol^{-1}} : $M(\ce{H}) = 1,0$ ; $M(\ce{O}) = 16,0$.
\end{exercice}

\begin{exercice}[Type Bac : étude cinétique de l'estérification et vie pratique]
\textit{(Sujet de synthèse, type Baccalauréat C/D/E)}

Une coopérative de la région du Littoral souhaite produire de l'éthanoate d'éthyle, solvant utilisé dans les arômes alimentaires. Un technicien étudie au laboratoire la réaction :
\[ \ce{CH3COOH + C2H5OH <=> CH3COOC2H5 + H2O} \]
Il prépare trois ballons identiques contenant chacun $n_a = \SI{0,20}{mol}$ d'acide éthanoïque et $n_b = \SI{0,20}{mol}$ d'éthanol, dans les conditions suivantes :
\begin{itemize}
\item \textbf{Ballon A :} chauffage à reflux à \SI{70}{\celsius} ;
\item \textbf{Ballon B :} chauffage à reflux à \SI{70}{\celsius} avec \SI{1}{mL} d'acide sulfurique concentré ;
\item \textbf{Ballon C :} mélange maintenu à \SI{20}{\celsius}.
\end{itemize}
Le dosage de l'acide restant par une solution d'hydroxyde de sodium permet de suivre la quantité d'ester formé. Les courbes $n(\text{ester}) = f(t)$ montrent que :
\begin{itemize}
\item les ballons A et B atteignent la même asymptote $n_{\infty} = \SI{0,134}{mol}$, le ballon B en \SI{3}{h} et le ballon A en \SI{12}{h} ;
\item le ballon C n'a pas encore atteint l'asymptote au bout de \SI{48}{h}.
\end{itemize}
\begin{enumerate}
\item Écris l'équation-bilan de la réaction et donne le nom de sa réaction inverse. Cite deux caractères de cette transformation.
\item Construis le tableau d'avancement du système et calcule l'avancement maximal $x_{\max}$. En déduire l'avancement final $x_f$ et le taux d'avancement final $\tau$ de la réaction dans les ballons A et B. La réaction est-elle totale ?
\item Explique, en t'appuyant sur les facteurs cinétiques, pourquoi l'asymptote est atteinte plus vite dans le ballon B que dans le ballon A, et plus lentement dans le ballon C. Précise la nature de la catalyse mise en jeu dans le ballon B.
\item Le technicien affirme : « L'acide sulfurique a augmenté le rendement de la production d'ester. » Cette affirmation est-elle correcte ? Justifie à l'aide des valeurs expérimentales.
\item Propose deux modifications du protocole permettant d'augmenter la valeur de l'asymptote, en justifiant par le déplacement de l'équilibre.
\item \textbf{Ouverture sur la vie quotidienne :} explique, en invoquant les facteurs cinétiques, chacune des trois pratiques suivantes : (a) la conservation du poisson au frais dans les glacières des vendeuses du marché ; (b) le broyage du charbon de bois pour faciliter l'allumage du feu ; (c) l'ajout de levure dans la pâte à beignets pour accélérer sa fermentation.
\end{enumerate}
\emph{Données :} masses molaires en \si{g.mol^{-1}} : $M(\ce{CH3COOH}) = 60,0$ ; $M(\ce{C2H5OH}) = 46,0$ ; $M(\ce{CH3COOC2H5}) = 88,0$ ; $M(\ce{H2O}) = 18,0$.
\end{exercice}



\newpage

\coursec{Corrigés des exercices}

\begin{corrige}[Exercice 1 — QCM : réactions rapides et lentes]
\begin{enumerate}
\item \textbf{b.} La réaction \ce{Ag+ + Cl- -> AgCl(s)} est une réaction de précipitation, quasi instantanée : c'est une réaction \cle{rapide}. Dès le mélange des solutions, un précipité blanc de chlorure d'argent apparaît (ce test sert d'ailleurs à caractériser les ions chlorure).
\item \textbf{c.} La formation de la rouille (corrosion du fer par le dioxygène en présence d'eau) s'étale sur des semaines, des mois, voire des années : c'est une réaction \cle{lente}. On l'observe bien sur les tôles et les barres de fer des chantiers exposées à l'air humide du littoral camerounais.
\item \textbf{c.} Le mélange \ce{H2} + \ce{O2} peut rester des années sans réagir à température ambiante : la réaction est \cle{très lente} (cinétiquement bloquée). Une simple étincelle ou une mousse de platine la déclenche brutalement.
\item \textbf{b.} Un catalyseur \cle{augmente la vitesse} de réaction sans modifier l'état final : il abaisse l'énergie d'activation, est régénéré en fin de mécanisme (donc non consommé globalement) et ne déplace pas l'équilibre (la constante d'équilibre $K$, qui ne dépend que de la température, est inchangée).
\end{enumerate}
\end{corrige}

\begin{corrige}[Exercice 2 — Vrai ou faux justifié]
\begin{enumerate}
\item \textbf{Vrai.} Par définition, la vitesse instantanée de formation d'un produit $P$ est la dérivée $v_{\text{form}}(P) = \dfrac{d[P]}{dt}$. Or la dérivée en un point est précisément le coefficient directeur (la pente) de la tangente à la courbe $[P] = f(t)$ en ce point. C'est le fondement de la détermination graphique des vitesses instantanées.
\item \textbf{Faux.} Par convention, une vitesse de réaction est toujours une grandeur \textbf{positive}. Pour un réactif $R$, on définit $v_{\text{disp}}(R) = -\dfrac{1}{\nu_R}\dfrac{d[R]}{dt} > 0$. C'est la dérivée $\dfrac{d[R]}{dt}$ qui est négative (car $[R]$ diminue) ; le signe moins de la définition rétablit une valeur positive.
\item \textbf{Vrai.} Le temps de demi-réaction $t_{1/2}$ est la date à laquelle l'avancement atteint la moitié de sa valeur finale : $x(t_{1/2}) = \dfrac{x_{\max}}{2}$ (pour une réaction totale) ou $\dfrac{x_f}{2}$ (pour une réaction limitée). À cet instant, la moitié du réactif limitant initialement présent a été consommée.
\item \textbf{Faux.} C'est l'inverse : l'élévation de température augmente l'énergie cinétique des particules, donc la fréquence et l'efficacité des chocs : la réaction est \textbf{accélérée}. C'est pourquoi on conserve les aliments au frais (ralentissement) et on cuit les aliments (accélération).
\item \textbf{Faux.} L'acide sulfurique concentré est un \cle{catalyseur} (catalyse homogène par les ions \ce{H+}). Il permet d'atteindre plus vite l'équilibre, mais ne modifie ni la constante d'équilibre ni la quantité d'ester à l'équilibre : l'asymptote de la courbe $n(\text{ester}) = f(t)$ reste la même. Pour augmenter le rendement, il faut déplacer l'équilibre (excès d'un réactif, élimination d'un produit).
\end{enumerate}
\end{corrige}

\begin{corrige}[Exercice 3 — Texte à trous]
Pour suivre l'évolution d'une réaction lente, on peut prélever des échantillons du milieu réactionnel que l'on soumet à une \textbf{trempe} (refroidissement brutal et dilution) afin de bloquer la réaction avant le dosage. On peut aussi utiliser des méthodes physiques comme la \textbf{colorimétrie} lorsqu'une espèce colorée se forme. La vitesse \textbf{volumique} de réaction est le quotient de la vitesse de réaction par le volume du milieu. La vitesse instantanée se détermine graphiquement grâce à la \textbf{tangente} à la courbe. Le temps de \textbf{demi-réaction} est la durée nécessaire pour consommer la moitié du réactif limitant. La catalyse est dite \textbf{homogène} lorsque le catalyseur et les réactifs sont dans la même phase (exemple : \ce{Fe^3+} sur le couple \ce{I-}/\ce{S2O8^2-}), et \textbf{hétérogène} lorsqu'ils sont dans des phases différentes (exemple : le platine sur la synthèse de l'eau). La concentration, la température, le catalyseur et la surface de contact sont des \textbf{facteurs cinétiques}.
\end{corrige}

\begin{corrige}[Exercice 4 — Définitions]
\begin{enumerate}
\item \textbf{Vitesse moyenne de formation d'un produit $P$} entre les dates $t_1$ et $t_2$ : c'est le quotient de la variation de concentration du produit par la durée écoulée :
\[ v_{\text{moy}}(P) = \frac{[P]_{t_2} - [P]_{t_1}}{t_2 - t_1} \]
Unité usuelle : \si{mol.L^{-1}.min^{-1}} (ou \si{mol.L^{-1}.s^{-1}}).
\item \textbf{Vitesse instantanée de disparition d'un réactif $R$} à la date $t$ : c'est l'opposé de la dérivée de sa concentration par rapport au temps (divisée par son coefficient stœchiométrique) :
\[ v_{\text{disp}}(R) = -\frac{1}{\nu_R}\frac{d[R]}{dt} \]
C'est une grandeur positive, égale en valeur absolue à la pente de la tangente à la courbe $[R] = f(t)$ à la date $t$. Unité : \si{mol.L^{-1}.min^{-1}}.
\item \textbf{Vitesse volumique de réaction} : c'est la dérivée de l'avancement $x$ par rapport au temps, rapportée au volume $V$ du milieu réactionnel :
\[ v = \frac{1}{V}\frac{dx}{dt} \]
Unité : \si{mol.L^{-1}.min^{-1}} (ou \si{mol.L^{-1}.s^{-1}}). Elle est indépendante du volume du milieu, ce qui permet de comparer des expériences.
\item \textbf{Temps de demi-réaction $t_{1/2}$} : c'est la durée au bout de laquelle l'avancement atteint la moitié de sa valeur finale : $x(t_{1/2}) = \dfrac{x_{\max}}{2}$. Unité : seconde, minute ou heure selon la réaction.
\item \textbf{Catalyseur} : c'est une espèce chimique qui augmente la vitesse d'une réaction sans être consommée globalement (il est régénéré au cours du mécanisme) et sans modifier l'état final du système. Il agit en abaissant l'énergie d'activation de la réaction.
\end{enumerate}
\end{corrige}

\begin{corrige}[Exercice 5 — Vitesse moyenne de formation du diiode]
\begin{enumerate}
\item Entre $t_1 = \SI{2}{min}$ et $t_2 = \SI{8}{min}$ :
\[ v_{\text{moy}}(\ce{I2}) = \frac{[\ce{I2}]_{8} - [\ce{I2}]_{2}}{t_2 - t_1} = \frac{(4,6 - 1,8)\times 10^{-3}}{8 - 2} = \frac{2,8\times 10^{-3}}{6} \approx \SI{4,7e-4}{mol.L^{-1}.min^{-1}} \]
\item Entre $t = 0$ et $t = \SI{10}{min}$ :
\[ v_{\text{moy}}(\ce{I2}) = \frac{(5,0 - 0)\times 10^{-3}}{10 - 0} = \SI{5,0e-4}{mol.L^{-1}.min^{-1}} \]
\item \textbf{Comparaison et interprétation.} La vitesse moyenne sur les dix premières minutes (\SI{5,0e-4}{mol.L^{-1}.min^{-1}}) est supérieure à celle mesurée entre 2 et 8 min (\SI{4,7e-4}{mol.L^{-1}.min^{-1}}). De manière générale, la vitesse de formation du diiode \textbf{diminue au cours du temps} : la courbe $[\ce{I2}] = f(t)$ est de plus en plus aplatie (sa pente décroît). \textbf{Explication :} au fur et à mesure de la réaction, les concentrations des réactifs \ce{I-} et \ce{S2O8^2-} diminuent ; or la concentration est un facteur cinétique : moins il y a d'ions en solution, moins les chocs efficaces sont fréquents, et plus la réaction ralentit.
\end{enumerate}
\end{corrige}

\begin{corrige}[Exercice 6 — Vitesse volumique et avancement]
\textbf{Données :} $V = \SI{100}{mL} = \SI{0,100}{L}$ ; $x(\SI{5}{min}) = \SI{4,0e-4}{mol}$ ; $x'(\SI{15}{min}) = \SI{8,0e-4}{mol}$.
\begin{enumerate}
\item Vitesse volumique moyenne entre $t = \SI{5}{min}$ et $t' = \SI{15}{min}$ :
\[ v_{\text{moy}} = \frac{1}{V}\cdot\frac{x' - x}{t' - t} = \frac{1}{0,100}\times\frac{(8,0 - 4,0)\times 10^{-4}}{15 - 5} = \frac{4,0\times 10^{-4}}{1,0} = \SI{4,0e-4}{mol.L^{-1}.min^{-1}} \]
\item Pour la réaction \ce{2I- + S2O8^2- -> I2 + 2SO4^2-}, la vitesse volumique s'écrit :
\[ v = -\frac{1}{2}\frac{d[\ce{I-}]}{dt} = -\frac{d[\ce{S2O8^2-}]}{dt} = \frac{d[\ce{I2}]}{dt} \]
On en déduit, sur le même intervalle :
\begin{itemize}
\item vitesse volumique moyenne de disparition des ions iodure :
\[ v_{\text{disp}}(\ce{I-}) = 2\,v = 2 \times \SI{4,0e-4}{} = \SI{8,0e-4}{mol.L^{-1}.min^{-1}} \]
\item vitesse volumique moyenne de disparition des ions peroxodisulfate :
\[ v_{\text{disp}}(\ce{S2O8^2-}) = v = \SI{4,0e-4}{mol.L^{-1}.min^{-1}} \]
\end{itemize}
\item D'après le tableau d'avancement, $n(\ce{I2}) = x'$. Donc :
\[ [\ce{I2}]_{t'} = \frac{x'}{V} = \frac{8,0\times 10^{-4}}{0,100} = \SI{8,0e-3}{mol.L^{-1}} \]
\end{enumerate}
\end{corrige}

\begin{attention}[Point de vigilance — Exercice 6]
Ne pas oublier les coefficients stœchiométriques : la disparition de \ce{I-} est deux fois plus rapide que la formation de \ce{I2}, car il faut 2 ions iodure pour former 1 molécule de diiode.
\end{attention}

\begin{corrige}[Exercice 7 — Interprétation d'expériences témoins]
\begin{enumerate}
\item Facteurs cinétiques mis en évidence :
\begin{itemize}
\item \textbf{Expérience 1 :} expérience témoin ; la réaction \ce{2I- + S2O8^2- -> I2 + 2SO4^2-} est lente dans les conditions ambiantes (concentrations et température ordinaires).
\item \textbf{Expérience 2 :} \cle{catalyse} par les ions \ce{Fe^3+} : la coloration brune du diiode apparaît beaucoup plus vite.
\item \textbf{Expérience 3 :} \cle{catalyse} par le platine : le mélange \ce{H2} + \ce{O2}, cinétiquement bloqué, réagit vivement en sa présence.
\item \textbf{Expérience 4 :} \cle{surface de contact} : le zinc en poudre offre une surface de contact avec l'acide bien plus grande que le zinc en bloc, d'où un dégagement de \ce{H2} plus rapide.
\end{itemize}
\item Nature des catalyses :
\begin{itemize}
\item \textbf{Expérience 2 — catalyse homogène :} les ions \ce{Fe^3+} et les réactifs \ce{I-} et \ce{S2O8^2-} sont tous en solution aqueuse : catalyseur et réactifs appartiennent à la \textbf{même phase} (liquide).
\item \textbf{Expérience 3 — catalyse hétérogène :} le platine est un \textbf{solide} alors que \ce{H2} et \ce{O2} sont des \textbf{gaz} : catalyseur et réactifs sont dans des \textbf{phases différentes}.
\end{itemize}
\item \textbf{Non}, l'ion \ce{Fe^3+} ne modifie pas la quantité finale de diiode. Un catalyseur accélère la réaction (il est régénéré en fin de mécanisme) mais n'apparaît pas dans l'équation-bilan et ne change ni les quantités initiales de réactifs, ni le réactif limitant. La quantité finale de \ce{I2}, fixée par l'avancement maximal, reste donc la même : seule la \textbf{durée} pour l'atteindre est raccourcie.
\end{enumerate}
\end{corrige}

\begin{corrige}[Exercice 8 — Courbes d'estérification]
\begin{enumerate}
\item Association courbe $\leftrightarrow$ condition :
\begin{itemize}
\item \textbf{Courbe (a) — condition (1)} (\SI{20}{\celsius}, sans catalyseur) : ni température élevée ni catalyseur, la réaction est la plus lente ; la pente initiale est la plus faible et l'asymptote n'est atteinte qu'au bout d'environ \SI{50}{h}.
\item \textbf{Courbe (b) — condition (2)} (\SI{60}{\celsius}) : l'élévation de température accélère fortement la réaction ; la pente initiale est la plus grande, l'asymptote est atteinte en \SI{8}{h} environ.
\item \textbf{Courbe (c) — condition (3)} (\SI{20}{\celsius} + \ce{H2SO4}) : le catalyseur accélère la réaction sans la rendre aussi rapide qu'à \SI{60}{\celsius} ; l'asymptote est atteinte en \SI{15}{h} environ.
\end{itemize}
Dans les trois cas, la pente initiale diffère (facteurs cinétiques différents) mais l'asymptote est identique.
\item Les trois courbes présentent la \textbf{même asymptote} car, conformément au programme, l'élévation de température et l'ajout de catalyseur n'entraînent pas de changement de la limite (l'avancement final) dans le cadre de cette étude. Le catalyseur ne déplace jamais un équilibre, il ne fait qu'accélérer son atteinte. La quantité d'ester à l'équilibre ne dépend que des quantités initiales et du taux d'avancement final.
Avec un rendement (taux d'avancement final) $\tau = \SI{67}{\%}$ :
\[ n_{\text{ester, éq}} = \tau \times n_0 = 0,67 \times 0,10 = \SI{6,7e-2}{mol} \]
\item Pour augmenter la valeur de l'asymptote (c'est-à-dire déplacer l'équilibre dans le sens direct), sans changer ni la température ni le catalyseur, on peut :
\begin{itemize}
\item introduire l'un des réactifs en \textbf{excès} (par exemple un excès d'éthanol, réactif bon marché) ;
\item \textbf{éliminer un produit} au fur et à mesure de sa formation (distillation de l'ester, le plus volatil, ou piégeage de l'eau par un agent déshydratant).
\end{itemize}
Conformément à la loi de Le Chatelier, l'équilibre se déplace alors dans le sens de la formation de l'ester et le rendement augmente.
\end{enumerate}
\end{corrige}

\begin{corrige}[Exercice 9 — Type Bac : cinétique de la réaction iodure–peroxodisulfate]
\begin{enumerate}
\item \textbf{Quantités initiales.} Le mélange a un volume total $V = V_1 + V_2 = \SI{100,0}{mL} = \SI{0,1000}{L}$.
\[ n_1(\ce{I-}) = C_1 V_1 = 0,50 \times 50,0\times 10^{-3} = \SI{2,5e-2}{mol} \]
\[ n_2(\ce{S2O8^2-}) = C_2 V_2 = 0,10 \times 50,0\times 10^{-3} = \SI{5,0e-3}{mol} \]
\textbf{Tableau d'avancement} (quantités en mol) :
\begin{center}
\begin{tabular}{|l|c|c|c|c|}
\hline
\rowcolor{popblueL} Équation & \ce{2I-} & \ce{+ S2O8^2- ->} & \ce{I2} & \ce{+ 2SO4^2-} \\
\hline
État initial ($x=0$) & $2,5\times 10^{-2}$ & $5,0\times 10^{-3}$ & $0$ & $0$ \\
\hline
En cours ($x$) & $2,5\times 10^{-2} - 2x$ & $5,0\times 10^{-3} - x$ & $x$ & $2x$ \\
\hline
État final ($x_{\max}$) & $1,5\times 10^{-2}$ & $0$ & $5,0\times 10^{-3}$ & $1,0\times 10^{-2}$ \\
\hline
\end{tabular}
\end{center}
Hypothèses : si \ce{I-} est limitant, $x_{\max} = \dfrac{2,5\times 10^{-2}}{2} = \SI{1,25e-2}{mol}$ ; si \ce{S2O8^2-} est limitant, $x_{\max} = \SI{5,0e-3}{mol}$. La plus petite valeur l'emporte : \ce{S2O8^2-} est le \textbf{réactif limitant} et
\[ x_{\max} = \SI{5,0e-3}{mol} \qquad [\ce{I2}]_{\max} = \frac{x_{\max}}{V} = \frac{5,0\times 10^{-3}}{0,1000} = \SI{5,0e-2}{mol.L^{-1}} \]
\item \textbf{Tracé de la courbe.} On obtient une courbe croissante, de pente décroissante, partant de l'origine et tendant vers l'asymptote horizontale $[\ce{I2}] = \SI{50e-3}{mol.L^{-1}}$ (les points expérimentaux n'atteignent pas encore l'asymptote à 60 min).
\item \textbf{Vitesse volumique moyenne de formation du diiode} entre $t_1$ et $t_2$ :
\[ v_{\text{moy}}(\ce{I2}) = \frac{[\ce{I2}]_{t_2} - [\ce{I2}]_{t_1}}{t_2 - t_1} \]
Entre $t_1 = \SI{10}{min}$ et $t_2 = \SI{30}{min}$ :
\[ v_{\text{moy}}(\ce{I2}) = \frac{(29,6 - 15,6)\times 10^{-3}}{30 - 10} = \frac{14,0\times 10^{-3}}{20} = \SI{7,0e-4}{mol.L^{-1}.min^{-1}} \]
\item \textbf{Vitesse instantanée à $t = \SI{20}{min}$.} On trace la tangente à la courbe au point d'abscisse $t = \SI{20}{min}$ et on détermine son coefficient directeur à partir de deux points éloignés de la tangente. Par exemple, si la tangente passe par les points $(0\ \text{min}\ ;\ 12,4\times 10^{-3}\ \si{mol.L^{-1}})$ et $(40\ \text{min}\ ;\ 36,4\times 10^{-3}\ \si{mol.L^{-1}})$ :
\[ v_{\text{inst}}(\ce{I2})_{20} = \frac{(36,4 - 12,4)\times 10^{-3}}{40 - 0} = \frac{24,0\times 10^{-3}}{40} = \SI{6,0e-4}{mol.L^{-1}.min^{-1}} \]
\textbf{Comparaison :} $v_{\text{inst}}(20\ \text{min}) < v_{\text{moy}}(10 \to 30\ \text{min})$. La vitesse instantanée diminue au cours du temps (la courbe s'aplatit) car les réactifs s'épuisent ; la vitesse moyenne sur $[10;30]$ intègre des vitesses plus grandes au début de l'intervalle, d'où une valeur supérieure à la vitesse au milieu de l'intervalle.
\item D'après la stœchiométrie : $v = v_{\text{form}}(\ce{I2}) = \dfrac{1}{2}v_{\text{disp}}(\ce{I-}) = v_{\text{disp}}(\ce{S2O8^2-})$. À $t = \SI{20}{min}$ :
\[ v_{\text{disp}}(\ce{I-}) = 2 \times 6,0\times 10^{-4} = \SI{1,2e-3}{mol.L^{-1}.min^{-1}} \]
\[ v_{\text{disp}}(\ce{S2O8^2-}) = v = \SI{6,0e-4}{mol.L^{-1}.min^{-1}} \]
\item \textbf{Temps de demi-réaction} : date à laquelle $x = \dfrac{x_{\max}}{2}$, c'est-à-dire $[\ce{I2}] = \dfrac{50\times 10^{-3}}{2} = \SI{25e-3}{mol.L^{-1}}$. Sur le graphique, on lit l'abscisse du point d'ordonnée \SI{25e-3}{mol.L^{-1}} : le point se situe juste après $t = \SI{20}{min}$ (où $[\ce{I2}] = \SI{24,4e-3}{mol.L^{-1}}$) :
\[ t_{1/2} \approx \SI{21}{min} \]
\item \textbf{Effet des ions \ce{Fe^3+}.} La nouvelle courbe (en pointillés) part de la même origine, monte \textbf{plus rapidement} (pente initiale plus grande) et rejoint la \textbf{même asymptote} $[\ce{I2}] = \SI{50e-3}{mol.L^{-1}}$ plus tôt. Justification : \ce{Fe^3+} est un catalyseur (catalyse homogène) qui abaisse l'énergie d'activation. Le temps de demi-réaction est \textbf{diminué} ($t'_{1/2} < t_{1/2}$). La concentration finale en diiode est \textbf{inchangée} : un catalyseur ne modifie ni le réactif limitant ni l'avancement maximal.
\end{enumerate}
\end{corrige}

\begin{attention}[Points de vigilance — Exercice 9 — Barème indicatif sur 10]
(1) 2 pts : quantités initiales, tableau d'avancement, réactif limitant ; (2) 1,5 pt : axes, échelles, points, courbe lissée ; (3) 1 pt : définition + calcul ; (4) 1,5 pt : tangente soignée + lecture de deux points éloignés ; (5) 1 pt : coefficients stœchiométriques ; (6) 1 pt : définition + lecture à \SI{25e-3}{mol.L^{-1}} ; (7) 2 pts : allure, $t_{1/2}$ diminué, asymptote inchangée. Erreurs fréquentes : oublier le facteur 2 pour \ce{I-}, confondre $t_{1/2}$ avec la moitié du temps total, tracer la tangente en un seul point.
\end{attention}

\begin{corrige}[Exercice 10 — Type Bac : suivi d'un dégagement gazeux par mesure de pression]
\begin{enumerate}
\item \textbf{Quantité initiale d'eau oxygénée :}
\[ n_0(\ce{H2O2}) = C_0 V_0 = 0,89 \times 20,0\times 10^{-3} = \SI{1,78e-2}{mol} \]
La réaction étant totale, d'après la stœchiométrie \ce{2H2O2 -> 2H2O + O2} :
\[ n_{\max}(\ce{O2}) = \frac{n_0}{2} = \SI{8,9e-3}{mol} \]
\item Le dioxygène est le seul gaz formé ; sa pression partielle est égale à la surpression $\Delta p$. L'équation des gaz parfaits appliquée à \ce{O2} dans le volume gazeux $V_g$ à la température $T$ donne :
\[ \Delta p \cdot V_g = n(\ce{O2})\,R\,T \quad\Longrightarrow\quad n(\ce{O2}) = \frac{\Delta p \cdot V_g}{R\,T} \]
avec $T = 25 + 273 = \SI{298}{K}$, $V_g = \SI{0,98}{L} = \SI{9,8e-4}{m^3}$, et $\Delta p$ en pascals.
À $t = \SI{10}{min}$ : $\Delta p = \SI{150}{hPa} = \SI{1,50e4}{Pa}$ :
\[ n(\ce{O2}) = \frac{1,50\times 10^{4} \times 9,8\times 10^{-4}}{8,314 \times 298} = \frac{14,7}{2478} \approx \SI{5,9e-3}{mol} \]
À $t = \SI{40}{min}$ : $\Delta p = \SI{316}{hPa} = \SI{3,16e4}{Pa}$ :
\[ n(\ce{O2}) = \frac{3,16\times 10^{4} \times 9,8\times 10^{-4}}{2478} = \frac{31,0}{2478} \approx \SI{1,25e-2}{mol} \]
\item D'après l'équation-bilan, $n(\ce{O2}) = x$, donc :
\[ x = n(\ce{O2}) = \frac{\Delta p \cdot V_g}{R\,T} \]
La quantité d'eau oxygénée restante est $n(\ce{H2O2}) = n_0 - 2x$, d'où :
\[ [\ce{H2O2}] = \frac{n_0 - 2x}{V_0} \]
À $t = \SI{10}{min}$ : $[\ce{H2O2}] = \dfrac{1,78\times 10^{-2} - 2\times 5,9\times 10^{-3}}{20,0\times 10^{-3}} = \dfrac{6,0\times 10^{-3}}{0,0200} \approx \SI{0,30}{mol.L^{-1}}$.
\item Tableau complété (avec $[\ce{H2O2}] = \dfrac{n_0 - 2n(\ce{O2})}{V_0}$) :
\begin{center}
\begin{tabularx}{\linewidth}{|l|*{7}{>{\centering\arraybackslash}X|}}
\hline
\rowcolor{popblueL} $t$ (min) & 0 & 5 & 10 & 15 & 20 & 30 & 40 \\
\hline
$n(\ce{O2})$ (\SI{e-3}{mol}) & 0 & 3,4 & 5,9 & 7,8 & 9,3 & 11,2 & 12,5 \\
\hline
$[\ce{H2O2}]$ (\si{mol.L^{-1}}) & 0,89 & 0,55 & 0,30 & 0,11 & 0 & 0 & 0 \\
\hline
\end{tabularx}
\end{center}
(À partir de $t \approx \SI{18}{min}$, le calcul donne $n(\ce{O2}) > n_{\max}$ : toute l'eau oxygénée est alors consommée et $[\ce{H2O2}] = 0$.) La courbe $[\ce{H2O2}] = f(t)$ est décroissante, de pente de plus en plus faible en valeur absolue, et rejoint l'axe des abscisses.
\item \textbf{Vitesse instantanée de disparition à $t = \SI{15}{min}$.} On trace la tangente à la courbe au point d'abscisse 15 min et on prend l'opposé de son coefficient directeur. Lecture typique : la tangente passe par $(5\ \text{min}\ ;\ 0,40\ \si{mol.L^{-1}})$ et $(25\ \text{min}\ ;\ 0\ \si{mol.L^{-1}})$ :
\[ v_{\text{disp}}(\ce{H2O2}) = -\frac{0 - 0,40}{25 - 5} = \frac{0,40}{20} = \SI{2,0e-2}{mol.L^{-1}.min^{-1}} \]
\item \textbf{Temps de demi-réaction.} À $t_{1/2}$, $[\ce{H2O2}] = \dfrac{0,89}{2} \approx \SI{0,45}{mol.L^{-1}}$. Lecture graphique : $t_{1/2} \approx \SI{6}{min}$.
\textbf{Vérification :} à $t_{1/2}$, $n(\ce{O2}) = \dfrac{n_{\max}}{2} = \SI{4,45e-3}{mol}$, soit $\Delta p_{1/2} = \dfrac{n_{\max}\,R\,T}{2V_g} = \dfrac{4,45\times 10^{-3}\times 2478}{9,8\times 10^{-4}} \approx \SI{1,13e4}{Pa} = \SI{113}{hPa}$. Le tableau donne $\Delta p = \SI{86}{hPa}$ à 5 min et \SI{150}{hPa} à 10 min : la valeur \SI{113}{hPa} est bien atteinte entre 5 et 10 min, vers 6 à 7 min : le résultat est cohérent.
\item Pour atteindre plus vite l'état final sans changer les quantités de réactifs : \textbf{élever la température} du milieu, ou \textbf{augmenter la quantité (ou la surface) du catalyseur}.
\end{enumerate}
\end{corrige}

\begin{attention}[Points de vigilance — Exercice 10 — Barème indicatif sur 10]
(1) 1 pt ; (2) 2 pts : conversion hPa $\to$ Pa et L $\to$ m$^3$ obligatoires ; (3) 1,5 pt ; (4) 2 pts : tableau + courbe ; (5) 1,5 pt : tangente + valeur absolue de la pente ; (6) 1 pt : lecture à $C_0/2$ ; (7) 1 pt. Pièges classiques : oublier le facteur 2 de la stœchiométrie ($n(\ce{H2O2}) = n_0 - 2x$), utiliser $V$ au lieu de $V_g$, laisser la pression en hPa.
\end{attention}

\begin{corrige}[Exercice 11 — Type Bac : étude cinétique de l'estérification et vie pratique]
\begin{enumerate}
\item \textbf{Équation-bilan :}
\[ \ce{CH3COOH + C2H5OH <=> CH3COOC2H5 + H2O} \]
La réaction inverse est l'\textbf{hydrolyse} de l'éthanoate d'éthyle. Caractères de la transformation : elle est \textbf{lente}, \textbf{limitée} (elle conduit à un équilibre) et \textbf{réversible}.
\item \textbf{Tableau d'avancement} (quantités en mol) :
\begin{center}
\begin{tabular}{|l|c|c|c|c|}
\hline
\rowcolor{popblueL} Équation & \ce{CH3COOH} & \ce{+ C2H5OH <=>} & \ce{CH3COOC2H5} & \ce{+ H2O} \\
\hline
État initial & $0,20$ & $0,20$ & $0$ & $0$ \\
\hline
En cours & $0,20 - x$ & $0,20 - x$ & $x$ & $x$ \\
\hline
État final & $0,20 - x_f$ & $0,20 - x_f$ & $x_f$ & $x_f$ \\
\hline
\end{tabular}
\end{center}
Le mélange est équimolaire et les coefficients sont tous égaux à 1 : $x_{\max} = \SI{0,20}{mol}$.
L'asymptote donne $x_f = n_{\infty}(\text{ester}) = \SI{0,134}{mol}$, d'où :
\[ \tau = \frac{x_f}{x_{\max}} = \frac{0,134}{0,20} = 0,67 = \SI{67}{\%} \]
$\tau < 1$ : la réaction n'est \textbf{pas totale}, elle est limitée par l'hydrolyse inverse (équilibre chimique).
\item \textbf{Ballon B plus rapide que A :} à température identique (\SI{70}{\celsius}), le ballon B contient de l'acide sulfurique, source d'ions \ce{H+} qui \textbf{catalysent} l'estérification : l'équilibre est atteint en \SI{3}{h} au lieu de \SI{12}{h}. Catalyseur et réactifs sont dans la même phase liquide : c'est une \cle{catalyse homogène}.
\textbf{Ballon C plus lent :} à \SI{20}{\celsius}, sans catalyseur, la température (facteur cinétique) est bien plus faible : les chocs efficaces sont rares et la réaction est très lente ; l'asymptote n'est pas atteinte en \SI{48}{h}.
\item \textbf{L'affirmation est fausse.} Les ballons A et B atteignent la \textbf{même asymptote} $n_{\infty} = \SI{0,134}{mol}$, donc le même rendement ($\tau = \SI{67}{\%}$). L'acide sulfurique n'a fait qu'\textbf{accélérer} la réaction : un catalyseur ne déplace pas un équilibre, il ne change pas le taux d'avancement final.
\item Pour augmenter l'asymptote, il faut déplacer l'équilibre dans le sens direct :
\begin{itemize}
\item utiliser un \textbf{excès} de l'un des réactifs (par exemple l'éthanol, peu coûteux) : l'excès d'un réactif déplace l'équilibre vers la formation des produits ;
\item \textbf{éliminer l'eau} (ou l'ester) au fur et à mesure de sa formation (agent déshydratant, distillation) : la disparition d'un produit déplace l'équilibre dans le sens direct.
\end{itemize}
\item \textbf{Applications quotidiennes :}
\begin{itemize}
\item \textbf{(a) Poisson au frais :} la \textbf{température} est un facteur cinétique ; le froid ralentit fortement les réactions de décomposition (oxydation, prolifération microbienne), ce qui prolonge la conservation.
\item \textbf{(b) Charbon broyé :} le broyage augmente la \textbf{surface de contact} entre le combustible solide et le dioxygène de l'air : la combustion démarre plus vite et plus facilement.
\item \textbf{(c) Levure dans la pâte :} les levures contiennent des \textbf{enzymes} qui jouent le rôle de catalyseurs biologiques (\cle{catalyse enzymatique}) de la fermentation ; le dégagement de \ce{CO2} qui fait lever la pâte est ainsi beaucoup plus rapide.
\end{itemize}
\end{enumerate}
\end{corrige}

\begin{attention}[Points de vigilance — Exercice 11 — Barème indicatif sur 10]
(1) 1,5 pt : double flèche \ce{<=>} exigée, deux caractères cités ; (2) 2 pts : tableau, $x_{\max}$, $\tau$ en \% ; (3) 2 pts : facteurs cinétiques nommés + catalyse homogène justifiée ; (4) 1,5 pt : argumentation appuyée sur les valeurs ($\SI{0,134}{mol}$ identique) ; (5) 1,5 pt ; (6) 1,5 pt : un facteur cinétique correctement identifié par pratique. Erreur fréquente : affirmer que le catalyseur augmente le rendement — il n'augmente que la vitesse.
\end{attention}

\newpage

\coursec{Fiche bilan}

\begin{popfigure}
\begin{center}
\begin{tikzpicture}[
  mindmap/.style={},
  central/.style={circle, draw=popdark, fill=popblue, text=white, font=\bfseries, minimum size=2.6cm, align=center, inner sep=2pt},
  branche/.style={rounded corners=4pt, draw=#1, fill=#1!18, text=popdark, font=\small\bfseries, align=center, inner sep=4pt},
  lien/.style={thick, draw=#1, line width=1.6pt}
]
\node[central] (C) {Cinétique\\ chimique};
\node[branche=popblue,   anchor=west]      (B1) at ([xshift=0.9cm,yshift=2.6cm]C.east)  {Réactions rapides,\\ lentes, très lentes};
\node[branche=poporange, anchor=west]      (B2) at ([xshift=1.6cm,yshift=0.6cm]C.east)  {Vitesses : moyenne\\ et instantanée};
\node[branche=popgreen,  anchor=east]      (B3) at ([xshift=-0.9cm,yshift=2.6cm]C.west) {Vitesse volumique\\ $v=\dfrac{1}{V}\dfrac{dx}{dt}$};
\node[branche=poppurple, anchor=east]      (B4) at ([xshift=-1.6cm,yshift=0.6cm]C.west) {Temps de\\ demi-réaction $t_{1/2}$};
\node[branche=popgold,   anchor=south]     (B5) at ([yshift=1.2cm]C.north)              {Suivi temporel :\\ trempe, colorimétrie,\\ conductimétrie};
\node[branche=poppink,   anchor=west]      (B6) at ([xshift=0.9cm,yshift=-2.4cm]C.east) {Facteurs cinétiques :\\ $C$, $T$, catalyseur,\\ surface de contact};
\node[branche=popblue!60!popgreen, anchor=east] (B7) at ([xshift=-0.9cm,yshift=-2.4cm]C.west) {Estérification :\\ équilibre lent,\\ asymptote inchangée};
\draw[lien=popblue]   (C) -- (B1);
\draw[lien=poporange] (C) -- (B2);
\draw[lien=popgreen]  (C) -- (B3);
\draw[lien=poppurple] (C) -- (B4);
\draw[lien=popgold]   (C) -- (B5);
\draw[lien=poppink]   (C) -- (B6);
\draw[lien=popblue!60!popgreen] (C) -- (B7);
\end{tikzpicture}
\end{center}
\legende{Carte mentale de la leçon : les sept idées-forces de la cinétique chimique.}
\end{popfigure}

\begin{bilan}[L'essentiel de la cinétique chimique]
\textbf{Définitions clés}
\begin{itemize}
  \item \cle{Réaction rapide} : achevée dès le mélange des réactifs (ex. \ce{Ag+ + Cl- -> AgCl}, réactions acide--base, précipitations).
  \item \cle{Réaction lente} : évolue de quelques secondes à plusieurs heures (ex. \ce{S2O8^2- + 2 I- -> 2 SO4^2- + I2}, estérification).
  \item \cle{Réaction très lente} : jours, mois, années (corrosion du fer, synthèse de l'eau sans catalyseur).
  \item \cle{Vitesse volumique} de formation d'un produit : $v = \dfrac{1}{V}\,\dfrac{dn_{\text{produit}}}{dt} = \dfrac{1}{V}\dfrac{dx}{dt}$ ; de disparition d'un réactif : $v = -\dfrac{1}{V}\,\dfrac{dn_{\text{réactif}}}{dt}$.
  \item \cle{Temps de demi-réaction} $t_{1/2}$ : durée au bout de laquelle l'avancement atteint la moitié de sa valeur finale : $x(t_{1/2}) = \dfrac{x_{\text{final}}}{2}$.
  \item \cle{Catalyseur} : espèce qui augmente la vitesse d'une réaction sans être consommée et sans déplacer l'état final (homogène : \ce{Fe^3+} ; hétérogène : platine ; enzymatique).
\end{itemize}

\textbf{Équations-types à connaître par c\oe ur}
\begin{itemize}
  \item Oxydation des iodures : \ce{S2O8^2- + 2 I- -> 2 SO4^2- + I2}
  \item Dosage du diiode (trempe) : \ce{I2 + 2 S2O3^2- -> 2 I- + S4O6^2-}
  \item Estérification : \ce{CH3COOH + C2H5OH <=> CH3COOC2H5 + H2O}
  \item Précipitation rapide : \ce{Ag+ + Cl- -> AgCl}
\end{itemize}

\textbf{Formules}
\begin{center}
\begin{tabularx}{\linewidth}{|l|X|l|}
\hline
\rowcolor{popblueL} \textbf{Grandeur} & \textbf{Expression} & \textbf{Unité} \\
\hline
Vitesse moyenne de formation & $v_{\text{moy}} = \dfrac{1}{V}\,\dfrac{\Delta n}{\Delta t} = \dfrac{[\text{P}]_{t_2}-[\text{P}]_{t_1}}{t_2-t_1}$ & \si{mol.L^{-1}.s^{-1}} \\
\hline
Vitesse instantanée & $v(t) = \dfrac{1}{V}\dfrac{dn}{dt}$ : coefficient directeur de la tangente à la courbe $n=f(t)$ en $t$ & \si{mol.L^{-1}.s^{-1}} \\
\hline
Vitesse de disparition & $v = -\dfrac{d[\text{R}]}{dt}$ (signe moins car $[\text{R}]$ décroît) & \si{mol.L^{-1}.s^{-1}} \\
\hline
Lien avec l'avancement & $v = \dfrac{1}{V}\dfrac{dx}{dt}$ & \si{mol.L^{-1}.s^{-1}} \\
\hline
Demi-réaction & $x(t_{1/2}) = \dfrac{x_{\text{final}}}{2}$ & \si{s} ou \si{min} \\
\hline
\end{tabularx}
\end{center}

\textbf{Méthodes express}
\begin{itemize}
  \item \emph{Vitesse instantanée graphique} : tracer la tangente à la courbe au point d'abscisse $t$, prendre deux points éloignés de la tangente, calculer $\dfrac{\Delta y}{\Delta x}$.
  \item \emph{Suivi par prélèvements} : prélever, \textbf{tremper} (dilution glacée pour bloquer la réaction), doser, reporter sur la courbe.
  \item \emph{Méthodes physiques} : colorimétrie/spectrophotométrie (diiode coloré), conductimétrie (ions), pression (gaz dégagé).
  \item \emph{Accélérer une réaction} : augmenter la concentration des réactifs, la température, la surface de contact, ou ajouter un catalyseur.
  \item \emph{Estérification} : température ou catalyseur (\ce{H2SO4}) $\Rightarrow$ asymptote atteinte plus vite, mais \textbf{limite inchangée} (même équilibre).
\end{itemize}
\end{bilan}

\begin{aretenir}[Checklist avant le Bac]
\begin{itemize}
  \item $\square$ Je sais citer un exemple de réaction rapide, lente et très lente.
  \item $\square$ Je sais écrire et équilibrer \ce{S2O8^2- + 2 I- -> 2 SO4^2- + I2} et le dosage \ce{I2 + 2 S2O3^2- -> 2 I- + S4O6^2-}.
  \item $\square$ Je sais définir la vitesse moyenne et la calculer entre deux dates.
  \item $\square$ Je sais tracer une tangente et en déduire la vitesse instantanée avec son unité.
  \item $\square$ Je sais que la vitesse de disparition porte un signe moins et reste positive.
  \item $\square$ Je connais la relation $v = \dfrac{1}{V}\dfrac{dx}{dt}$ et je sais utiliser un tableau d'avancement.
  \item $\square$ Je sais définir et lire graphiquement le temps de demi-réaction $t_{1/2}$.
  \item $\square$ Je sais expliquer le rôle de la trempe dans un suivi par prélèvements.
  \item $\square$ Je connais les quatre facteurs cinétiques et je sais les illustrer expérimentalement.
  \item $\square$ Je sais montrer le rôle catalytique de \ce{Fe^3+} (homogène) et du platine (hétérogène).
  \item $\square$ Je sais qu'un catalyseur accélère la réaction sans modifier l'état final.
  \item $\square$ Je sais interpréter la courbe d'estérification : asymptote, effet de la température et de \ce{H2SO4}.
\end{itemize}
\end{aretenir}

\findecours

\end{document}
